% Lesson 2 — Grammar: Revise the problem areas in grammar — Anglais Tle A — Cours signé OnBuch+
% Programme officiel MINESEC. Compiler avec : tectonic EN02-grammar-revise-the-problem-areas-in-grammar.tex
\newcommand{\DOCMATIERE}{Anglais}
\newcommand{\DOCNIVEAU}{Tle A}
\newcommand{\DOCMODULE}{Chapter 1 : Family and Social Life}
\newcommand{\DOCLECON}{EN02}
\newcommand{\DOCTITRE}{Lesson 2 — Grammar: Revise the problem areas in grammar}
% OnBuch+ — préambule « cours signés » ENRICHI (compilé avec tectonic / XeLaTeX)
% Système visuel « Pop » d'OnBuch+ : fond crème (jamais blanc), bordures encre
% épaisses, ombres « dures » décalées, titres Archivo Black en capitales.
% Version MATHÉMATIQUES : graphes (pgfplots/tikz), boîtes pédagogiques
% (définition, propriété, méthode, attention, le savais-tu, exercice résolu,
% exercice, TP, bilan), page de garde et signature OnBuch+ ancrée en bas.
%
% Macros attendues dans main.tex AVANT \input{preamble} :
%   \DOCMATIERE  \DOCNIVEAU  \DOCMODULE  \DOCLECON (numéro)  \DOCTITRE
\documentclass[11pt]{article}

\usepackage{fontspec}
\usepackage[english,french]{babel} % français = langue principale ; l'anglais via \eng{..} / textebox
\usepackage[a4paper,margin=2cm,top=3.4cm,bottom=2.8cm,headheight=1.4cm,footskip=1.3cm]{geometry}
\usepackage{amsmath,amssymb,amsfonts}
\usepackage{mathtools} % \xmapsto, \coloneqq... souvent utilisés par les modèles
\usepackage{cancel} % simplifications barrées (bilans de forces)
% \abs{z} (module, valeur absolue) et \norm{u} (norme), utilisés par les modèles
\providecommand{\abs}{}\let\abs\relax\DeclarePairedDelimiter{\abs}{\lvert}{\rvert}
\providecommand{\norm}{}\let\norm\relax\DeclarePairedDelimiter{\norm}{\lVert}{\rVert}
\usepackage[table]{xcolor} % [table] : fournit aussi \rowcolors (alternance de lignes)
\usepackage{pagecolor}
\usepackage{tikz}
\usetikzlibrary{arrows.meta,positioning,calc,shapes.geometric,shapes.misc,shapes.arrows,decorations.pathmorphing,decorations.pathreplacing,decorations.markings,patterns,shadows,mindmap,trees,fit,backgrounds,matrix,intersections,angles,quotes}
\usepackage{pgfplots}
\usepackage[version=4]{mhchem} % équations nucléaires \ce{^{235}_{92}U}
\usepackage[european,siunitx]{circuitikz} % schémas électriques
\pgfplotsset{compat=1.18}
\usepgfplotslibrary{fillbetween,groupplots} % aires entre courbes (calcul intégral)
\usepackage{enumitem}
\usepackage{fancyhdr}
% [most] charge toutes les bibliothèques tcolorbox (listings, minted, raster,
% documentation...) et gonfle inutilement la mémoire TeX sur un document long
% (~300 boîtes) : on ne charge que ce qui est réellement utilisé.
\usepackage[skins,breakable]{tcolorbox}
\usepackage{graphicx}
\usepackage{colortbl}
\usepackage{booktabs}
\usepackage{tabularx}
\usepackage{diagbox} % case d'en-tête diagonale des tableaux à double entrée
% Colonnes extensibles centrée / gauche / droite (C, L, R), souvent utilisées par les modèles
\newcolumntype{C}{>{\centering\arraybackslash}X}
\newcolumntype{L}{>{\raggedright\arraybackslash}X}
\newcolumntype{R}{>{\raggedleft\arraybackslash}X}
\usepackage{array}
\usepackage{multicol}
\usepackage{multirow}
\usepackage{siunitx}
% parse-numbers=false : accepte aussi \SI{2,5 \times 10^{-3}}{...}
\sisetup{locale=FR,output-decimal-marker={,},group-minimum-digits=4,parse-numbers=false}
\DeclareSIUnit\ohm{\ensuremath{\symup{\Omega}}} % U+2126 absent de la police
\DeclareSIUnit\division{div} % réglages d'oscilloscope (ms/div, V/div)
\DeclareSIUnit\tr{tr} % tours (vitesse de rotation, ex. \SI{1500}{\tr\per\minute})
\DeclareSIUnit\molar{M} % concentration molaire (ex. \SI{3}{\molar}, \SI{1}{\micro\molar})
\DeclareSIUnit\an{an} % année (le modèle confond parfois avec \year, un compteur TeX) — ex. \SI{5}{\kilo\meter\per\an}
\DeclareSIUnit\FCFA{FCFA} % franc CFA (ex. \SI{500}{\FCFA\per\kilo\gram})
\DeclareSIUnit\degreeBrix{\textdegree Brix} % teneur en sucre d'un sirop/jus (réfractométrie)
\DeclareSIUnit\month{mois} % le modèle utilise parfois \month, un compteur TeX, comme unité de durée
\usepackage{qrcode}
% \degree : commande fréquemment utilisée par les modèles pour un angle ou une
% température en dehors de \SI{}{} (non fournie par les paquets chargés).
\providecommand{\degree}{^{\circ}}
% \qed (fin de démonstration), utilisé par les modèles sans amsthm
\providecommand{\qed}{\hfill$\square$}
% \barre : double barre des valeurs interdites dans un tableau de variations (tkz-tab)
\providecommand{\barre}{\ensuremath{\|}}
% environnement proof (amsthm non chargé) utilisé par les modèles
\ifdefined\proof\else\newenvironment{proof}[1][Démonstration]{\par\noindent\textit{#1.}\ }{\qed\par}\fi
\usepackage{lastpage}
\usepackage{hyperref}
\hypersetup{hidelinks}

% ============================================================
% Palette « Pop » OnBuch+ — mêmes hex que la classe P (plus_ui.dart)
% ============================================================
\definecolor{poporange}{HTML}{F59321}
\definecolor{popdark}{HTML}{E07A0C}
\definecolor{popcream}{HTML}{FFF6E8}
\definecolor{popink}{HTML}{1C1714}
\definecolor{poppurple}{HTML}{7A5AE0}
\definecolor{popgreen}{HTML}{2E9E6A}
\definecolor{poppink}{HTML}{E0457B}
\definecolor{popblue}{HTML}{2F7DE1}
\definecolor{popgold}{HTML}{C98A12}
\definecolor{popmuted}{HTML}{8A7F76}
\definecolor{popline}{HTML}{EADFD2}
\definecolor{popgray}{HTML}{8A7F76}
\colorlet{popgrey}{popgray}
% Alias tolérants (le modèle nomme parfois les couleurs différemment)
\colorlet{popwhite}{white}
\colorlet{popblack}{popink}
\colorlet{popred}{poppink}
\colorlet{popyellow}{popgold}
% Teintes claires (fonds de tableaux/encadrés) — le modèle nomme parfois
% les couleurs avec un suffixe L (ex. popblueL) sans les avoir définies.
\colorlet{poporangeL}{poporange!20}
\colorlet{popdarkL}{popdark!20}
\colorlet{poppurpleL}{poppurple!20}
\colorlet{popgreenL}{popgreen!20}
\colorlet{poppinkL}{poppink!20}
\colorlet{popblueL}{popblue!20}
\colorlet{popgoldL}{popgold!20}
\colorlet{popredL}{popred!20}
\colorlet{popgrayL}{popgray!20}
% Teintes claires pour fonds de tableaux / zones
\colorlet{popblueL}{popblue!10!white}
\colorlet{poporangeL}{poporange!14!white}
\colorlet{popgreenL}{popgreen!12!white}
\colorlet{poppurpleL}{poppurple!12!white}
\colorlet{poppinkL}{poppink!12!white}
\colorlet{popgoldL}{popgold!14!white}

\pagecolor{popcream}

% ============================================================
% Typographie
% ============================================================
\defaultfontfeatures{Ligatures=TeX}
\setmainfont{PlusJakartaSans-Medium.ttf}[Path=../../../fonts/,BoldFont=PlusJakartaSans-Bold.ttf]
% Symboles, API et emojis absents de la police principale : repli sur DejaVu Sans (caractères actifs)
\newfontface\symfont{DejaVuSans.ttf}[Path=../../../fonts/]
\makeatletter
\newcommand\symactive[1]{\catcode#1=13 \begingroup\lccode`\~=#1 \lowercase{\endgroup\def~}{{\symfont\char#1}}}
\newcommand\symblank[1]{\catcode#1=13 \begingroup\lccode`\~=#1 \lowercase{\endgroup\def~}{}}
\newcommand\symhyph[1]{\catcode#1=13 \begingroup\lccode`\~=#1 \lowercase{\endgroup\def~}{\mbox{-}}}
\newcount\symc
\newcommand\symrange[2]{\symc=#1 \@whilenum\symc<#2 \do{\expandafter\symactive\expandafter{\the\symc}\advance\symc by 1 }}
\newcommand\symblankrange[2]{\symc=#1 \@whilenum\symc<#2 \do{\expandafter\symblank\expandafter{\the\symc}\advance\symc by 1 }}
\makeatother
\symrange{"0250}{"02FF}\symactive{"2713}\symactive{"2714}\symactive{"2717}\symactive{"2718}\symactive{"2610}\symactive{"2611}\symactive{"2612}
\symhyph{"2011}\symhyph{"2E1A}\symblankrange{"1F300}{"1FAFF}\symblank{"FE0F}
% Maths en sans-serif (Fira Math) avec chiffres et lettres droites de Plus
% Jakarta, pour que les formules se fondent dans le texte.
\usepackage[math-style=ISO,bold-style=ISO]{unicode-math}
\setmathfont{FiraMath-Regular.otf}[Path=../../../fonts/]
\setmathfont{PlusJakartaSans-Medium.ttf}[Path=../../../fonts/,range={up/num,up/{Latin,latin},\mathup}]
\usepackage{icomma} % virgule décimale sans espace en mode math
% Caractères mathématiques Unicode absents des polices de texte (Plus Jakarta,
% Archivo Black) : rendus par leur équivalent mathématique, en texte comme en
% formule — sinon ils s'affichent en carré vide (ex. « Arithmétique dans ℤ »).
\usepackage{newunicodechar}
\newunicodechar{ℝ}{\ensuremath{\mathbb{R}}}
\newunicodechar{ℤ}{\ensuremath{\mathbb{Z}}}
\newunicodechar{ℂ}{\ensuremath{\mathbb{C}}}
\newunicodechar{ℕ}{\ensuremath{\mathbb{N}}}
\newunicodechar{ℚ}{\ensuremath{\mathbb{Q}}}
\newunicodechar{𝔻}{\ensuremath{\mathbb{D}}}
\newunicodechar{α}{\ensuremath{\alpha}}
\newunicodechar{β}{\ensuremath{\beta}}
\newunicodechar{γ}{\ensuremath{\gamma}}
\newunicodechar{δ}{\ensuremath{\delta}}
\newunicodechar{ε}{\ensuremath{\varepsilon}}
\newunicodechar{θ}{\ensuremath{\theta}}
\newunicodechar{λ}{\ensuremath{\lambda}}
\newunicodechar{μ}{\ensuremath{\mu}}
\newunicodechar{σ}{\ensuremath{\sigma}}
\newunicodechar{τ}{\ensuremath{\tau}}
\newunicodechar{φ}{\ensuremath{\varphi}}
\newunicodechar{ψ}{\ensuremath{\psi}}
\newunicodechar{ω}{\ensuremath{\omega}}
\newunicodechar{Σ}{\ensuremath{\Sigma}}
% majuscules produites par \MakeUppercase dans les titres (α-aminés -> Α)
\newunicodechar{Α}{\ensuremath{\alpha}}
\newunicodechar{Β}{\ensuremath{\beta}}
\newunicodechar{Γ}{\ensuremath{\Gamma}}
\newunicodechar{Δ}{\ensuremath{\Delta}}
\newunicodechar{Ε}{\ensuremath{\varepsilon}}
\newunicodechar{Θ}{\ensuremath{\Theta}}
\newunicodechar{Λ}{\ensuremath{\Lambda}}
\newunicodechar{Μ}{\ensuremath{\mu}}
\newunicodechar{Τ}{\ensuremath{\tau}}
\newunicodechar{Φ}{\ensuremath{\Phi}}
\newunicodechar{Ψ}{\ensuremath{\Psi}}
\newunicodechar{Ω}{\ensuremath{\Omega}}
\newunicodechar{↦}{\ensuremath{\mapsto}}
\newunicodechar{∈}{\ensuremath{\in}}
\newunicodechar{∉}{\ensuremath{\notin}}
\newunicodechar{∧}{\ensuremath{\wedge}}
\newunicodechar{⇔}{\ensuremath{\Leftrightarrow}}
\newunicodechar{⟺}{\ensuremath{\Longleftrightarrow}}
\newunicodechar{⇒}{\ensuremath{\Rightarrow}}
\newunicodechar{⟹}{\ensuremath{\Longrightarrow}}
\newunicodechar{∘}{\ensuremath{\circ}}
\newunicodechar{∪}{\ensuremath{\cup}}
\newunicodechar{∩}{\ensuremath{\cap}}
\newunicodechar{≡}{\ensuremath{\equiv}}
\newunicodechar{⊂}{\ensuremath{\subset}}
\newunicodechar{⊥}{\ensuremath{\perp}}
\newunicodechar{∃}{\ensuremath{\exists}}
\newunicodechar{∀}{\ensuremath{\forall}}
\newunicodechar{∛}{\ensuremath{\sqrt[3]{\,}}}
\newunicodechar{∜}{\ensuremath{\sqrt[4]{\,}}}
\newunicodechar{⃗}{\ensuremath{{}^{\to}}}
\newunicodechar{ⁿ}{\ifmmode^{n}\else\textsuperscript{\ensuremath{n}}\fi}
\newunicodechar{ˣ}{\ifmmode^{x}\else\textsuperscript{\ensuremath{x}}\fi}
\newunicodechar{ᵇ}{\ifmmode^{b}\else\textsuperscript{\ensuremath{b}}\fi}
\newunicodechar{ʳ}{\ifmmode^{r}\else\textsuperscript{\ensuremath{r}}\fi}
\newunicodechar{ᵘ}{\ifmmode^{u}\else\textsuperscript{\ensuremath{u}}\fi}
\newunicodechar{ʸ}{\ifmmode^{y}\else\textsuperscript{\ensuremath{y}}\fi}
\newunicodechar{ᵃ}{\ifmmode^{a}\else\textsuperscript{\ensuremath{a}}\fi}
\newunicodechar{ᵉ}{\ifmmode^{e}\else\textsuperscript{\ensuremath{e}}\fi}
\newunicodechar{ᵏ}{\ifmmode^{k}\else\textsuperscript{\ensuremath{k}}\fi}
\newunicodechar{ᵖ}{\ifmmode^{p}\else\textsuperscript{\ensuremath{p}}\fi}
\newunicodechar{ᴺ}{\ifmmode^{N}\else\textsuperscript{\ensuremath{N}}\fi}
\newunicodechar{ᶜ}{\ifmmode^{c}\else\textsuperscript{\ensuremath{c}}\fi}
\newunicodechar{ʰ}{\ifmmode^{h}\else\textsuperscript{\ensuremath{h}}\fi}
\newunicodechar{ᵠ}{\ifmmode^{q}\else\textsuperscript{\ensuremath{q}}\fi}
\newunicodechar{ₙ}{\ifmmode_{n}\else\textsubscript{\ensuremath{n}}\fi}
\newunicodechar{ₐ}{\ifmmode_{a}\else\textsubscript{\ensuremath{a}}\fi}
\newunicodechar{ᵢ}{\ifmmode_{i}\else\textsubscript{\ensuremath{i}}\fi}
\newunicodechar{ᵣ}{\ifmmode_{r}\else\textsubscript{\ensuremath{r}}\fi}
\newunicodechar{ₖ}{\ifmmode_{k}\else\textsubscript{\ensuremath{k}}\fi}
\newunicodechar{ᵦ}{\ifmmode_{\beta}\else\textsubscript{\ensuremath{\beta}}\fi}
\newunicodechar{ₓ}{\ifmmode_{x}\else\textsubscript{\ensuremath{x}}\fi}
\newfontfamily\popdisplay{ArchivoBlack-Regular.ttf}[Path=../../../fonts/]
\newfontfamily\poplogo{SpaceGrotesk-Bold.ttf}[Path=../../../fonts/]
\newfontfamily\popsemi{PlusJakartaSans-SemiBold.ttf}[Path=../../../fonts/]
\newfontfamily\popxbold{PlusJakartaSans-ExtraBold.ttf}[Path=../../../fonts/]
\newfontfamily\popxboldls{PlusJakartaSans-ExtraBold.ttf}[Path=../../../fonts/,LetterSpace=3.0]

\setlist[enumerate]{leftmargin=1.7em,itemsep=5pt,topsep=4pt}
\setlist[itemize]{leftmargin=1.7em,itemsep=4pt,topsep=4pt,label={\color{poporange}\textbullet}}
\setlength{\parindent}{0pt}
\setlength{\parskip}{5pt}
\renewcommand{\arraystretch}{1.25}

% pgfplots : style Pop par défaut
\pgfplotsset{
  every axis/.append style={axis line style={popink,line width=1pt},tick style={popink},
    grid style={popline,line width=0.5pt},label style={font=\small},tick label style={font=\footnotesize}},
  popcurve/.style={poporange,line width=1.8pt,no markers,smooth},
}
% Même style disponible dans un \draw TikZ simple (hors pgfplots)
\tikzset{popcurve/.style={poporange,line width=1.8pt}}
\tikzset{popgrid/.style={popline,very thin}}
\colorlet{popcurve}{poporange} % le nom du style est parfois utilisé comme couleur

% ============================================================
% Pill / squiggle
% ============================================================
\newcommand{\poppill}[3]{%
  \tikz[baseline=-0.55ex]\node[draw=popink,line width=1.2pt,rounded corners=2.6pt,%
    fill=#2,inner xsep=6.5pt,inner ysep=3pt]{%
    \popxboldls\fontsize{7}{7}\selectfont\color{#3}\MakeUppercase{#1}};%
}
\newcommand{\popsquiggle}[1]{%
  \tikz[baseline=-0.4ex]{
    \draw[#1,line width=2.6pt,line cap=round]
      plot [smooth] coordinates {(0,0.09) (0.34,0.01) (0.68,0.21) (1.02,0.01) (1.36,0.10)};
  }%
}
% Mot-clé surligné (marqueur orange sous le texte)
\newcommand{\cle}[1]{{\popxbold\color{popdark}#1}}

% ============================================================
% En-tête / pied
% ============================================================
\pagestyle{fancy}
\fancyhf{}
\renewcommand{\headrulewidth}{0pt}
\renewcommand{\footrulewidth}{0pt}
\lhead{\raisebox{-0.15ex}{\poplogo\fontsize{13}{13}\selectfont\color{popink}OnBuch\color{poporange}+}}
\rhead{\poppill{\DOCMATIERE\ \textbullet\ \DOCNIVEAU\ \textbullet\ Leçon \DOCLECON}{white}{popink}}
\lfoot{\popxbold\fontsize{7.6}{7.6}\selectfont\color{popmuted}\MakeUppercase{Cours signé OnBuch+}\ \textbullet\ {\popsemi\DOCTITRE}}
\rfoot{\tikz[baseline=-0.6ex]\node[rounded corners=8.5pt,fill=popink,minimum height=17pt,minimum width=17pt,inner xsep=5pt,inner sep=0]{\popxbold\fontsize{8}{8}\selectfont\color{popcream}\thepage\,/\,\pageref*{LastPage}};}

% ============================================================
% Titres
% ============================================================
\newcommand{\chaptitre}[2]{%
  \begin{tcolorbox}[enhanced,colback=poporange,colframe=popink,boxrule=2.6pt,arc=11pt,
      drop shadow=popink,left=15pt,right=15pt,top=12pt,bottom=11pt,before skip=3pt,after skip=18pt]
    \poppill{#2}{popink}{popcream}\par\vspace{9pt}
    {\raggedright\hyphenpenalty=10000\exhyphenpenalty=10000\popdisplay\fontsize{22}{24}\selectfont\color{popcream}\MakeUppercase{#1}\par}
  \end{tcolorbox}
}
% Section numérotée : pastille orange avec le numéro + titre Archivo
\newcounter{popsec}
\newcommand{\coursec}[1]{%
  \stepcounter{popsec}\setcounter{popsub}{0}%
  \par\vspace{16pt}\noindent
  \begin{minipage}{\linewidth}
  \tikz[baseline=-0.7ex]\node[draw=popink,line width=1.6pt,fill=poporange,rounded corners=4pt,minimum size=22pt,inner sep=0]{\popdisplay\fontsize{12}{12}\selectfont\color{popcream}\thepopsec};%
  \hspace{8pt}{\popdisplay\fontsize{15}{17}\selectfont\color{popink}\MakeUppercase{#1}}\par
  \vspace{4pt}\noindent{\color{poporange}\rule{36pt}{3.6pt}}
  \end{minipage}\par\nopagebreak\vspace{8pt}
}
\newcounter{popsub}
\newcommand{\courssub}[1]{%
  \stepcounter{popsub}%
  \par\vspace{9pt}\noindent{\popxbold\fontsize{11.5}{13}\selectfont\color{popink}{\color{poporange}\thepopsec.\thepopsub}\hspace{6pt}#1}\par\nopagebreak\vspace{4pt}
}

% ============================================================
% Boîtes pédagogiques — même recette : fond blanc, bordure épaisse colorée,
% ombre dure encre, badge plein. Titre optionnel [#1].
% ============================================================
\tcbset{popbase/.style={breakable,enhanced,colback=white,boxrule=2.3pt,arc=9pt,
  shadow={3pt}{-3pt}{0pt}{popink},left=14pt,right=14pt,top=11pt,bottom=11pt,before skip=11pt,after skip=15pt,
  fontupper=\popsemi\fontsize{10.6}{14.5}\selectfont\color{popink},
  fontlower=\popsemi\fontsize{10.6}{14.5}\selectfont\color{popink}}}
% Coche dessinée (le glyphe ✓ n'existe pas dans les polices du document)
\newcommand{\popcheck}[1]{\tikz[baseline=-0.5ex]\draw[#1,line width=1.6pt,line cap=round,line join=round](-0.13,0.0)--(-0.04,-0.09)--(0.13,0.1);}
\providecommand{\coche}{\popcheck{popgreen}}
\providecommand{\resultat}[1]{\par\smallskip\noindent\fcolorbox{popgreen}{popgreenL}{\parbox{\dimexpr\linewidth-2\fboxsep-2\fboxrule\relax}{\textbf{Résultat.} #1}}\par\smallskip}
\newcommand{\popboxhead}[3]{\poppill{#1}{#2}{white}\ifx\relax#3\relax\else\hspace{6pt}{\popxbold\fontsize{10.6}{12}\selectfont\color{popink}#3}\fi\par\vspace{7pt}}

\newtcolorbox{aretenir}[1][]{popbase,colframe=popgreen,before upper={\popboxhead{À retenir}{popgreen}{#1}}}
% Texte en anglais (document de lecture, dialogue, extrait) : cadre
% d'étude. Un titre optionnel (source, auteur) s'affiche dans l'en-tête.
\newtcolorbox{textebox}[1][]{popbase,colframe=popblue,before upper={\popboxhead{Text}{popblue}{#1}\begin{otherlanguage}{english}},after upper={\end{otherlanguage}}}
% Marques ✓ / ✗ (forme correcte / fautive) souvent utilisées par les modèles
\providecommand{\cross}{\textcolor{poppink}{\ensuremath{\times}}}
\providecommand{\xmark}{\cross}\providecommand{\crossmark}{\cross}\providecommand{\cmark}{\ensuremath{\checkmark}}
% Ligne de réponse / justification à compléter (inventée par les modèles)
\providecommand{\justification}[1]{\par\noindent\textit{Justification~:} \dotfill\par}
% \textipa (package tipa non chargé) : la police ne couvre pas l'API -> simple passage
\providecommand{\textipa}[1]{#1}
% Soulignement (ulem non chargé) : \uline, \ul
\providecommand{\uline}[1]{\underline{#1}}\providecommand{\ul}[1]{\underline{#1}}
% \glossaire{\item ...} inventée par les modèles : liste de glossaire
\providecommand{\glossaire}[1]{\begin{itemize}#1\end{itemize}}
% \alert (beamer) inventée par les modèles : texte mis en évidence
\providecommand{\alert}[1]{\textbf{\textcolor{poppink}{#1}}}
% variantes ASCII d'ordinaux écrites en macro
\providecommand{\iere}{\textsuperscript{re}}\providecommand{\ieme}{\textsuperscript{e}}\providecommand{\ere}{\textsuperscript{re}}\providecommand{\eme}{\textsuperscript{e}}\providecommand{\er}{\textsuperscript{er}}\providecommand{\ie}{\textsuperscript{e}}
% Ordinaux français écrits en macro par les modèles : 1\ière, 3\ième
\newcommand{\ière}{\textsuperscript{re}}\newcommand{\ième}{\textsuperscript{e}}
% \exemple (inventée par les modèles) : étiquette « Exemple : »
\providecommand{\exemple}{\textbf{Exemple~:}\ }
% \square utilisable aussi hors mode math (cases à cocher)
\AtBeginDocument{\let\squareMath\square\renewcommand{\square}{\ensuremath{\squareMath}}}
% Mot ou phrase en anglais à l'intérieur d'un paragraphe français
\newcommand{\eng}[1]{\ifmmode\text{\textit{#1}}\else\begingroup\selectlanguage{english}\itshape #1\endgroup\fi}
\let\esp\eng % alias toléré
% flèches d'intonation inventées par les modèles
\providecommand{\textsearrow}{}\renewcommand{\textsearrow}{\ensuremath{\searrow}}\providecommand{\textnearrow}{}\renewcommand{\textnearrow}{\ensuremath{\nearrow}}
% \fr{..} : mot français (inventé par les modèles)
\providecommand{\fr}[1]{\textit{#1}}
\newtcolorbox{exemplebox}[1][]{popbase,colframe=poporange,before upper={\popboxhead{Exemple}{poporange}{#1}}}
\newtcolorbox{definition}[1][]{popbase,colframe=popblue,before upper={\popboxhead{Définition}{popblue}{#1}}}
\newtcolorbox{propriete}[1][]{popbase,colframe=poppurple,before upper={\popboxhead{Propriété}{poppurple}{#1}}}
\newtcolorbox{methode}[1][]{popbase,colframe=popgold,before upper={\popboxhead{Méthode}{popgold}{#1}}}
\newtcolorbox{attention}[1][]{popbase,colframe=poppink,before upper={\popboxhead{Attention}{poppink}{#1}}}
\newtcolorbox{experience}[1][]{popbase,colframe=popblue,colback=popblueL,before upper={\popboxhead{Expérience}{popblue}{#1}}}
% « Le savais-tu ? » : carte violette pleine (contexte, culture, Cameroun)
\newtcolorbox{savaistu}[1][]{popbase,colback=poppurple,colframe=popink,
  fontupper=\popsemi\fontsize{10.4}{14.2}\selectfont\color{white},
  before upper={\poppill{Le savais-tu\,?}{white}{poppurple}\ifx\relax#1\relax\else\hspace{6pt}{\popxbold\color{white}#1}\fi\par\vspace{7pt}}}
% Exercice résolu : énoncé en haut, \tcblower puis la solution sur fond crème
\newcounter{popexo}\newcounter{popexr}
\newtcolorbox{exoresolu}[1][]{popbase,colframe=poporange,colbacklower=popcream,
  segmentation style={popink,line width=1.2pt,dashed},
  before upper={\stepcounter{popexr}\popboxhead{Exercice résolu \thepopexr}{poporange}{#1}},
  before lower={\poppill{Solution}{popink}{popcream}\par\vspace{6pt}}}
% Exercice (non résolu en place) — la correction est dans la partie Corrigés
\newtcolorbox{exercice}[1][]{popbase,colframe=popink,
  before upper={\stepcounter{popexo}\popboxhead{Exercice \thepopexo}{popink}{#1}}}
% Corrigé
\newtcolorbox{corrige}[1][]{popbase,colframe=popgreen,colback=popgreenL,
  before upper={\popboxhead{Corrigé}{popgreen}{#1}}}
% Objectifs / prérequis / bilan
\newtcolorbox{objectifs}{popbase,colframe=popink,colback=poporangeL,
  before upper={\popboxhead{Objectifs de la leçon}{popink}{}}}
\newtcolorbox{prerequis}{popbase,colframe=popink,colback=popblueL,
  before upper={\popboxhead{Prérequis}{popblue}{}}}
\newtcolorbox{bilan}{popbase,colframe=popink,colback=white,boxrule=2.8pt,
  before upper={\popboxhead{Fiche bilan}{poporange}{L'essentiel à connaître pour l'examen}}}
% Figure encadrée avec légende
\newtcolorbox{popfigure}[1][]{enhanced,breakable=false,colback=white,colframe=popline,boxrule=1.4pt,arc=8pt,
  left=10pt,right=10pt,top=10pt,bottom=8pt,before skip=10pt,after skip=12pt,halign=center,
  lower separated=false,halign lower=center,
  fontlower=\popsemi\fontsize{9}{11.5}\selectfont\color{popmuted},
  #1}
\newcommand{\legende}[1]{\par\vspace{6pt}{\popsemi\fontsize{9}{11.5}\selectfont\color{popmuted}#1}}

% ============================================================
% Signature OnBuch+
% ============================================================
\newcommand{\poplogoinline}[1]{{\poplogo\fontsize{#1}{#1}\selectfont\color{popink}OnBuch\color{poporange}+}}
\newcommand{\signatureonbuch}{%
  \begin{tcolorbox}[enhanced,colback=popink,colframe=popink,boxrule=2.2pt,arc=10pt,
      drop shadow=poporange,left=14pt,right=14pt,top=10pt,bottom=10pt,before skip=0pt,after skip=0pt]
    \tikz[baseline=-0.7ex]\node[circle,fill=popgreen,draw=popcream,line width=1.3pt,minimum size=22pt,inner sep=0]{\popcheck{white}};%
    \hspace{9pt}\begin{minipage}[c]{0.62\linewidth}
      {\popxbold\fontsize{12}{13}\selectfont\color{popcream}Signé\ \textbullet\ {\poplogo OnBuch\color{poporange}+}}\par\vspace{2pt}
      {\raggedright\popsemi\fontsize{8}{10.5}\selectfont\color{popline}\MakeUppercase{Équipe pédagogique OnBuch+}\\\MakeUppercase{Conforme au programme officiel MINESEC}\par}
    \end{minipage}\hfill
    {\poplogo\fontsize{22}{22}\selectfont\color{popcream}OnBuch\color{poporange}+}
  \end{tcolorbox}%
}

% ============================================================
% Page de garde
% #1 titre · #2 matière · #3 niveau · #4 module · #5 n° leçon · #6 durée ·
% #7 description / objectifs (texte ou itemize)
% ============================================================
\newcommand{\pagegarde}[7]{%
  \thispagestyle{empty}
  \newgeometry{a4paper,margin=1.7cm,top=1.6cm,bottom=1.4cm}%
  \begin{tikzpicture}[remember picture,overlay]
    \fill[poporange!18] ([xshift=-3.2cm,yshift=-3.6cm]current page.north east) circle (4.6cm);
    \fill[poppurple!14] ([xshift=2.4cm,yshift=7.6cm]current page.south west) circle (3.8cm);
  \end{tikzpicture}%
  {\poplogo\fontsize{30}{30}\selectfont\color{popink}OnBuch\color{poporange}+}\hfill
  \raisebox{0.6ex}{\poppill{Cours signé \textbullet\ Premium}{popink}{popcream}}\par
  \vspace{1pt}\popsquiggle{poppurple}\par
  \vspace{8pt}
  \begin{tcolorbox}[enhanced,colback=poporange,colframe=popink,boxrule=2.8pt,arc=13pt,
      drop shadow=popink,left=18pt,right=18pt,top=12pt,bottom=13pt,after skip=10pt]
    \poppill{#2 \textbullet\ #3}{popink}{popcream}\hspace{5pt}\poppill{#4}{white}{popink}\par\vspace{8pt}
    {\popdisplay\fontsize{14}{14}\selectfont\color{popink}LEÇON #5}\par\vspace{4pt}
    {\raggedright\hyphenpenalty=10000\exhyphenpenalty=10000\popdisplay\fontsize{25}{27}\selectfont\color{popcream}\MakeUppercase{#1}\par}
    \vspace{8pt}
    {\popxbold\fontsize{9.5}{9.5}\selectfont\color{popink}\MakeUppercase{Durée conseillée : #6}}
  \end{tcolorbox}
  \begin{tcolorbox}[enhanced,colback=white,colframe=popink,boxrule=2.2pt,arc=10pt,
      drop shadow=popink,left=16pt,right=16pt,top=10pt,bottom=10pt,after skip=8pt]
    \poppill{Dans cette leçon}{poppurple}{white}\par\vspace{5pt}
    \popsemi\fontsize{9.3}{12.6}\selectfont\color{popink}#7
  \end{tcolorbox}
  \noindent\begin{minipage}[t]{0.487\linewidth}
    \begin{tcolorbox}[enhanced,colback=popgreen,colframe=popink,boxrule=2.2pt,arc=10pt,
        drop shadow=popink,left=11pt,right=11pt,top=10pt,bottom=10pt,height=3.1cm,valign=center]
      \tcbox[colback=white,colframe=popink,boxrule=1.3pt,arc=5pt,boxsep=0pt,nobeforeafter,
        left=3pt,right=3pt,top=3pt,bottom=3pt,valign=center]{\qrcode[height=1.9cm]{https://play.google.com/store/apps/details?id=com.luvvix.onbuch&hl=fr}}%
      \hspace{8pt}\begin{minipage}[c]{0.5\linewidth}
        {\popdisplay\fontsize{11}{12.5}\selectfont\color{white}\MakeUppercase{Télécharge OnBuch}}\par\vspace{4pt}
        {\raggedright\popsemi\fontsize{8.4}{11}\selectfont\color{white}Gratuit \textbullet\ Google Play\\Tuteur IA, annales, concours, résultats\par}
      \end{minipage}
    \end{tcolorbox}
  \end{minipage}\hfill
  \begin{minipage}[t]{0.487\linewidth}
    \begin{tcolorbox}[enhanced,colback=poppurple,colframe=popink,boxrule=2.2pt,arc=10pt,
        drop shadow=popink,left=11pt,right=11pt,top=10pt,bottom=10pt,height=3.1cm,valign=center]
      \tikz[baseline=-0.7ex]\node[circle,fill=white,draw=popink,line width=1.3pt,minimum size=1.6cm,inner sep=0]{\popdisplay\fontsize{24}{24}\selectfont\color{poppurple}+};%
      \hspace{8pt}\begin{minipage}[c]{0.6\linewidth}
        {\popdisplay\fontsize{11}{12.5}\selectfont\color{white}\MakeUppercase{Passe à OnBuch+}}\par\vspace{4pt}
        {\raggedright\popsemi\fontsize{8.4}{11}\selectfont\color{white}Tous les cours signés, Tuteur IA illimité, fiches PDF — onglet OnBuch+ de l'app\par}
      \end{minipage}
    \end{tcolorbox}
  \end{minipage}
  \vspace{8pt}
  \popcontenu
  \vfill
  \signatureonbuch
  \restoregeometry
  \newpage
}

% Rangée « Ce cours contient » (page de garde)
\newcommand{\poptile}[3]{% #1 couleur #2 gros texte #3 légende
  \begin{tcolorbox}[enhanced,colback=#1,colframe=popink,boxrule=2pt,arc=9pt,shadow={2.5pt}{-2.5pt}{0pt}{popink},
      left=6pt,right=6pt,top=9pt,bottom=9pt,halign=center,valign=center,height=1.75cm,width=\linewidth,nobeforeafter]
    {\popdisplay\fontsize{12.5}{13}\selectfont\color{white}\MakeUppercase{#2}}\par\vspace{3pt}
    {\centering\popsemi\fontsize{7.8}{9.5}\selectfont\color{white}#3\par}
  \end{tcolorbox}}
\newcommand{\popcontenu}{%
  \par\noindent{\popxboldls\fontsize{7.5}{7.5}\selectfont\color{popmuted}CE COURS SIGNÉ CONTIENT}\par\vspace{7pt}
  \noindent\begin{minipage}[t]{0.235\linewidth}\poptile{popblue}{Cours}{complet, illustré, pas à pas}\end{minipage}\hfill
  \begin{minipage}[t]{0.235\linewidth}\poptile{poporange}{Méthodes}{et exercices résolus}\end{minipage}\hfill
  \begin{minipage}[t]{0.235\linewidth}\poptile{poppink}{Exercices}{type examen + corrigés}\end{minipage}\hfill
  \begin{minipage}[t]{0.235\linewidth}\poptile{popgreen}{Fiche bilan}{l'essentiel à retenir}\end{minipage}\par}

% Sommaire
\newtcolorbox{sommaire}{enhanced,colback=white,colframe=popink,boxrule=2.2pt,arc=10pt,
  drop shadow=popink,left=18pt,right=18pt,top=13pt,bottom=13pt,before skip=6pt,after skip=20pt,
  before upper={\poppill{Sommaire}{poppurple}{white}\par\vspace{9pt}\popsemi\fontsize{10.6}{14.5}\selectfont\color{popink}}}

% Signature de fin de cours — poussée tout en bas de la dernière page
\newcommand{\findecours}{%
  \par\vfill\nopagebreak
  \begin{center}\popsquiggle{poporange}\end{center}\vspace{4pt}
  \signatureonbuch
}
\AtBeginDocument{\def\llbracket{\mathopen{[\mkern-3.5mu[}}\def\rrbracket{\mathclose{]\mkern-3.5mu]}}} % intervalles d'entiers (glyphe ⟦ absent de la police)
% Glyphes absents de Fira Math : redéfinis à partir de glyphes disponibles
\AtBeginDocument{%
  \def\setminus{\mathbin{\backslash}}\let\smallsetminus\setminus
  \def\vdots{\mathord{\vbox{\baselineskip3.2pt\lineskiplimit0pt\kern4pt\hbox{.}\hbox{.}\hbox{.}}}}%
  \def\ddots{\mathinner{\mkern1mu\raise6.5pt\hbox{.}\mkern2mu\raise3.6pt\hbox{.}\mkern2mu\raise0.7pt\hbox{.}\mkern1mu}}%
  \def\triangle{\mathord{\scalebox{0.9}{$\symup{\Delta}$}}}%
  \def\nabla{\mathord{\raisebox{\depth}{\scalebox{1}[-1]{$\symup{\Delta}$}}}}%
  \def\checkmark{\popcheck{popdark}}%
  \def\star{\mathbin{\ast}}\def\bigstar{\mathord{\ast}}%
  \def\diamond{\mathbin{\scalebox{0.6}{$\Diamond$}}}%
  \def\bigcup{\mathop{\mathchoice{\raisebox{-0.3em}{\scalebox{1.7}{$\cup$}}}{\scalebox{1.25}{$\cup$}}{\cup}{\cup}}\displaylimits}%
  \def\bigcap{\mathop{\mathchoice{\raisebox{-0.3em}{\scalebox{1.7}{$\cap$}}}{\scalebox{1.25}{$\cap$}}{\cap}{\cap}}\displaylimits}%
  \def\lesseqqgtr{\mathrel{\lessgtr}}\def\bigoplus{\mathop{\oplus}\displaylimits}%
}
\providecommand{\ssi}{si et seulement si} % abréviation fréquente
\AtBeginDocument{\providecommand{\wideparen}{\overparen}} % arc de cercle

%% FIGFIX-BEGIN v3 — correctifs globaux des figures (docs/onbuch-generation/tools/figfix)
\usepackage{adjustbox}\usepackage{etoolbox}
% toute figure TikZ/pgfplots plus large que la ligne est réduite à la largeur disponible
\BeforeBeginEnvironment{tikzpicture}{\begin{adjustbox}{max width=\linewidth}}
\AfterEndEnvironment{tikzpicture}{\end{adjustbox}}
% légendes pgfplots lisibles par-dessus les courbes
\pgfplotsset{every axis/.append style={legend style={fill=white,fill opacity=0.92,text opacity=1,draw=black!15,inner sep=3pt}}}
% étiquettes d'arêtes (syntaxe edge["..."]) avec fond blanc
\tikzset{every edge quotes/.append style={fill=white,inner sep=1.2pt}}
% bulles de cartes mentales : texte à la ligne dans la bulle, bulle un peu plus grande
\tikzset{every concept/.append style={text width=2.0cm,align=center,minimum size=2.3cm,inner sep=2pt}}
%% FIGFIX-END
\begin{document}

\pagegarde{Lesson 2 — Grammar: Revise the problem areas in grammar}{Anglais}{Tle A}{Chapter 1 : Family and Social Life}{EN02}{2 h}{Cette leçon de grammaire révise les zones à problèmes les plus fréquentes chez les francophones : accord sujet-verbe, emploi des temps (present perfect vs past simple), articles, prépositions et structures à double possibilité. À travers l'observation d'exemples, des règles synthétiques et des exercices de transformation, l'élève consolide ses acquis en vue du Bac.
\vspace{4pt}
\begin{multicols}{2}\small
\begin{itemize}
\item À la fin de cette leçon, je sais identifier les erreurs grammaticales fréquentes dans une phrase ou un court texte
\item À la fin de cette leçon, je sais appliquer correctement l'accord sujet-verbe, y compris avec les sujets difficiles (everyone, neither... nor, les noms collectifs)
\item À la fin de cette leçon, je sais choisir entre present perfect et past simple selon le contexte temporel
\item À la fin de cette leçon, je sais employer correctement les articles a/an/the et l'article zéro
\item À la fin de cette leçon, je sais utiliser les prépositions de temps et de lieu (in, on, at, for, since, by) sans les confondre
\item À la fin de cette leçon, je sais éviter les erreurs induites par le français (faux amis, traduction mot à mot, double négation)
\end{itemize}\end{multicols}}

\begin{sommaire}
\begin{multicols}{2}
\begin{enumerate}
\item Observation et diagnostic
\item Subject-verb agreement
\item Tenses: present perfect vs past simple
\item Articles and prepositions
\item French interference and false friends
\item Application and transformations
\item Activité d'intégration
\item Méthodes et exercices résolus
\item Exercices
\item Corrigés des exercices
\item Fiche bilan
\end{enumerate}
\end{multicols}
\end{sommaire}

\begin{objectifs}
\begin{itemize}
\item À la fin de cette leçon, je sais identifier les erreurs grammaticales fréquentes dans une phrase ou un court texte
\item À la fin de cette leçon, je sais appliquer correctement l'accord sujet-verbe, y compris avec les sujets difficiles (everyone, neither... nor, les noms collectifs)
\item À la fin de cette leçon, je sais choisir entre present perfect et past simple selon le contexte temporel
\item À la fin de cette leçon, je sais employer correctement les articles a/an/the et l'article zéro
\item À la fin de cette leçon, je sais utiliser les prépositions de temps et de lieu (in, on, at, for, since, by) sans les confondre
\item À la fin de cette leçon, je sais éviter les erreurs induites par le français (faux amis, traduction mot à mot, double négation)
\item À la fin de cette leçon, je sais transformer des phrases (affirmative → négative/interrogative, actif → passif, discours direct → indirect)
\item À la fin de cette leçon, je sais rédiger un court paragraphe grammaticalement correct sur la vie familiale et sociale
\end{itemize}
\end{objectifs}

\begin{prerequis}
\begin{itemize}
\item Les temps de base étudiés en Première (present simple, present continuous, past simple, future)
\item Les auxiliaires do/does/did et have/has
\item Les articles définis et indéfinis (notions de Seconde)
\item Le discours direct et indirect (notions de Première)
\item Le vocabulaire de la famille et de la vie sociale (Chapter 1, Lesson 1)
\end{itemize}
\end{prerequis}

\newpage

\coursec{Observation et diagnostic}

\begin{experience}[La composition de Manka'a]
Pendant l'examen blanc, Manka'a, élève de Tle A à Bamenda, relit sa composition. Elle a écrit :

\eng{My father have been living in Douala since ten years. He don't like when peoples arrives late.}

Son professeur a souligné \textbf{trois erreurs en rouge}. Pourtant, Manka'a connaît les règles ! Elle sait que \eng{my father} est une troisième personne du singulier, elle sait que \eng{since} se construit avec un point de départ, et elle sait que \eng{people} est déjà pluriel. Mais sous la pression de l'examen, son cerveau francophone a pris le dessus : « mon père \emph{ont} vécu » non, « il \emph{ne} aime pas » non... les réflexes du français se glissent partout.

Comme elle, la plupart des candidats au Bac perdent des points non pas sur des notions inconnues, mais sur des \cle{zones à problèmes} récurrentes : \textbf{accord sujet-verbe}, \textbf{concordance des temps}, \textbf{articles}, \textbf{prépositions} et \textbf{faux amis}. Cette leçon est une session de « réparation ciblée » : nous reprenons chaque point faible, nous le diagnostiquons, puis nous le corrigeons définitivement.

\textbf{Problématique :} comment repérer, classer et corriger systématiquement les erreurs de grammaire qui reviennent dans les copies des candidats francophones ?
\end{experience}

\courssub{Un texte truffé d'erreurs : à vous de jouer les correcteurs}

Lisons d'abord un court texte écrit par un élève de Première sur la vie familiale au Cameroun. Ce texte contient \textbf{huit erreurs typiques}. Lisez-le attentivement et soulignez tout ce qui vous semble incorrect, sans encore chercher à corriger.

\begin{textebox}[A text to correct — « Family life in Yaoundé »]
My uncle and his wife lives in Yaoundé since 2015. They have three childs. The last saturday, they have visited my grandmother in the village. She was very happy because she doesn't see them since long time. My uncle is a teacher in a secondary school and my aunt sell vegetables at the Mokolo market. Their childrens are very intelligent: the older is doctor in Buea and the two others are still at school. I am agree with my mother when she says that family is the most important thing in the life.
\end{textebox}

\begin{definition}[Glossaire]
\begin{itemize}
\item \eng{uncle} (nom) : oncle ; \eng{aunt} (nom) : tante.
\item \eng{grandmother} (nom) : grand-mère.
\item \eng{to sell} (verbe) : vendre.
\item \eng{vegetables} (nom pluriel) : légumes.
\item \eng{market} (nom) : marché.
\item \eng{still} (adverbe) : encore, toujours.
\end{itemize}
\end{definition}

\courssub{Repérage guidé : souligner avant de corriger}

Avant toute correction, il faut d'abord \emph{voir} l'erreur. C'est une compétence en soi, et c'est exactement celle que l'examinateur évalue dans votre propre production écrite. Voici les huit passages fautifs du texte, tels qu'un correcteur les aurait soulignés :

\begin{enumerate}
\item \eng{My uncle and his wife \textbf{lives}}
\item \eng{\textbf{since} 2015} (avec le présent simple \eng{lives})
\item \eng{three \textbf{childs}}
\item \eng{\textbf{The} last saturday}
\item \eng{they \textbf{have visited}}
\item \eng{she \textbf{doesn't see} them \textbf{since} long time}
\item \eng{my aunt \textbf{sell}}
\item \eng{I \textbf{am agree}}
\end{enumerate}

Avez-vous tout trouvé ? Si vous en avez repéré au moins six, votre « radar à erreurs » fonctionne bien. Sinon, pas de panique : c'est précisément l'objectif de cette leçon de le régler.

\begin{methode}[Diagnostiquer une erreur en quatre étapes]
Face à une phrase suspecte, procédez toujours dans le même ordre :
\begin{enumerate}
\item \textbf{Identifier} le mot ou le groupe de mots fautif : soulignez-le précisément (pas la phrase entière).
\item \textbf{Nommer} la règle concernée : est-ce un problème d'accord sujet-verbe ? de temps ? d'article ? de préposition ? de vocabulaire ?
\item \textbf{Reformuler} correctement la phrase complète, à voix haute ou par écrit.
\item \textbf{Vérifier} en traduisant en français : si la traduction « sonne » comme la phrase fautive, méfiance — c'est souvent le signe d'une interférence du français.
\end{enumerate}
\end{methode}

\courssub{Classement des erreurs par catégories}

Appliquons la méthode. Chaque erreur du texte appartient à l'une de nos cinq catégories :

\begin{center}
\begin{tabularx}{\linewidth}{|X|X|X|}
\hline
\rowcolor{popblueL} Erreur dans le texte & Correction & Catégorie \\
\hline
\eng{My uncle and his wife lives} & \eng{My uncle and his wife live} & Accord sujet-verbe (sujet pluriel) \\
\hline
\eng{lives ... since 2015} & \eng{have lived ... since 2015} & Temps (present perfect + \eng{since}) \\
\hline
\eng{three childs} & \eng{three children} & Pluriel irrégulier (lexique/forme) \\
\hline
\eng{The last saturday} & \eng{last Saturday} & Article + majuscule des jours \\
\hline
\eng{they have visited} (avec \eng{last Saturday}) & \eng{they visited} & Temps (past simple + moment précis) \\
\hline
\eng{she doesn't see them since long time} & \eng{she hasn't seen them for a long time} & Temps + préposition (\eng{since}/\eng{for}) \\
\hline
\eng{my aunt sell} & \eng{my aunt sells} & Accord sujet-verbe (3\textsuperscript{e} pers. du sing., -s) \\
\hline
\eng{I am agree} & \eng{I agree} & Faux ami / interférence du français (« je suis d'accord ») \\
\hline
\end{tabularx}
\end{center}

\begin{exemplebox}[Le texte corrigé]
\eng{My uncle and his wife \textbf{have lived} in Yaoundé since 2015. They have three \textbf{children}. \textbf{Last Saturday}, they \textbf{visited} my grandmother in the village. She was very happy because she \textbf{hadn't seen} them \textbf{for} a long time. My uncle is a teacher in a secondary school and my aunt \textbf{sells} vegetables at the Mokolo market. Their \textbf{children} are very intelligent: the \textbf{eldest} is \textbf{a} doctor in Buea and the two others are still at school. I \textbf{agree} with my mother when she says that family is the most important thing in \textbf{life}.}

« Mon oncle et sa femme vivent à Yaoundé depuis 2015. Ils ont trois enfants. Samedi dernier, ils ont rendu visite à ma grand-mère au village. Elle était très heureuse car elle ne les avait pas vus depuis longtemps. Mon oncle est enseignant dans un lycée et ma tante vend des légumes au marché Mokolo. Leurs enfants sont très intelligents : l'aîné est médecin à Buea et les deux autres sont encore à l'école. Je suis d'accord avec ma mère quand elle dit que la famille est la chose la plus importante dans la vie. »
\end{exemplebox}

\courssub{Le constat : des erreurs qui suivent des schémas}

Regardez bien le tableau : aucune de ces erreurs n'est \cle{aléatoire}. Chacune obéit à un schéma prévisible, et ce schéma a presque toujours la même origine : \textbf{l'influence du français}.

\begin{itemize}
\item « Je suis d'accord » → \eng{I am agree} : le français utilise l'auxiliaire \emph{être}, l'anglais utilise un verbe ordinaire, \eng{to agree}.
\item « Il habite à Yaoundé depuis 2015 » → \eng{he lives ... since 2015} : le français utilise le présent, l'anglais exige le present perfect.
\item « Samedi dernier, ils ont visité » → \eng{they have visited} : le passé composé français ressemble au present perfect anglais, mais l'anglais exige le past simple dès qu'un moment précis du passé est mentionné.
\item « Ma tante vend » → \eng{my aunt sell} : le français marque la 3\textsuperscript{e} personne à l'oreille mais pas toujours à l'écrit ; l'anglais exige le \textbf{-s} du présent simple, sans exception.
\end{itemize}

\begin{attention}[Le piège de la ressemblance]
Le danger n° 1 pour un francophone n'est pas ce qu'il ignore : c'est ce qui \emph{ressemble} au français. Le passé composé et le present perfect ont la même forme apparente (auxiliaire + participe), mais pas les mêmes emplois. De même, \eng{to be agree}, \eng{to have 20 years} ou \eng{to assist to a meeting} sont des calques du français. Règle d'or : \textbf{quand une phrase anglaise vous semble « trop facile » parce qu'elle calque le français, vérifiez-la deux fois.}
\end{attention}

\begin{savaistu}[Ce que disent les correcteurs du Bac]
Les jurys et correcteurs du Baccalauréat, session après session, font le même constat : plus de 60 \% des erreurs de grammaire relevées dans les copies des candidats francophones portent sur seulement \textbf{cinq points} — exactement ceux de cette leçon : l'accord sujet-verbe, l'emploi des temps (surtout present perfect contre past simple), les articles, les prépositions et les interférences du français. Autrement dit, un candidat qui maîtrise ces cinq zones élimine la majorité de ses fautes potentielles. C'est un investissement à très haut rendement pour l'examen !
\end{savaistu}

\courssub{Les cinq zones à problèmes qui structurent la leçon}

Voici la carte de notre parcours de réparation. Chaque branche correspond à une section de la leçon :

\begin{popfigure}
\begin{tikzpicture}[mindmap, grow cyclic, every node/.style={concept, text=white, align=center},
root concept/.append style={concept color=popink, font=\bfseries},
level 1/.append style={level distance=3.2cm, sibling angle=72},
level 1 concept/.append style={font=\small\bfseries}]
\node[root concept, align=center]{Problem\\areas}
  child[concept color=popblue] { node[align=center]{Subject-verb\\agreement} }
  child[concept color=poporange] { node[align=center]{Tenses\\present perfect\\vs past simple} }
  child[concept color=popgreen] { node[align=center]{Articles\\a / an / the / Ø} }
  child[concept color=poppurple] { node[align=center]{Prepositions\\since / for / in...} }
  child[concept color=poppink] { node[align=center]{French\\interference\\false friends} };
\end{tikzpicture}
\legende{Carte mentale des cinq zones à problèmes de la leçon}
\end{popfigure}

\begin{aretenir}[Ce qu'il faut retenir de cette première étape]
\begin{itemize}
\item Les erreurs des candidats francophones ne sont \textbf{pas aléatoires} : elles suivent des schémas liés à l'influence du français.
\item Le diagnostic se fait en quatre étapes : \textbf{identifier → nommer la règle → reformuler → vérifier en traduisant}.
\item Cinq zones concentrent l'essentiel des fautes : \cle{accord sujet-verbe}, \cle{temps}, \cle{articles}, \cle{prépositions}, \cle{interférences du français}.
\item Savoir \emph{repérer} une erreur est aussi important que savoir la corriger : c'est la même compétence, appliquée à sa propre copie.
\end{itemize}
\end{aretenir}

\begin{exoresolu}[Diagnostic complet d'une phrase]
\textbf{Énoncé.} Voici une phrase extraite d'une copie de Tle A :

\eng{My mother work in a hospital since 2010 and she have three childs.}

Appliquez la méthode en quatre étapes pour diagnostiquer et corriger toutes les erreurs.

\tcblower
\textbf{Solution détaillée.}

\textbf{Étape 1 — Identifier.} Trois groupes sont fautifs : \eng{My mother work}, \eng{work ... since 2010}, \eng{three childs}.

\textbf{Étape 2 — Nommer la règle.}
\begin{itemize}
\item \eng{My mother work} : problème d'\textbf{accord sujet-verbe} — à la 3\textsuperscript{e} personne du singulier, le présent simple prend un \textbf{-s}.
\item \eng{work ... since 2010} : problème de \textbf{temps} — \eng{since} + point de départ exige le \textbf{present perfect} (action commencée dans le passé, encore vraie).
\item \eng{three childs} : problème de \textbf{pluriel} — \eng{child} a un pluriel irrégulier : \eng{children}.
\end{itemize}

\textbf{Étape 3 — Reformuler.}

\eng{My mother \textbf{has worked} in a hospital since 2010 and she \textbf{has} three \textbf{children}.}

\textbf{Étape 4 — Vérifier en traduisant.} « Ma mère travaille dans un hôpital depuis 2010 et elle a trois enfants. » Le français utilise le présent (« travaille ») : c'est exactement le piège qui a produit l'erreur \eng{work}. La traduction confirme donc que la phrase fautive était un calque du français — et que notre correction avec le present perfect est bien nécessaire.
\end{exoresolu}

\begin{exercice}[À vous de diagnostiquer]
Pour chaque phrase, soulignez l'erreur, nommez la catégorie (accord, temps, article, préposition, faux ami), puis corrigez.

\begin{enumerate}
\item \eng{My parents has been married for twenty years.}
\item \eng{Last year, we have spent our holidays in Kribi.}
\item \eng{She is the more intelligent girl of the class.}
\item \eng{I am living in Bamenda since five years.}
\item \eng{He is agree with his father about the family meeting.}
\end{enumerate}
\end{exercice}

\begin{corrige}[Exercice « À vous de diagnostiquer »]
\begin{enumerate}
\item \eng{My parents \textbf{have} been married for twenty years.} — Accord sujet-verbe : sujet pluriel (\eng{my parents}) → auxiliaire \eng{have}.
\item \eng{Last year, we \textbf{spent} our holidays in Kribi.} — Temps : \eng{last year} est un moment précis et révolu → past simple, jamais le present perfect.
\item \eng{She is the \textbf{most} intelligent girl \textbf{in} the class.} — Superlatif (adjectif long → \eng{the most}) et préposition (\eng{in the class}).
\item \eng{I \textbf{have been living} in Bamenda \textbf{for} five years.} — Temps + préposition : durée qui continue → present perfect ; \eng{for} + durée, \eng{since} + point de départ.
\item \eng{He \textbf{agrees} with his father about the family meeting.} — Faux ami : « il est d'accord » ne se traduit pas par \eng{he is agree} ; \eng{to agree} est un verbe ordinaire (avec le -s de la 3\textsuperscript{e} personne).
\end{enumerate}
\end{corrige}

Maintenant que notre « radar à erreurs » est réglé et que les cinq zones à problèmes sont identifiées, passons à la première d'entre elles : l'accord sujet-verbe.

\coursec{Subject-verb agreement}

Dans la section précédente, nous avons observé un texte et repéré plusieurs erreurs typiques des élèves francophones. La première famille d'erreurs — et la plus fréquente à l'examen — concerne l'\cle{accord sujet-verbe} (\eng{subject-verb agreement}). C'est le point de grammaire que nous allons réviser maintenant, en profondeur, car il conditionne la correction de toutes vos productions écrites et orales.

\courssub{Observation : le verbe change-t-il selon son sujet ?}

\begin{textebox}[A family in Bafoussam]
The Nguemo family lives in Bafoussam. Everybody in the neighbourhood knows them because the father is a teacher and the mother sells vegetables at the market. Each of their children attends a different school. The news about their eldest son is wonderful: he has won a scholarship to study in Buea. Mathematics is his favourite subject, and physics is easy for him too. People in the street are proud of him. The police were even invited to the celebration party, because the uncle is a police officer. Neither his brothers nor his sister was jealous: everyone was happy. There is a lot of joy in the house, and there are many visitors every evening.
\end{textebox}

\textbf{Glossaire.} \eng{neighbourhood} (n., quartier) ; \eng{scholarship} (n., bourse d'études) ; \eng{proud} (adj., fier) ; \eng{jealous} (adj., jaloux) ; \eng{joy} (n., joie).

\textbf{Questions d'observation.}
\begin{enumerate}
\item Soulignez tous les verbes conjugués au present simple. Quels sujets prennent un verbe en \eng{-s} ?
\item Pourquoi dit-on \eng{people are} mais \eng{the news is} ?
\item Dans \eng{Neither his brothers nor his sister was jealous}, avec quel nom le verbe s'accorde-t-il ?
\end{enumerate}

Observation : en anglais, le verbe change de forme selon son sujet, mais beaucoup moins qu'en français. Au present simple, une seule forme change : la troisième personne du singulier, qui prend un \cle{-s}. Toute la difficulté consiste donc à savoir si le sujet est \emph{singulier} ou \emph{pluriel} — et l'anglais et le français ne sont pas toujours d'accord sur ce point !

\courssub{La règle de base : le verbe s'accorde avec son sujet en nombre}

\begin{propriete}[Règle fondamentale de l'accord sujet-verbe]
En anglais, le verbe s'accorde en \cle{nombre} avec son sujet :
\begin{itemize}
\item sujet \textbf{singulier} $\rightarrow$ verbe singulier : au present simple, le verbe prend \textbf{-s} (ou \textbf{-es}) à la 3e personne du singulier (\eng{he works}, \eng{she watches}, \eng{it rains}) ;
\item sujet \textbf{pluriel} $\rightarrow$ verbe pluriel : le verbe reste à la forme de base, sans \textbf{-s} (\eng{they work}, \eng{the children play}).
\end{itemize}
Aux autres temps, l'accord est surtout visible avec \eng{be} (\eng{is/are}, \eng{was/were}) et \eng{have} (\eng{has/have}).
\end{propriete}

\begin{center}
\begin{tabular}{|l|l|l|}
\hline
\rowcolor{popblueL} Personne & Present simple & Avec \eng{be} / \eng{have} \\
\hline
I / you / we / they & work, play, go & am / are ; have \\
\hline
he / she / it & work\textbf{s}, play\textbf{s}, go\textbf{es} & is ; has \\
\hline
\end{tabular}
\end{center}

\begin{exemplebox}[Le -s de la 3e personne : ne jamais l'oublier]
\eng{My father works in a bank in Douala.} « Mon père travaille dans une banque à Douala. »

\eng{She watches television every evening.} « Elle regarde la télévision tous les soirs. »

\eng{The bus leaves at six o'clock.} « Le car part à six heures. »

\eng{He has two sisters.} « Il a deux sœurs. »
\end{exemplebox}

\begin{attention}[L'erreur n° 1 des francophones : oublier le -s]
En français, le « s » de « il travaille\emph{s} » ne se prononce pas ; en anglais, le \eng{-s} de \eng{he works} se prononce (« ouèrk\emph{s} ») et \textbf{s'écrit obligatoirement}. Écrire \eng{my father work} est une faute grave, immédiatement sanctionnée à l'examen. Prononcez le -s à voix haute quand vous relisez : votre oreille vous aidera.
\end{attention}

\courssub{Les sujets indéfinis singuliers : everyone, everybody, someone, nobody, each, every}

\begin{propriete}[Les indéfinis en every-, some-, no- sont singuliers]
Les pronoms et déterminants indéfinis \eng{everyone}, \eng{everybody}, \eng{someone}, \eng{somebody}, \eng{no one}, \eng{nobody}, \eng{anyone}, \eng{anybody}, ainsi que \eng{each} et \eng{every + nom singulier}, sont toujours suivis d'un verbe au \textbf{singulier}.
\end{propriete}

\begin{exemplebox}[Exemples traduits]
\eng{Everybody in my village knows my father.} « Tout le monde dans mon village connaît mon père. »

\eng{Someone is waiting for you at the door.} « Quelqu'un t'attend à la porte. »

\eng{Nobody was absent yesterday.} « Personne n'était absent hier. »

\eng{Each student has a textbook.} « Chaque élève a un manuel. »

\eng{Every child likes football.} « Chaque enfant aime le football. »
\end{exemplebox}

Comparaison avec le français : la bonne nouvelle, c'est que le français fait pareil (« tout le monde \emph{connaît} », « chacun \emph{a} »). Le piège vient de l'idée de pluralité que ces mots évoquent : \eng{everybody} désigne plusieurs personnes, mais grammaticalement il est \textbf{singulier}. On dit donc \eng{everybody is}, jamais \eng{everybody are}.

\courssub{Neither... nor / either... or : l'accord avec le sujet le plus proche}

\begin{propriete}[Accord de proximité]
Avec \eng{neither... nor} (ni... ni) et \eng{either... or} (ou... ou), le verbe s'accorde avec le sujet \textbf{le plus proche} du verbe.
\end{propriete}

\begin{exemplebox}[Exemples traduits]
\eng{Neither my mother nor my brothers were at home.} « Ni ma mère ni mes frères n'étaient à la maison. » (accord avec \eng{brothers}, pluriel)

\eng{Neither my brothers nor my mother was at home.} « Ni mes frères ni ma mère n'était à la maison. » (accord avec \eng{mother}, singulier)

\eng{Either the teacher or the students are wrong.} « Ou le professeur ou les élèves ont tort. »

\eng{Either you or I am mistaken.} « Ou toi ou moi nous nous trompons. »
\end{exemplebox}

\courssub{Noms collectifs, noms toujours pluriels, there is / there are, et les cas piégeux en -s}

\begin{propriete}[Quatre cas particuliers à connaître par cœur]
\begin{enumerate}
\item \textbf{Noms collectifs} (\eng{family}, \eng{government}, \eng{team}, \eng{class}) : en anglais britannique, on emploie en général le \textbf{singulier} quand on considère le groupe comme une unité : \eng{My family is large.} « Ma famille est grande. »
\item \textbf{Noms toujours pluriels} : \eng{people} (les gens), \eng{police} (la police), \eng{cattle} (le bétail) prennent toujours un verbe au \textbf{pluriel} : \eng{The police are here.}
\item \textbf{There is / there are} : le verbe s'accorde avec le nom qui \textbf{suit} : \eng{There is a book on the table.} / \eng{There are three books on the table.}
\item \textbf{Faux pluriels en -s} : \eng{news} (les nouvelles), \eng{mathematics}, \eng{physics}, \eng{economics} sont \textbf{singuliers} malgré leur -s final : \eng{The news is good.}
\end{enumerate}
\end{propriete}

\begin{exemplebox}[Exemples traduits]
\eng{The police are investigating the case.} « La police enquête sur l'affaire. »

\eng{The news is good.} « Les nouvelles sont bonnes. »

\eng{Mathematics is difficult for many students.} « Les mathématiques sont difficiles pour beaucoup d'élèves. »

\eng{People are friendly in Bamenda.} « Les gens sont accueillants à Bamenda. »

\eng{There is a market near my school, and there are many shops around it.} « Il y a un marché près de mon école, et il y a beaucoup de boutiques autour. »
\end{exemplebox}

\begin{attention}[Erreur fréquente : « the people is »]
En français, « le peuple » et « la police » sont singuliers ; les élèves écrivent donc \eng{the people is}, \eng{the police is}. En anglais, \eng{people} et \eng{police} sont \textbf{toujours pluriels} : \eng{people are}, \eng{the police are}. À l'inverse, « les nouvelles \emph{sont} bonnes » est pluriel en français, mais \eng{news} est singulier : \eng{the news \textbf{is} good}. Méfiez-vous du -s final : il ne signale pas toujours un pluriel (\eng{news}, \eng{physics}), et son absence ne signale pas toujours un singulier (\eng{people}, \eng{police}).
\end{attention}

\begin{center}
\begin{tabularx}{\linewidth}{|X|X|X|}
\hline
\rowcolor{popblueL} Français & English & Remarque \\
\hline
Tout le monde connaît mon père. & Everybody knows my father. & singulier dans les deux langues \\
\hline
La police enquête. & The police are investigating. & pluriel en anglais \\
\hline
Les gens sont gentils. & People are kind. & pluriel en anglais \\
\hline
Les nouvelles sont bonnes. & The news is good. & singulier en anglais \\
\hline
Les mathématiques sont utiles. & Mathematics is useful. & singulier en anglais \\
\hline
Ma famille est grande. & My family is big. & collectif : singulier \\
\hline
Il y a des visiteurs. & There are visitors. & accord avec le nom qui suit \\
\hline
\end{tabularx}
\end{center}

\begin{popfigure}
\begin{tikzpicture}[node distance=0.35cm]
\node[draw=popblue, fill=popblueL, rounded corners, font=\bfseries, minimum width=6.2cm] (sing) {Singular subjects (verb + s / is / was)};
\node[draw=poporange, fill=poporangeL, rounded corners, font=\bfseries, minimum width=6.2cm, right=1.2cm of sing] (plur) {Plural subjects (verb / are / were)};
\node[draw=popblue, rounded corners, below=0.4cm of sing, align=left, text width=5.8cm] (singlist) {everyone / everybody\\ someone / nobody\\ each student / every child\\ the news / mathematics\\ the family (collective)};
\node[draw=poporange, rounded corners, below=0.4cm of plur, align=left, text width=5.8cm] (plurlist) {people\\ the police\\ cattle\\ the children\\ my brothers and I};
\end{tikzpicture}
\legende{Sujets singuliers et sujets pluriels : les cas à mémoriser.}
\end{popfigure}

\begin{exoresolu}[Choisissez la bonne forme du verbe]
Énoncé. Complétez avec la forme correcte du verbe entre parenthèses :
\begin{enumerate}
\item Everybody in the class \dots\ (be) ready for the test.
\item The police \dots\ (be) looking for the thief.
\item The news on the radio \dots\ (be) surprising.
\item Neither the students nor the teacher \dots\ (be) in the classroom.
\item There \dots\ (be) two mangoes on the table.
\item Each of the girls \dots\ (have) a new uniform.
\end{enumerate}
\tcblower
Solution détaillée :
\begin{enumerate}
\item \eng{Everybody in the class \textbf{is} ready for the test.} — \eng{everybody} est un indéfini singulier.
\item \eng{The police \textbf{are} looking for the thief.} — \eng{police} est toujours pluriel en anglais.
\item \eng{The news on the radio \textbf{is} surprising.} — \eng{news} est singulier malgré le -s final.
\item \eng{Neither the students nor the teacher \textbf{is} in the classroom.} — avec \eng{neither... nor}, accord avec le sujet le plus proche (\eng{the teacher}, singulier).
\item \eng{There \textbf{are} two mangoes on the table.} — \eng{there be} s'accorde avec le nom qui suit (\eng{two mangoes}, pluriel).
\item \eng{Each of the girls \textbf{has} a new uniform.} — \eng{each} est singulier, même suivi de \eng{of + pluriel}.
\end{enumerate}
\end{exoresolu}

\begin{exercice}[À vous de jouer]
Mettez les verbes entre parenthèses à la forme correcte :
\begin{enumerate}
\item Someone \dots\ (steal) my bicycle every year in this neighbourhood.
\item The government \dots\ (be) building a new school in our village.
\item Neither my father nor my uncles \dots\ (be) farmers.
\item Physics \dots\ (be) my favourite subject.
\item There \dots\ (be) a lot of people at the stadium yesterday.
\item Every student in Terminale \dots\ (work) hard for the baccalauréat.
\end{enumerate}
\end{exercice}

\begin{corrige}[Exercice « À vous de jouer »]
\begin{enumerate}
\item \eng{steals} — \eng{someone} est singulier ; n'oubliez pas le -s.
\item \eng{is} — nom collectif considéré comme une unité : singulier.
\item \eng{are} — accord avec le sujet le plus proche (\eng{my uncles}, pluriel).
\item \eng{is} — les matières en -ics sont singulières.
\item \eng{were} — \eng{people} est pluriel, et \eng{yesterday} impose le prétérit.
\item \eng{works} — \eng{every + nom singulier} : verbe singulier avec -s.
\end{enumerate}
\end{corrige}

\begin{aretenir}[L'essentiel de la section]
\begin{itemize}
\item Au present simple, le verbe prend \textbf{-s} à la 3e personne du singulier : \eng{he works}, \eng{she has}, \eng{it is}.
\item \eng{everyone, everybody, someone, nobody, each, every + nom} $\rightarrow$ verbe \textbf{singulier}.
\item \eng{people, police, cattle} $\rightarrow$ verbe \textbf{pluriel}.
\item \eng{news, mathematics, physics} $\rightarrow$ verbe \textbf{singulier} malgré le -s.
\item \eng{neither... nor / either... or} : accord avec le sujet le \textbf{plus proche}.
\item \eng{there is / there are} : accord avec le nom qui \textbf{suit}.
\end{itemize}
\end{aretenir}

Maintenant que l'accord sujet-verbe est consolidé, passons à un autre problème majeur des francophones : le choix entre le \emph{present perfect} et le \emph{past simple}.

\coursec{Tenses: present perfect vs past simple}

Dans la section précédente, nous avons vérifié que le verbe s'accorde toujours avec son sujet. Mais choisir la bonne \emph{forme} du verbe ne suffit pas : il faut aussi choisir le bon \emph{temps}. C'est précisément là que les francophones commettent le plus d'erreurs, car le français utilise un seul temps — le passé composé — là où l'anglais en exige deux : le \cle{past simple} et le \cle{present perfect}. Cette section vous apprend à les distinguer sans hésiter.

\courssub{Observation : un texte pour découvrir}

\begin{textebox}[A letter from Buea]
Dear Ngum,

How is life in Yaoundé? I have lived in Buea for three years now, and I still love it. I arrived here in 2021, when I started my studies at the university. Since then, I have made many friends and I have visited the mountain twice. Last Saturday, my friends and I climbed to the second hut — it was exhausting but wonderful! I have never seen such a beautiful landscape. Have you ever been to Buea? You came here once, didn't you? I think it was in 2019. I have just received your photos from Douala; they are great!

Write soon.

Your friend, Ali
\end{textebox}

\textbf{Glossaire :} \emph{to climb} (v., grimper), \emph{hut} (n., cabane, refuge), \emph{exhausting} (adj., épuisant), \emph{landscape} (n., paysage), \emph{once} (adv., une fois).

\textbf{Questions de compréhension :}
\begin{enumerate}
\item When did Ali arrive in Buea?
\item How long has he lived there?
\item When did he climb the mountain?
\item Has Ngum ever visited Buea? When?
\end{enumerate}

\textbf{Observation grammaticale.} Soulignons les verbes : \eng{I arrived here in 2021} (past simple, date précise) mais \eng{I have lived in Buea for three years} (present perfect, action qui continue). \eng{Last Saturday, we climbed} (past simple, moment révolu) mais \eng{I have never seen} (present perfect, expérience de vie non datée). Le choix du temps ne dépend pas du « degré de passé » : il dépend du \textbf{lien avec le présent}.

\courssub{Le past simple : un passé terminé et daté}

\begin{propriete}[Règle d'emploi du past simple]
Le \cle{past simple} exprime une action passée \textbf{terminée}, située à un \textbf{moment précis et révolu}, \textbf{coupée du présent}. Le moment est connu ou sous-entendu : on sait \emph{quand} cela s'est passé.

\textbf{Forme :} sujet + verbe à la 2\textsuperscript{e} forme (\eng{visited}, \eng{went}, \eng{saw}) ; négation et question avec \eng{did} + base verbale : \eng{Did you go? — I did not go.}

\textbf{Marqueurs typiques :} \eng{yesterday}, \eng{last week/month/year}, \eng{in 2010}, \eng{two days ago}, \eng{when I was young}, \eng{that morning}.
\end{propriete}

\begin{exemplebox}[Exemples traduits]
\eng{I visited my aunt last December.} « J'ai rendu visite à ma tante en décembre dernier. »

\eng{She passed her exam in 2023.} « Elle a réussi son examen en 2023. »

\eng{We watched the football match two hours ago.} « Nous avons regardé le match de football il y a deux heures. »

\eng{Did you see the President's speech yesterday?} « As-tu vu le discours du Président hier ? »
\end{exemplebox}

\courssub{Le present perfect : un passé lié au présent}

\begin{propriete}[Règle d'emploi du present perfect]
Le \cle{present perfect} exprime :
\begin{enumerate}
\item le \textbf{bilan présent} d'un passé \textbf{non daté} (expérience de vie, résultat visible maintenant) ;
\item une action \textbf{commencée dans le passé et qui continue} jusqu'au présent.
\end{enumerate}

\textbf{Forme :} sujet + \eng{have/has} + participe passé (\eng{have lived}, \eng{has gone}). Négation : \eng{have not (haven't)} ; question : \eng{Have you ever...?}

\textbf{Marqueurs typiques :} \eng{already} (déjà), \eng{just} (venir de), \eng{never} (jamais), \eng{ever} (déjà, interrogation), \eng{since} (depuis + point de départ), \eng{for} (depuis + durée), \eng{recently}, \eng{so far} (jusqu'ici), \eng{yet} (encore/déjà, en fin de phrase négative ou interrogative).
\end{propriete}

\begin{exemplebox}[Exemples traduits]
\eng{I have just visited my aunt.} « Je viens de rendre visite à ma tante. »

\eng{We have lived in Douala since 2018.} « Nous vivons à Douala depuis 2018. » (et nous y vivons encore)

\eng{She has already finished her homework.} « Elle a déjà fini ses devoirs. »

\eng{Have you ever eaten eru?} « As-tu déjà mangé de l'eru ? »

\eng{I have never travelled by plane.} « Je n'ai jamais voyagé en avion. »
\end{exemplebox}

\begin{popfigure}
\begin{center}
\begin{tikzpicture}[scale=1]
\draw[-{Stealth[length=3mm]}, thick, popdark] (0,0) -- (11,0);
\fill[popdark] (9.5,0) circle (0.08);
\node[below, popdark, font=\bfseries] at (9.5,-0.1) {NOW};
\node[below, popmuted] at (0.3,-0.1) {PAST};
% past simple : croix isolée datée
\draw[popink, very thick] (2.6,-0.18) -- (3.0,0.18);
\draw[popink, very thick] (2.6,0.18) -- (3.0,-0.18);
\node[above, popink, align=center, font=\small] at (2.8,0.25) {past simple\\ \eng{I went in 2022.}};
% present perfect : flèche jusqu'à NOW
\draw[-{Stealth[length=3mm]}, very thick, popgreen] (5.2,0.55) -- (9.4,0.55);
\fill[popgreen] (5.2,0.55) circle (0.07);
\node[above, popgreen, align=center, font=\small] at (7.3,0.65) {present perfect : jusqu'à maintenant\\ \eng{I have lived here since 2018.}};
\end{tikzpicture}
\legende{Le past simple est une croix isolée et datée dans le passé ; le present perfect relie le passé au présent.}
\end{center}
\end{popfigure}

\courssub{For et since : la durée et le point de départ}

\begin{propriete}[For + durée / since + point de départ]
Avec le present perfect, on distingue :
\begin{itemize}
\item \cle{for} + une \textbf{durée} : \eng{for ten years} (depuis dix ans), \eng{for two months}, \eng{for a long time} ;
\item \cle{since} + un \textbf{point de départ} : \eng{since 2015} (depuis 2015), \eng{since Monday}, \eng{since I was a child}.
\end{itemize}
\eng{They have known each other for ten years.} « Ils se connaissent depuis dix ans. »
\eng{They have known each other since 2015.} « Ils se connaissent depuis 2015. »
\end{propriete}

\begin{attention}[L'erreur n° 1 des francophones]
En français on dit « J'habite ici \textbf{depuis} 2018 » avec un \emph{présent}. En anglais, c'est impossible : avec \eng{since} ou \eng{for} et une action qui continue, il faut le \textbf{present perfect}.

Incorrect : \eng{I live here since 2018.}

Correct : \eng{I have lived here since 2018.}

De même : « Je l'attends depuis une heure » se dit \eng{I have been waiting for her for an hour}, jamais \eng{I wait for her since one hour}.
\end{attention}

\courssub{Have been ou have gone ?}

\begin{propriete}[Have been to / have gone to]
\begin{itemize}
\item \eng{He has \textbf{gone} to Buea.} = Il est parti à Buea et il \textbf{y est encore} (ou il est en route).
\item \eng{He has \textbf{been} to Buea.} = Il est allé à Buea et il \textbf{en est revenu} (expérience).
\end{itemize}
\eng{My father has gone to the market — he will be back at noon.} « Mon père est allé au marché — il rentrera à midi. »
\eng{I have been to Bamenda three times.} « Je suis allé à Bamenda trois fois. » (et j'en suis revenu)
\end{propriete}

\courssub{Méthode : comment choisir entre past simple et present perfect ?}

\begin{methode}[Choisir le bon temps en 3 étapes]
\begin{enumerate}
\item \textbf{Repérez le marqueur temporel} (s'il y en a un) : \eng{yesterday}, \eng{last week}, \eng{in 2020}, \eng{ago} $\rightarrow$ \textbf{past simple} ; \eng{since}, \eng{for}, \eng{just}, \eng{already}, \eng{never}, \eng{ever}, \eng{yet}, \eng{recently} $\rightarrow$ \textbf{present perfect}.
\item \textbf{Si aucun marqueur explicite} : demandez-vous si l'action est \textbf{datée mentalement} (je sais \emph{quand}) $\rightarrow$ past simple, ou si elle concerne mon \textbf{expérience / mon présent} (résultat, durée) $\rightarrow$ present perfect.
\item \textbf{Vérifiez la bascule} : dans une conversation, on commence souvent au present perfect (\eng{Have you ever...?}) et on bascule au past simple dès qu'on date l'événement (\eng{I went there in 2022}).
\end{enumerate}
\end{methode}

\courssub{Comparaison systématique avec le français}

Le passé composé français recouvre les deux emplois anglais. C'est le contexte (et surtout la présence d'une date révolue) qui impose le choix :

\begin{center}
\begin{tabularx}{\linewidth}{|X|X|X|}
\hline
\rowcolor{popblueL} Français & English & Remarque \\
\hline
J'ai vu ce film hier. & I \textbf{saw} that film yesterday. & Date révolue : past simple \\
\hline
J'ai déjà vu ce film. & I \textbf{have already seen} that film. & Pas de date : present perfect \\
\hline
Il a plu la semaine dernière. & It \textbf{rained} last week. & \eng{last} = past simple \\
\hline
Il a beaucoup plu cette année. & It \textbf{has rained} a lot this year. & Période non terminée : present perfect \\
\hline
Elle est partie il y a deux jours. & She \textbf{left} two days ago. & \eng{ago} impose le past simple \\
\hline
Elle vient de partir. & She \textbf{has just left}. & \eng{just} : present perfect \\
\hline
\end{tabularx}
\end{center}

\begin{aretenir}[Tableau récapitulatif des marqueurs temporels]
\begin{center}
\begin{tabular}{|l|l|}
\hline
\rowcolor{popblueL} Past simple (passé coupé du présent) & Present perfect (lien avec le présent) \\
\hline
yesterday & already, just \\
\hline
last night / week / year & never, ever \\
\hline
in 2010, in December & since + point de départ \\
\hline
two days ago & for + durée \\
\hline
when I was young & recently, so far \\
\hline
that morning, then & this week, today (non terminé), yet \\
\hline
\end{tabular}
\end{center}
\textbf{Exception absolue :} avec \cle{ago}, toujours le past simple, jamais le present perfect : \eng{I met her three years ago.} « Je l'ai rencontrée il y a trois ans. » — jamais \eng{I have met her three years ago.}
\end{aretenir}

\begin{textebox}[Mini-dialogue : experience vs date]
— Have you ever been to London?

— Yes, I have. I went there in 2022.

— Really? Did you enjoy it?

— Yes, it was fantastic! We visited the British Museum and we saw a musical.

— I have never travelled to England, but I have just bought my ticket for next July!
\end{textebox}

« — As-tu déjà été à Londres ? — Oui. J'y suis allé en 2022. — Vraiment ? Ça t'a plu ? — Oui, c'était fantastique ! Nous avons visité le British Museum et vu une comédie musicale. — Je ne suis jamais allé en Angleterre, mais je viens d'acheter mon billet pour juillet prochain ! »

\textbf{Observation :} on commence au present perfect pour évoquer l'\emph{expérience} (\eng{Have you ever been...?}) ; dès qu'on \textbf{date} l'événement, on bascule au past simple (\eng{I went there in 2022}). C'est la règle d'or de la conversation anglaise.

\begin{experience}[Jeu de rôle : "Life experiences interview"]
\textbf{Consigne :} En binôme, menez une interview de 2 minutes. L'élève A est journaliste, l'élève B est une personnalité camerounaise (footballeur, artiste, entrepreneur).
\begin{enumerate}
\item Le journaliste pose des questions au \textbf{present perfect} pour connaître le parcours (\eng{Have you ever played in a World Cup? Have you visited the North Region?}).
\item La personnalité répond, et \textbf{date} les événements précis au \textbf{past simple} (\eng{Yes, I played in 2022. I visited Maroua last year.}).
\item Changez de rôle.
\end{enumerate}
\textbf{Objectif :} Automatiser la bascule \eng{Have you ever...?} $\rightarrow$ \eng{I did it in...}.
\end{experience}

\begin{exoresolu}[Choisissez le bon temps]
Complétez avec le past simple ou le present perfect du verbe entre parenthèses.

1. I (visit) \dots\ my grandmother last Sunday.

2. She (live) \dots\ in Bafoussam since 2020.

3. They (never / see) \dots\ the Waza National Park.

4. We (buy) \dots\ this house ten years ago.

5. (you / ever / taste) \dots\ ndolé?

6. He (just / finish) \dots\ his composition.
\tcblower
\textbf{Solution détaillée :}

1. \eng{I \textbf{visited} my grandmother last Sunday.} — \eng{last Sunday} : moment révolu $\Rightarrow$ past simple.

2. \eng{She \textbf{has lived} in Bafoussam since 2020.} — \eng{since} + action qui continue $\Rightarrow$ present perfect.

3. \eng{They \textbf{have never seen} the Waza National Park.} — \eng{never}, expérience de vie non datée $\Rightarrow$ present perfect.

4. \eng{We \textbf{bought} this house ten years ago.} — \eng{ago} impose toujours le past simple.

5. \eng{\textbf{Have you ever tasted} ndolé?} — \eng{ever} : expérience non datée $\Rightarrow$ present perfect.

6. \eng{He \textbf{has just finished} his composition.} — \eng{just} (« venir de ») $\Rightarrow$ present perfect.
\end{exoresolu}

\begin{attention}[Pièges à éviter]
\begin{itemize}
\item Ne dites jamais \eng{I have seen him yesterday.} Dès qu'une date révolue apparaît (\eng{yesterday}, \eng{last week}, \eng{in 2019}, \eng{ago}), le present perfect est impossible : \eng{I saw him yesterday.}
\item Attention à la place de \eng{ago} : il se place \textbf{après} la durée — \eng{two days ago} (il y a deux jours), jamais \eng{ago two days}.
\item \eng{for} ne signifie pas « pendant » au sens d'une action terminée : « J'ai vécu à Garoua pendant trois ans » (c'est fini) = \eng{I lived in Garoua for three years} (past simple) ; mais « Je vis à Garoua depuis trois ans » = \eng{I have lived in Garoua for three years}.
\end{itemize}
\end{attention}

\begin{savaistu}[Pourquoi « ago » va à la fin ?]
En anglais, \eng{ago} est un adverbe qui se place toujours \emph{après} la durée qu'il modifie : \eng{a week ago}, \eng{five years ago}, \eng{long ago} (il y a longtemps). Le français fait l'inverse : « il y a \emph{une semaine} ». Ce petit mot vient de l'ancien anglais \emph{agan} (« passé »). Retenez aussi l'expression figée \eng{once upon a time, long, long ago} — « il était une fois, il y a très, très longtemps » — qui ouvre les contes anglais, exactement comme nos contes camerounais commencent par une formule rituelle.
\end{savaistu}

\begin{exercice}[À vous de jouer]
Mettez les verbes entre parenthèses au past simple ou au present perfect, puis traduisez les phrases 1, 3 et 5 en français.

1. My uncle (work) \dots\ in Douala for twenty years; he still works there.

2. The students (write) \dots\ their essays yesterday afternoon.

3. I (never / be) \dots\ to Limbe, but my sister (go) \dots\ there in 2021.

4. She (lose) \dots\ her keys this morning and she (not / find) \dots\ them yet.

5. We (know) \dots\ each other since primary school.

6. They (move) \dots\ to Bamenda three years ago.
\end{exercice}

\begin{corrige}[Exercice « À vous de jouer »]
1. \eng{My uncle \textbf{has worked} in Douala for twenty years; he still works there.} « Mon oncle travaille à Douala depuis vingt ans ; il y travaille encore. » (action qui continue : present perfect)

2. \eng{The students \textbf{wrote} their essays yesterday afternoon.} (\eng{yesterday} : past simple)

3. \eng{I \textbf{have never been} to Limbe, but my sister \textbf{went} there in 2021.} « Je ne suis jamais allé à Limbé, mais ma sœur y est allée en 2021. » (expérience non datée, puis date précise)

4. \eng{She \textbf{lost} her keys this morning and she \textbf{has not found} them yet.} (le matin est passé ; la clé reste introuvable maintenant)

5. \eng{We \textbf{have known} each other since primary school.} « Nous nous connaissons depuis l'école primaire. »

6. \eng{They \textbf{moved} to Bamenda three years ago.} (\eng{ago} : past simple obligatoire)
\end{corrige}

\coursec{Articles and prepositions}

\courssub{Transition et objectifs}

Après avoir révisé l'opposition entre \eng{present perfect} et \eng{past simple} — distinction cruciale entre un passé révolu et un passé lié au présent — nous abordons maintenant deux piliers de la phrase anglaise qui sont sources d'erreurs systématiques chez les francophones : les \cle{articles} (définis, indéfinis, zéro) et les \cle{prépositions} (temps, lieu, durée). En français, l'article est quasi obligatoire et les prépositions sont souvent imposées par le verbe (\eng{se marier \textbf{avec}}, \eng{répondre \textbf{à}}, \eng{entrer \textbf{dans}}). En anglais, l'article zéro exprime la généralité, et de nombreux verbes sont \cle{transitifs directs} sans préposition. Cette section vous donnera les repères visuels et les règles de décision pour ne plus hésiter.

\courssub{Les articles : a/an, the, zéro article}

\begin{textebox}[Observation : trois phrases, trois réalités]
My sister married \textbf{a} teacher last year. \\
The teacher she married is from Bafoussam. \\
\textbf{Teachers} work hard in Cameroon.
\end{textebox}

\begin{propriete}[Règle des articles — a/an, the, zéro]
\textbf{1. Article indéfini \cle{a / an} (singulier seulement)} \\
\emph{Forme :} \eng{a} devant son consonantique (\eng{a book}, \eng{a university} /juː/), \eng{an} devant son vocalique (\eng{an apple}, \eng{an hour} — \eng{h} muet, \eng{an honest man}). \\
\emph{Emploi :} désigne \textbf{un élément non identifié} parmi une classe, une première mention, une appartenance à une catégorie. \\
\eng{She married a teacher.} (un enseignant, n'importe lequel, première mention)

\textbf{2. Article défini \cle{the}} \\
\emph{Forme :} invariable, se place devant singulier ou pluriel. \\
\emph{Emploi :} désigne un \textbf{élément identifié}, connu de l'énonciateur et du destinataire (déjà mentionné, unique au monde, défini par une relative ou le contexte). \\
\eng{The teacher she married is from Bafoussam.} (celui-là précisément, identifié par la relative) \\
\eng{The sun rises in the east.} (unique) \\
\eng{The President of the Republic.} (fonction unique dans le contexte)

\textbf{3. Zéro article (\cle{Ø})} \\
\emph{Emploi :} \textbf{généralités} (noms pluriels comptables ou noms abstraits/non comptables sans déterminant), noms de repas, matières, langues, jours, mois, lieux institutionnels (\eng{school, hospital, church, prison, university, bed}) quand on y va pour leur fonction première. \\
\eng{Teachers work hard.} (les enseignants en général) \\
\eng{Life is short.} (la vie en général, abstrait) \\
\eng{I like \textbf{rice}.} (pas \eng{the rice} pour la généralité) \\
\eng{She goes to \textbf{school} by bike.} (fonction d'élève) mais \eng{She went to \textbf{the} school to meet the headmaster.} (bâtiment précis)
\end{propriete}

\begin{exemplebox}[Exemples traduits — Articles]
\eng{My sister married a teacher.} \\
« Ma sœur a épousé un enseignant. » (indéfini, première mention)

\eng{The teacher she married is from Bafoussam.} \\
« L'enseignant qu'elle a épousé est de Bafoussam. » (défini, identifié par la relative)

\eng{Teachers work hard.} \\
« Les enseignants travaillent dur. » (généralité, zéro article + pluriel)

\eng{She loves \textbf{music} and \textbf{dancing}.} \\
« Elle aime la musique et la danse. » (abstraits, généralité, zéro article)

\eng{We have \textbf{lunch} at 1 p.m.} \\
« Nous déjeunons à 13 h. » (repas, zéro article)
\end{exemplebox}

\begin{center}
\begin{tabularx}{\linewidth}{|X|X|X|}
\hline
\rowcolor{popblueL} \textbf{Contexte} & \textbf{Français} & \textbf{English} \\
\hline
Généralité (pluriel) & Les enfants aiment les jeux. & \eng{Children love games.} (\textbf{Ø}) \\
\hline
Généralité (abstrait) & La vie est courte. & \eng{Life is short.} (\textbf{Ø}) \\
\hline
Repas / Matière / Langue & J'aime le riz. / J'étudie l'anglais. & \eng{I like rice.} / \eng{I study English.} (\textbf{Ø}) \\
\hline
Première mention & J'ai vu \textbf{un} accident. & \eng{I saw \textbf{an} accident.} \\
\hline
Élément identifié & \textbf{L'}accident était grave. & \eng{\textbf{The} accident was serious.} \\
\hline
Unique / Connu de tous & Le soleil se lève à l'est. & \eng{\textbf{The} sun rises in \textbf{the} east.} \\
\hline
Lieu institutionnel (fonction) & Il est à l'école / à l'hôpital / au lit. & \eng{He is at \textbf{school} / in \textbf{hospital} / in \textbf{bed}.} (\textbf{Ø}) \\
\hline
Lieu institutionnel (bâtiment) & Il est \textbf{à l'}école (le bâtiment). & \eng{He is at \textbf{the} school.} \\
\hline
\end{tabularx}
\end{center}

\begin{attention}[Piège majeur : généralité et noms de repas/matières/langues]
En français, la généralité prend l'article défini (\eng{le, la, les}). En anglais : \textbf{zéro article}.
\begin{itemize}
    \item \eng{I like \textbf{chocolate}.} (pas \eng{the chocolate})
    \item \eng{She teaches \textbf{mathematics}.} (pas \eng{the mathematics})
    \item \eng{They speak \textbf{French}.} (pas \eng{the French})
    \item \eng{We have \textbf{breakfast} at 7.} (pas \eng{the breakfast})
\end{itemize}
\emph{Exception :} si le nom est particularisé (ex. \eng{The French spoken in Cameroon is slightly different}).
\end{attention}

\courssub{Prépositions de temps : at, on, in — du point à la période}

\begin{popfigure}
\begin{tikzpicture}[
    node distance=1.2cm,
    box/.style={draw=popblue, thick, rounded corners, fill=popblueL, minimum width=3.2cm, minimum height=2.2cm, align=center, font=\sffamily},
    arrow/.style={-Stealth, thick, poporange}
]
\node[box, align=center] (at){\textbf{AT}\\[0.3em] \cle{Point précis}\\[0.2em] \eng{at 7 o'clock}\\\eng{at noon / midnight}\\\eng{at the moment}};
\node[box, right=of at, align=center] (on){\textbf{ON}\\[0.3em] \cle{Jour / Date / Surface}\\[0.2em] \eng{on Monday}\\\eng{on 5th May}\\\eng{on Monday morning}\\\eng{on the table}};
\node[box, right=of on, align=center] (in){\textbf{IN}\\[0.3em] \cle{Période / Espace clos}\\[0.2em] \eng{in July / in 2024}\\\eng{in the morning}\\\eng{in Yaoundé}\\\eng{in the box}};
\draw[arrow] (at) -- (on) node[midway, above, font=\tiny, popdark] {plus large};
\draw[arrow] (on) -- (in) node[midway, above, font=\tiny, popdark] {plus large};
\end{tikzpicture}
\legende{Échelle de précision : AT (point) $\to$ ON (jour/date/surface) $\to$ IN (période/mois/année/espace clos).}
\end{popfigure}

\begin{propriete}[Prépositions de temps — Règles de choix]
\textbf{AT} : \textbf{heure précise} (\eng{at 7:30 a.m.}, \eng{at noon}, \eng{at midnight}, \eng{at the moment}, \eng{at night} — exception idiomatique).
\textbf{ON} : \textbf{jour, date, jour férié, partie de journée précise} (\eng{on Monday}, \eng{on 20th May}, \eng{on Christmas Day}, \eng{on Monday morning}, \eng{on the weekend} — UK ; US souvent \eng{on the weekend} ou \eng{over the weekend}).
\textbf{IN} : \textbf{mois, année, saison, partie de journée (sauf night), durée future} (\eng{in July}, \eng{in 2024}, \eng{in summer}, \eng{in the morning/afternoon/evening}, \eng{in two weeks} = dans deux semaines).
\end{propriete}

\begin{exemplebox}[Exemples traduits — Temps]
\eng{The wedding is \textbf{on} Saturday \textbf{at} 10 a.m. \textbf{in} December.} \\
« Le mariage a lieu samedi à 10 heures, en décembre. »

\eng{I usually study \textbf{in} the evening, but \textbf{on} Friday evening I go out.} \\
« J'étudie d'habitude le soir, mais vendredi soir je sors. »

\eng{The train arrives \textbf{at} midnight.} \\
« Le train arrive à minuit. »

\eng{We will finish the project \textbf{in} three days.} \\
« Nous terminerons le projet dans trois jours. »
\end{exemplebox}

\begin{attention}[Pièges fréquents « in the morning » mais « at night » ; « on Monday morning »]
\begin{itemize}
    \item \eng{in the morning / afternoon / evening} MAIS \eng{\textbf{at} night} (idiomatique).
    \item \eng{on Monday morning} (jour précis \textbf{prime} sur la partie de journée) — pas \eng{in Monday morning}.
    \item \eng{at the weekend} (UK) / \eng{on the weekend} (US) — les deux existent, suivez votre variété de référence.
    \item \eng{in time} (à temps, avant l'échéance) vs \eng{on time} (ponctuel, à l'heure prévue).
\end{itemize}
\end{attention}

\courssub{Prépositions de lieu : at, in, on — même logique}

\begin{propriete}[Prépositions de lieu — Règles de choix]
\textbf{AT} : \textbf{point précis}, adresse, événement, arrêt (\eng{at the market}, \eng{at the bus stop}, \eng{at 10 Main Street}, \eng{at the party}, \eng{at home}, \eng{at school} — fonction).
\textbf{ON} : \textbf{surface, ligne, rue, étage, rive} (\eng{on the table}, \eng{on the wall}, \eng{on the map}, \eng{on the first floor}, \eng{on the left}, \eng{on the river}, \eng{on Main Street}).
\textbf{IN} : \textbf{espace clos, volume, ville, pays, quartier, bâtiment (à l'intérieur)} (\eng{in Yaoundé}, \eng{in Cameroon}, \eng{in the classroom}, \eng{in the box}, \eng{in the car}, \eng{in the street} — vu comme un espace clos).
\end{propriete}

\begin{center}
\begin{tabularx}{\linewidth}{|X|X|X|}
\hline
\rowcolor{popblueL} \textbf{Préposition} & \textbf{Image mentale} & \textbf{Exemples (EN / FR)} \\
\hline
\textbf{AT} & Point sur une carte & \eng{at the market} (au marché) / \eng{at the bus stop} (à l'arrêt de bus) / \eng{at home} (à la maison) \\
\hline
\textbf{ON} & Surface / Ligne & \eng{on the table} (sur la table) / \eng{on the 2nd floor} (au 2e étage) / \eng{on the beach} (sur la plage) \\
\hline
\textbf{IN} & Volume / Espace clos & \eng{in Yaoundé} (à Yaoundé) / \eng{in the classroom} (dans la classe) / \eng{in the car} (dans la voiture) \\
\hline
\end{tabularx}
\end{center}

\begin{exemplebox}[Exemples traduits — Lieu]
\eng{Meet me \textbf{at} the entrance of the supermarket \textbf{in} Douala.} \\
« Retrouve-moi à l'entrée du supermarché à Douala. »

\eng{The picture is \textbf{on} the wall \textbf{in} the living room.} \\
« Le tableau est sur le mur dans le salon. »

\eng{She lives \textbf{on} Mvog-Ada Street \textbf{in} Yaoundé.} \\
« Elle habite rue Mvog-Ada à Yaoundé. »
\end{exemplebox}

\courssub{For, since, by, until — durée, point de départ, échéance}

\begin{propriete}[For / Since / By / Until — Distinctions]
\textbf{FOR} + \textbf{durée} (\eng{for two hours}, \eng{for a long time}, \eng{for ages}). Répond à \eng{How long?}. S'utilise avec \textbf{tous les temps}.
\textbf{SINCE} + \textbf{point de départ dans le passé} (\eng{since 2010}, \eng{since Monday}, \eng{since I arrived}). Exige généralement \textbf{present perfect} (ou past perfect) : \eng{I have known him since 2015.}
\textbf{BY} + \textbf{échéance (deadline)} = « au plus tard le/avant le » (\eng{by Monday}, \eng{by 5 p.m.}, \eng{by the time you arrive}). Action \textbf{terminée} avant ce moment.
\textbf{UNTIL / TILL} + \textbf{point final} = « jusqu'à » (continuité de l'action ou de l'état jusqu'à ce moment). \eng{We waited until midnight.} \eng{He works till 6 p.m.}
\end{propriete}

\begin{exemplebox}[Exemples traduits — For / Since / By / Until]
\eng{I have lived in Buea \textbf{for} ten years.} \\
« J'habite à Buea depuis dix ans. » (durée)

\eng{She has been a teacher \textbf{since} 2018.} \\
« Elle est enseignante depuis 2018. » (point de départ)

\eng{Please submit your assignment \textbf{by} Friday.} \\
« Veuillez rendre votre devoir au plus tard vendredi. » (échéance)

\eng{The shop stays open \textbf{until} 9 p.m.} \\
« Le magasin reste ouvert jusqu'à 21 h. » (continuité)
\end{exemplebox}

\begin{attention}[Erreur fréquente : \eng{for} vs \eng{since} avec le present perfect]
\begin{itemize}
    \item \eng{I have been here \textbf{for} two hours.} (durée = 2 heures)
    \item \eng{I have been here \textbf{since} 10 a.m.} (départ = 10 h)
    \item \textbf{Jamais} \eng{*I am here since two hours.} (présent + since = erreur). Utilisez \eng{I have been here for two hours.}
\end{itemize}
\end{attention}

\courssub{Verbes transitifs directs en anglais / prépositionnels en français}

C'est un point de \cle{transfert négatif} majeur. En français, le verbe appelle une préposition (\eng{à, avec, dans}) ; en anglais, le verbe est \textbf{transitif direct} (pas de préposition, COD immédiat).

\begin{propriete}[Verbes sans préposition en anglais (avec préposition en français)]
\begin{center}
\begin{tabular}{|l|l|l|}
\hline
\rowcolor{popblueL} \textbf{Verbe anglais} & \textbf{Construction FR} & \textbf{Exemple EN / FR} \\
\hline
\eng{marry someone} & se marier \textbf{avec} qqn & \eng{She married John.} / « Elle a épousé John. » \\
\hline
\eng{answer someone / something} & répondre \textbf{à} qqn / qqch & \eng{Answer the question.} / « Réponds à la question. » \\
\hline
\eng{phone / call someone} & téléphoner \textbf{à} qqn & \eng{I phoned my mother.} / « J'ai appelé ma mère. » \\
\hline
\eng{enter a place} & entrer \textbf{dans} un lieu & \eng{He entered the room.} / « Il est entré dans la pièce. » \\
\hline
\eng{reach a place} & arriver \textbf{à} / atteindre & \eng{We reached Yaoundé at noon.} / « Nous sommes arrivés à Yaoundé à midi. » \\
\hline
\eng{approach someone/something} & s'approcher \textbf{de} & \eng{The train approached the station.} / « Le train s'approchait de la gare. » \\
\hline
\eng{await someone/something} (formel) & attendre qqn/qqch & \eng{We await your reply.} / « Nous attendons votre réponse. » \\
\hline
\eng{resemble someone} & ressembler \textbf{à} qqn & \eng{He resembles his father.} / « Il ressemble à son père. » \\
\hline
\end{tabular}
\end{center}
\end{propriete}

\begin{attention}[Erreur classique : \eng{*marry WITH someone}]
\begin{itemize}
    \item \textbf{Faux :} \eng{*She married with John.} \quad \textbf{Juste :} \eng{She married John.}
    \item \textbf{Faux :} \eng{*He entered in the room.} \quad \textbf{Juste :} \eng{He entered the room.}
    \item \textbf{Faux :} \eng{*I phoned to her.} \quad \textbf{Juste :} \eng{I phoned her.}
    \item \textbf{Faux :} \eng{*Answer to me.} \quad \textbf{Juste :} \eng{Answer me.}
\end{itemize}
\emph{Astuce :} Apprenez le verbe \textbf{avec son objet direct} : \eng{marry someone}, \eng{enter somewhere}, \eng{phone someone}.
\end{attention}

\courssub{Exercice résolu : transformation et choix}

\begin{exoresolu}[Articles, prépositions, verbes transitifs]
\textbf{Énoncé :} Complétez avec \eng{a/an, the, Ø} (zéro article), \eng{at, on, in, for, since, by, until} ou rien (verbe transitif direct). Traduisez.

1. My brother is \_\_\_ teacher. He works \_\_\_ a secondary school \_\_\_ Yaoundé. \\
2. We have known each other \_\_\_ 2015. We met \_\_\_ a party \_\_\_ Christmas Day. \\
3. The meeting starts \_\_\_ 9 a.m. \_\_\_ Monday and finishes \_\_\_ 1 p.m. \\
4. She \_\_\_ (marry) \_\_\_ doctor last year. \_\_\_ doctor is from Bafoussam. \\
5. Please \_\_\_ (answer) \_\_\_ question \_\_\_ Friday. \\
6. I usually have \_\_\_ lunch \_\_\_ 1 p.m., but \_\_\_ Sundays I eat \_\_\_ noon.

\tcblower
\textbf{Solution détaillée :}

1. \eng{My brother is \textbf{a} teacher. He works \textbf{at} a secondary school \textbf{in} Yaoundé.} \\
« Mon frère est enseignant. Il travaille dans un lycée à Yaoundé. » \\
\emph{Explication :} \eng{a teacher} (catégorie, indéfini) ; \eng{at} (point précis : l'école comme lieu de travail) ; \eng{in} (ville = espace clos).

2. \eng{We have known each other \textbf{since} 2015. We met \textbf{at} a party \textbf{on} Christmas Day.} \\
« Nous nous connaissons depuis 2015. Nous nous sommes rencontrés à une fête le jour de Noël. » \\
\emph{Explication :} \eng{since} + point de départ (present perfect) ; \eng{at a party} (événement/point) ; \eng{on} + jour férié/date.

3. \eng{The meeting starts \textbf{at} 9 a.m. \textbf{on} Monday and finishes \textbf{by} 1 p.m.} \\
« La réunion commence à 9 h lundi et se termine au plus tard à 13 h. » \\
\emph{Explication :} \eng{at} (heure précise) ; \eng{on} (jour) ; \eng{by} (échéance = fin au plus tard à cette heure).

4. \eng{She \textbf{married a} doctor last year. \textbf{The} doctor is from Bafoussam.} \\
« Elle a épousé un médecin l'an dernier. Le médecin est de Bafoussam. » \\
\emph{Explication :} \eng{married} (transitif direct, pas de \eng{with}) + \eng{a} (1re mention) ; \eng{The} (identifié, déjà mentionné).

5. \eng{Please \textbf{answer the} question \textbf{by} Friday.} \\
« Veuillez répondre à la question d'ici vendredi. » \\
\emph{Explication :} \eng{answer} (transitif direct, pas de \eng{to}) + \eng{the} (question identifiée) ; \eng{by} (échéance).

6. \eng{I usually have \textbf{Ø} lunch \textbf{at} 1 p.m., but \textbf{on} Sundays I eat \textbf{at} noon.} \\
« Je déjeune d'habitude à 13 h, mais le dimanche je mange à midi. » \\
\emph{Explication :} \eng{Ø lunch} (repas, généralité) ; \eng{at} (heure) ; \eng{on Sundays} (jour répétitif) ; \eng{at noon} (heure précise).
\end{exoresolu}

\courssub{Exercices d'application}

\begin{exercice}[Entraînement — Articles et prépositions]
\textbf{A. Choisissez la forme correcte.}
1. I have been waiting for you \_\_\_ (since / for) two hours. \\
2. The conference is \_\_\_ (at / on / in) 10 a.m. \_\_\_ (at / on / in) Tuesday. \\
3. He \_\_\_ (married / married with) his childhood friend. \\
4. We arrived \_\_\_ (at / in / to) Douala \_\_\_ (at / on / in) midnight. \\
5. \_\_\_ (The / Ø) honesty is \_\_\_ (a / the / Ø) rare quality.

\textbf{B. Traduisez en anglais.}
1. Elle est entrée dans la salle sans frapper. \\
2. Réponds-moi tout de suite ! \\
3. J'étudie l'anglais depuis trois ans. \\
4. Le match commence à 16 h samedi. \\
5. Les enseignants sont patients. (généralité)

\textbf{C. Corrigez les erreurs (une par phrase).}
1. She is married with a lawyer. \\
2. I have been here since three weeks. \\
3. We discussed about the problem. \\
4. He entered in the house quietly. \\
5. The life is beautiful in Cameroon.
\end{exercice}

\begin{corrige}[Exercice — Articles et prépositions]
\textbf{A.}
1. \eng{for} (durée) — \eng{I have been waiting for you \textbf{for} two hours.}
2. \eng{at 10 a.m. \textbf{on} Tuesday} (heure précise + jour)
3. \eng{married} (transitif direct) — \eng{He \textbf{married} his childhood friend.}
4. \eng{in Douala \textbf{at} midnight} (ville = espace clos ; heure précise = at)
5. \eng{\textbf{Ø} honesty is \textbf{a} rare quality.} (abstrait généralité + indéfini singularité)

\textbf{B.}
1. \eng{She \textbf{entered the room} without knocking.}
2. \eng{\textbf{Answer me} right now!}
3. \eng{I \textbf{have studied English for} three years.} (ou \eng{since 2021} si année de référence)
4. \eng{The match starts \textbf{at 4 p.m. on} Saturday.}
5. \eng{\textbf{Teachers are patient.}} (zéro article + pluriel pour généralité)

\textbf{C.}
1. \eng{She is married \textbf{to} a lawyer.} (participe adjectival \eng{married to}) ou \eng{She \textbf{married} a lawyer.} (verbe acte)
2. \eng{I have been here \textbf{for} three weeks.}
3. \eng{We \textbf{discussed} the problem.} (\eng{discuss} = transitif direct)
4. \eng{He \textbf{entered the house} quietly.}
5. \eng{\textbf{Life is beautiful in Cameroon.}} (zéro article pour généralité)
\end{corrige}

\begin{savaistu}[Le \eng{the} « générique » vs zéro article]
Parfois \eng{the} + singulier exprime une généralité sur une espèce ou une invention : \eng{The elephant is a massive animal.} (L'éléphant en tant qu'espèce). \eng{The telephone changed communication.} Mais pour les noms abstraits, matières, repas, langues : \textbf{zéro article} obligatoire. \eng{The French} = les Français (peuple) ; \eng{French} = la langue.
\end{savaistu}

\begin{methode}[Comment choisir sa préposition de temps/lieu en 3 secondes]
\begin{enumerate}
    \item \textbf{Temps :} C'est une \textbf{heure exacte} ? $\to$ \eng{AT}. C'est un \textbf{jour/date} ? $\to$ \eng{ON}. C'est un \textbf{mois/année/saison/partie de journée/durée future} ? $\to$ \eng{IN}.
    \item \textbf{Lieu :} C'est un \textbf{point/adresse/événement} ? $\to$ \eng{AT}. C'est une \textbf{surface/ligne/étage/rue} ? $\to$ \eng{ON}. C'est un \textbf{espace clos/ville/pays/pièce} ? $\to$ \eng{IN}.
    \item \textbf{Durée/Échéance :} \textbf{Durée} ($\to$ \eng{FOR}) ; \textbf{Départ passé} ($\to$ \eng{SINCE} + perfect) ; \textbf{Deadline} ($\to$ \eng{BY}) ; \textbf{Continuité jusqu'à} ($\to$ \eng{UNTIL}).
\end{enumerate}
\end{methode}

\coursec{French interference and false friends}

Après avoir revu les articles et les prépositions, nous abordons maintenant un domaine où l'interférence du français est la plus sournoise : les \textbf{faux amis}, les \textbf{calques syntaxiques} et les \textbf{différences structurelles} qui transforment une phrase correcte en anglais en une erreur flagrante. Cette section est décisive pour l'épreuve écrite du Baccalauréat : chaque point perdu sur ces pièges classiques est un point offert aux correcteurs.

\courssub{Double negation : l'interdit absolu}

\begin{textebox}[Observation : la négation en anglais vs en français]
\textbf{French:} Je ne sais \textbf{rien}. / Je \textbf{ne} connais \textbf{personne}. / Je n'ai \textbf{jamais} vu ça.

\textbf{English:} I \textbf{don't know anything}. / I \textbf{don't know anyone}. / I \textbf{have never} seen that.
\end{textebox}

\begin{propriete}[Règle : une seule négation par proposition]
En anglais standard, \textbf{un seul marqueur de négation} suffit pour rendre la phrase négative. Ce marqueur peut être :
\begin{itemize}
    \item l'auxiliaire négatif : \eng{don't, doesn't, didn't, haven't, isn't, can't, won't, etc.}
    \item l'adverbe \eng{never, hardly, scarcely, rarely, seldom}
    \item le déterminant \eng{no, none, neither}
    \item le pronom \eng{nothing, nobody, no one, nowhere}
\end{itemize}
\textbf{Jamais} deux marqueurs négatifs dans la même proposition (sauf rhétorique populaire \eng{I ain't got nothing} — à proscrire à l'écrit scolaire).
\end{propriete}

\begin{center}
\begin{tabularx}{\linewidth}{|X|X|X|}
\hline
\rowcolor{popblueL} \textbf{Français (double négation)} & \textbf{Anglais INCORRECT (calque)} & \textbf{Anglais CORRECT} \\
\hline
Je ne sais rien. & I don't know nothing. & \eng{I don't know \textbf{anything}.} \\
Je ne vois personne. & I don't see nobody. & \eng{I don't see \textbf{anyone / anybody}.} \\
Je n'ai jamais fait ça. & I haven't never done that. & \eng{I \textbf{have never} done that.} (ou \eng{I haven't \textbf{ever} done that.}) \\
Il n'y a plus de pain. & There isn't no more bread. & \eng{There is \textbf{no} more bread.} / \eng{There isn't \textbf{any} more bread.} \\
\hline
\end{tabularx}
\end{center}

\begin{exemplebox}[Exemples traduits]
\eng{She didn't say anything.} « Elle n'a rien dit. » \\
\eng{We haven't received any news.} « Nous n'avons reçu aucune nouvelle. » \\
\eng{Nobody came to the meeting.} « Personne n'est venu à la réunion. » (sujet négatif = verbe affirmatif) \\
\eng{He hardly speaks English.} « Il parle à peine l'anglais. » (\eng{hardly} = à peine, presque pas ; déjà négatif)
\end{exemplebox}

\begin{attention}[Piège : \eng{anything / anyone / anywhere} vs \eng{nothing / nobody / nowhere}]
Après une négation (\eng{not, never, without}), on utilise la série en \textbf{-thing / -one / -where} : \eng{anything, anyone, anywhere}.
\begin{itemize}
    \item \eng{I don't have \textbf{any} money.} (pas \eng{no money} après \eng{don't})
    \item \eng{Without \textbf{any} hesitation.} (pas \eng{without no hesitation})
\end{itemize}
La série en \textbf{no-} (\eng{nothing, nobody, nowhere}) s'emploie \textbf{seule}, sans \eng{not} : \eng{Nothing happened.} / \eng{I saw nobody.}
\end{attention}

\courssub{Adjectifs invariables et position}

\begin{textebox}[Observation]
\textbf{French:} Des enfants \textbf{grands}. / Une maison \textbf{bleue}. / Des problèmes \textbf{difficiles}.

\textbf{English:} \textbf{Tall} children. / A \textbf{blue} house. / \textbf{Difficult} problems.
\end{textebox}

\begin{propriete}[Deux règles d'or]
\begin{enumerate}
    \item \textbf{Invariabilité totale :} L'adjectif anglais ne s'accorde \textbf{jamais} en genre ni en nombre. \eng{A tall boy}, \eng{tall boys}, \eng{a tall girl}, \eng{tall girls}.
    \item \textbf{Position anteposée :} L'adjectif qualificatif se place \textbf{devant} le nom qu'il qualifie (sauf cas particuliers : \eng{something new}, \eng{the president elect}, adjectifs en \eng{-able/-ible} postposés pour le sens).
\end{enumerate}
\end{propriete}

\begin{exemplebox}[Exemples traduits]
\eng{She bought three beautiful dresses.} « Elle a acheté trois belles robes. » \\
\eng{The old man walked slowly.} « Le vieil homme marchait lentement. » \\
\eng{We need fresh ideas.} « Nous avons besoin d'idées nouvelles. »
\end{exemplebox}

\begin{attention}[Erreur fréquente : \eng{*talls children}, \eng{*blues eyes}, \eng{*differents ideas}]
\begin{itemize}
    \item \eng{*My parents are very poors.} $\rightarrow$ \eng{My parents are very \textbf{poor}.}
    \item \eng{*She has blues eyes.} $\rightarrow$ \eng{She has \textbf{blue} eyes.}
    \item \eng{*We discussed differents topics.} $\rightarrow$ \eng{We discussed \textbf{different} topics.}
\end{itemize}
\textbf{Astuce :} Si vous hésitez, rappelez-vous que l'adjectif anglais se comporte comme un adjectif \textbf{invariable} en français (\eng{sympa, super, génial} — on ne dit pas \eng{*des gens sympas} au féminin pluriel).
\end{attention}

\courssub{Les faux amis classiques : tableau de survie}

\begin{center}
\begin{tabularx}{\linewidth}{|X|X|X|}
\hline
\rowcolor{popblueL} \textbf{Français} & \textbf{Faux ami anglais} & \textbf{Traduction correcte} \\ \hline
Actuellement & \eng{actually (= en fait)} & \eng{currently} \\ \hline
Assister à & \eng{to assist (= aider)} & \eng{to attend} \\ \hline
Sensible (émotif) & \eng{sensible (= raisonnable)} & \eng{sensitive} \\ \hline
Éventuellement & \eng{eventually (= finalement)} & \eng{possibly / maybe} \\ \hline
Prétendre & \eng{to pretend (= faire semblant)} & \eng{to claim} \\ \hline
Demander & \eng{to demand (= exiger)} & \eng{to ask (for)} \\ \hline
Blesser & \eng{to bless (= bénir)} & \eng{to hurt / to injure} \\ \hline
Une librairie & \eng{a library (= bibliothèque)} & \eng{a bookshop / bookstore} \\ \hline
Un collège & \eng{a college (= université)} & \eng{a secondary school} \\ \hline
Une location & \eng{a location (= lieu)} & \eng{a rental} \\ \hline
Un journal & \eng{a journal (= revue)} & \eng{a newspaper} \\ \hline
Une figure & \eng{a figure (= chiffre, silhouette)} & \eng{a face / a character} \\ \hline
Rester & \eng{to rest (= se reposer)} & \eng{to stay / to remain} \\ \hline
Passer un examen & \eng{to pass an exam (= réussir)} & \eng{to take / to sit an exam} \\ \hline
\end{tabularx}
\end{center}
\textit{Les 14 faux amis les plus piégeants pour les francophones : à connaître par cœur pour le Bac.}

\begin{exemplebox}[Focus sur les trois « champions » du Bac]
\eng{He \textbf{actually} lives in Bamenda.} « En fait, il vit à Bamenda. » \\
\eng{He \textbf{currently} lives in Bamenda.} « Il vit \textbf{actuellement} à Bamenda. » \\[4pt]
\eng{The nurse \textbf{assisted} the doctor.} « L'infirmière a \textbf{aidé} le médecin. » \\
\eng{Many parents \textbf{attended} the meeting.} « De nombreux parents ont \textbf{assisté à} la réunion. » \\[4pt]
\eng{Be \textbf{sensible} and wear a helmet.} « Sois \textbf{raisonnable} et porte un casque. » \\
\eng{She is very \textbf{sensitive} to criticism.} « Elle est très \textbf{sensible} (émotivement) à la critique. »
\end{exemplebox}

\begin{aretenir}[Autres faux amis à haute fréquence]
\begin{itemize}
    \item \eng{to \textbf{pretend}} = faire semblant (≠ prétendre $\rightarrow$ \eng{to claim})
    \item \eng{to \textbf{demand}} = exiger (≠ demander $\rightarrow$ \eng{to ask})
    \item \eng{\textbf{eventually}} = finalement / à la fin (≠ éventuellement $\rightarrow$ \eng{possibly, maybe})
    \item \eng{\textbf{bless}} = bénir (≠ blesser $\rightarrow$ \eng{to hurt, to injure})
    \item \eng{a \textbf{library}} = bibliothèque (≠ librairie $\rightarrow$ \eng{a bookshop / bookstore})
    \item \eng{a \textbf{college}} = université / établissement post-bac (≠ collège $\rightarrow$ \eng{secondary school / high school})
    \item \eng{to \textbf{pass} an exam} = réussir un examen (≠ passer un examen $\rightarrow$ \eng{to take / sit an exam})
\end{itemize}
\end{aretenir}

\courssub{Traductions mot à mot à bannir (calques syntaxiques)}

\begin{textebox}[Observation : le piège de la traduction littérale]
\textbf{French:} J'ai 15 ans. / Faire la fête. / Prendre une décision. / Manquer le bus. / Avoir faim/soif/chaud/froid. / C'est moi qui l'ai fait. / Je m'appelle Paul.
\end{textebox}

\begin{propriete}[Les 10 calques interdits — forme correcte obligatoire]
\begin{center}
\begin{tabularx}{\linewidth}{|X|X|X|}
\hline
\rowcolor{popblueL} \textbf{Calque français (INCORRECT)} & \textbf{Forme anglaise CORRECTE} & \textbf{Traduction / Note} \\
\hline
\eng{*I have 15 years.} & \eng{I \textbf{am} 15 years old.} & \eng{to be} + âge \\
\eng{*to make a party} & \eng{to \textbf{have / throw} a party} & \eng{throw} = organiser (familier) \\
\eng{*to take a decision} & \eng{to \textbf{make} a decision} & collocation figée \\
\eng{*to miss the bus} (dans le sens « rater ») & \eng{to \textbf{miss} the bus} & \textbf{Attention :} \eng{to miss} = rater (le bus) MAIS \eng{to miss someone} = manquer à qqn (sentiment) \\
\eng{*I have hunger / thirst / hot / cold} & \eng{I \textbf{am} hungry / thirsty / hot / cold} & \eng{to be} + adjectif \\
\eng{*It is me who did it.} & \eng{\textbf{I'm the one who did it.}} / \eng{\textbf{It was me who did it.}} & structure cleft \\
\eng{*I call myself Paul.} & \eng{\textbf{My name is} Paul.} / \eng{\textbf{I'm} Paul.} & \eng{to call oneself} = s'appeler (réfléchi) rare \\
\eng{*Explain me the rule.} & \eng{\textbf{Explain the rule \textbf{to} me.}} & \eng{explain sth \textbf{to} sb} \\
\eng{*Listen me.} & \eng{\textbf{Listen \textbf{to} me.}} & \eng{listen \textbf{to}} \\
\eng{*Ask to him.} & \eng{\textbf{Ask him.}} / \eng{\textbf{Ask \textbf{for} it.}} & \eng{ask sb} / \eng{ask \textbf{for} sth} \\
\hline
\end{tabularx}
\end{center}
\end{propriete}

\begin{exemplebox}[Exemples traduits en contexte]
\eng{My sister \textbf{is twenty years old}.} « Ma sœur a vingt ans. » \\
\eng{They \textbf{threw a huge party} for her birthday.} « Ils ont organisé une énorme fête pour son anniversaire. » \\
\eng{The manager \textbf{made a difficult decision}.} « Le directeur a pris une décision difficile. » \\
\eng{\textbf{I'm hungry}. Let's go eat.} « J'ai faim. Allons manger. » \\
\eng{\textbf{Explain the problem to me}, please.} « Explique-moi le problème, s'il te plaît. »
\end{exemplebox}

\begin{attention}[Piège « manquer » : deux sens, deux verbes]
\begin{itemize}
    \item \eng{to \textbf{miss} the train / the bus / a class} = rater (ne pas attraper / ne pas être présent)
    \item \eng{to \textbf{miss} someone} = avoir le mal du pays / regretter la présence de qqn (sentiment)
    \item \eng{to \textbf{lack}} = manquer (faire défaut) : \eng{This soup \textbf{lacks} salt.} « Il manque du sel dans cette soupe. »
    \item \eng{to \textbf{be missing}} = manquer à l'appel / avoir disparu : \eng{Two pages \textbf{are missing}.} « Il manque deux pages. »
\end{itemize}
\end{attention}

\courssub{Noms indénombrables trompeurs : singulier obligatoire}

\begin{textebox}[Observation]
\textbf{French:} Une information. / Des informations. / Un conseil. / Des conseils. / Un meuble. / Des meubles. / Un devoir. / Des devoirs.

\textbf{English:} \textbf{A piece of information.} / \textbf{Information} (uncountable). / \textbf{A piece of advice.} / \textbf{Advice} (uncountable). / \textbf{A piece of furniture.} / \textbf{Furniture} (uncountable). / \textbf{Homework} (uncountable).
\end{textebox}

\begin{propriete}[Noms indénombrables (uncountable) — toujours au singulier]
Ces noms n'ont \textbf{pas de pluriel} et ne s'utilisent \textbf{jamais} avec \eng{a/an} ni avec des nombres. Pour les compter, on utilise \eng{a piece of}, \eng{an item of}, \eng{a bit of}.
\begin{center}
\begin{tabular}{|l|l|l|}
\hline
\rowcolor{popblueL} \textbf{Nom indénombrable} & \textbf{Français} & \textbf{Quantificateur} \\
\hline
\eng{information} & une information / des informations & \eng{a piece of information} \\
\eng{advice} & un conseil / des conseils & \eng{a piece of advice} \\
\eng{furniture} & un meuble / des meubles & \eng{a piece of furniture} \\
\eng{homework} & un devoir / des devoirs & \eng{a piece of homework / an assignment} \\
\eng{luggage / baggage} & un bagage / des bagages & \eng{a piece of luggage} \\
\eng{equipment} & un équipement / du matériel & \eng{a piece of equipment} \\
\eng{news} & une nouvelle / des nouvelles & \eng{a piece of news} \\
\eng{progress} & un progrès / des progrès & (pas de \eng{a piece of}) \\
\eng{knowledge} & une connaissance / des connaissances & (pas de \eng{a piece of}) \\
\eng{evidence} & une preuve / des preuves & \eng{a piece of evidence} \\
\hline
\end{tabular}
\end{center}
\end{propriete}

\begin{attention}[Erreurs fatales au Bac : \eng{*an information}, \eng{*advices}, \eng{*furnitures}, \eng{*homeworks}]
\begin{itemize}
    \item \eng{*The teacher gave us many \textbf{advices}.} $\rightarrow$ \eng{The teacher gave us \textbf{much advice} / \textbf{a lot of advice}.}
    \item \eng{*I have three \textbf{informations} for you.} $\rightarrow$ \eng{I have \textbf{three pieces of information} for you.}
    \item \eng{*She bought new \textbf{furnitures}.} $\rightarrow$ \eng{She bought new \textbf{furniture}.}
    \item \eng{*We have too many \textbf{homeworks}.} $\rightarrow$ \eng{We have \textbf{too much homework}.}
\end{itemize}
\textbf{Rappel :} \eng{much} (indénombrables) vs \eng{many} (dénombrables) ; \eng{a lot of / lots of} fonctionnent pour les deux.
\end{attention}

\begin{exemplebox}[Exemples traduits]
\eng{Can you give me \textbf{some information} about the train times?} « Peux-tu me donner des informations sur les horaires des trains ? » \\
\eng{My father gave me \textbf{a good piece of advice}.} « Mon père m'a donné un bon conseil. » \\
\eng{All the \textbf{furniture} in this room is old.} « Tous les meubles de cette pièce sont vieux. » \\
\eng{Have you finished your \textbf{homework}?} « As-tu fini tes devoirs ? »
\end{exemplebox}

\courssub{Possessifs avec les parties du corps}

\begin{textebox}[Observation]
\textbf{French:} Il a levé \textbf{la} main. / Elle s'est lavé \textbf{les} mains. / Il s'est cassé \textbf{la} jambe.

\textbf{English:} He raised \textbf{his} hand. / She washed \textbf{her} hands. / He broke \textbf{his} leg.
\end{textebox}

\begin{propriete}[Règle : possessif obligatoire]
En anglais, les parties du corps (et les vêtements, effets personnels) s'accompagnent \textbf{systématiquement} de l'adjectif possessif (\eng{my, your, his, her, its, our, their}) lorsque le possesseur est le sujet de la phrase. L'article défini \eng{the} est incorrect dans ce contexte.
\end{propriete}

\begin{center}
\begin{tabularx}{\linewidth}{|X|X|X|}
\hline
\rowcolor{popblueL} \textbf{Français (article défini)} & \textbf{Anglais INCORRECT} & \textbf{Anglais CORRECT} \\
\hline
Il a levé la main. & He raised the hand. & \eng{He raised \textbf{his} hand.} \\
Elle s'est lavé les mains. & She washed the hands. & \eng{She washed \textbf{her} hands.} \\
Il s'est cassé la jambe. & He broke the leg. & \eng{He broke \textbf{his} leg.} \\
Elle s'est brossé les dents. & She brushed the teeth. & \eng{She brushed \textbf{her} teeth.} \\
J'ai mal à la tête. & I have pain in the head. & \eng{I have a headache.} / \eng{\textbf{My} head hurts.} \\
\hline
\end{tabularx}
\end{center}

\begin{exemplebox}[Exemples traduits]
\eng{The boxer \textbf{raised his arms} in victory.} « Le boxeur a levé les bras en signe de victoire. » \\
\eng{She \textbf{injured her knee} playing football.} « Elle s'est blessée au genou en jouant au football. » \\
\eng{Don't put \textbf{your hands} in your pockets.} « Ne mets pas les mains dans les poches. »
\end{exemplebox}

\begin{attention}[Cas particulier : voix passive ou partie du corps détachée]
Si la partie du corps n'est pas « possédée » par le sujet (contexte médical, objet détaché), \eng{the} redevient possible :
\begin{itemize}
    \item \eng{The surgeon examined \textbf{the} heart.} (le cœur comme organe, pas celui du patient en particulier)
    \item \eng{They found \textbf{a} hand in the river.} (une main coupée, macabre)
\end{itemize}
Mais dans 99 \% des phrases courantes : \textbf{possessif obligatoire}.
\end{attention}

\courssub{Exercice résolu : chasse aux erreurs francophones}

\begin{exoresolu}[Correction de 10 phrases « à la française»]
\textbf{Énoncé :} Corrigez les erreurs suivantes (toutes sont des calques du français). Expliquez chaque correction.

\begin{enumerate}
    \item I don't know nothing about Cameroon politics.
    \item She is a very intelligents girl.
    \item Actually, I am living in Douala since 2020.
    \item The nurse assisted at the operation yesterday.
    \item My father gave me some good advices.
    \item We made a party for his birthday.
    \item I have 20 years and I am student.
    \item He raised the hand to answer the question.
    \item Can you explain me this grammar rule?
    \item She has many luggages.
\end{enumerate}

\tcblower

\textbf{Solution détaillée :}

\begin{enumerate}
    \item \eng{I don't know \textbf{anything} about \textbf{Cameroonian} politics.}
    \begin{itemize}
        \item Double négation interdite : \eng{don't} + \eng{nothing} $\rightarrow$ \eng{anything}.
        \item \eng{Cameroon} (nom propre) $\rightarrow$ adjectif \eng{Cameroonian}.
    \end{itemize}
    \item \eng{She is a very \textbf{intelligent} girl.}
    \begin{itemize}
        \item Adjectif invariable : pas de \eng{-s} au pluriel.
    \end{itemize}
    \item \eng{\textbf{Currently} / \textbf{At present}, I \textbf{have been living} in Douala \textbf{since} 2020.}
    \begin{itemize}
        \item \eng{Actually} = en fait (faux ami) ; pour « actuellement » : \eng{currently, at present}.
        \item Présent perfect continu \eng{have been living} avec \eng{since} (action commencée dans le passé et qui continue).
    \end{itemize}
    \item \eng{The nurse \textbf{attended} the operation yesterday.}
    \begin{itemize}
        \item \eng{to assist} = aider ; \eng{to attend} = assister à (être présent).
    \end{itemize}
    \item \eng{My father gave me \textbf{some good advice} / \textbf{a good piece of advice}.}
    \begin{itemize}
        \item \eng{advice} indénombrable : pas de pluriel, pas de \eng{an}.
    \end{itemize}
    \item \eng{We \textbf{had} / \textbf{threw} a party for his birthday.}
    \begin{itemize}
        \item \eng{to make a party} $\rightarrow$ \eng{to have / throw a party}.
    \end{itemize}
    \item \eng{I \textbf{am 20 years old} and I \textbf{am a student}.}
    \begin{itemize}
        \item Âge : \eng{to be} + \eng{years old}.
        \item \eng{student} = nom comptable singulier $\rightarrow$ \eng{a student}.
    \end{itemize}
    \item \eng{He raised \textbf{his} hand to answer the question.}
    \begin{itemize}
        \item Partie du corps $\rightarrow$ adjectif possessif.
    \end{itemize}
    \item \eng{Can you explain this grammar rule \textbf{to me}?}
    \begin{itemize}
        \item \eng{explain sth \textbf{to} sb} (pas \eng{explain sb sth}).
    \end{itemize}
    \item \eng{She has \textbf{a lot of luggage} / \textbf{much luggage}.}
    \begin{itemize}
        \item \eng{luggage} indénombrable : pas de pluriel, pas de \eng{many}.
    \end{itemize}
\end{enumerate}
\end{exoresolu}

\courssub{Entraînement : transformation et pièges}

\begin{exercice}[Mettez les phrases suivantes en anglais correct]
\begin{enumerate}
    \item Je n'ai personne à qui parler. (2 façons)
    \item Mes parents m'ont donné de bons conseils. (2 façons)
    \item Actuellement, il travaille à Yaoundé.
    \item Elle a assisté à la conférence hier.
    \item Ils ont organisé une fête surprise.
    \item J'ai 18 ans et je suis lycéen.
    \item Il s'est lavé les mains avant de manger.
    \item Explique-moi pourquoi tu es en retard.
    \item Il n'y a plus de meubles dans cette pièce.
    \item Elle est très sensible aux critiques. (deux sens possibles)
\end{enumerate}
\end{exercice}

\begin{corrige}[Exercice n°1]
\begin{enumerate}
    \item \eng{I have \textbf{no one to talk to}.} / \eng{I \textbf{don't have anyone to talk to}.}
    \item \eng{My parents gave me \textbf{some good advice}.} / \eng{My parents gave me \textbf{good pieces of advice}.}
    \item \eng{\textbf{Currently} / \textbf{At present}, he works in Yaoundé.}
    \item \eng{She \textbf{attended} the conference yesterday.}
    \item \eng{They \textbf{threw} / \textbf{had} a surprise party.}
    \item \eng{I \textbf{am 18 years old} and I \textbf{am a high school student}.}
    \item \eng{He washed \textbf{his hands} before eating.}
    \item \eng{Explain \textbf{to me} why you are late.}
    \item \eng{There is \textbf{no furniture} left in this room.} / \eng{There isn't \textbf{any furniture} in this room.}
    \item \eng{She is very \textbf{sensitive} to criticism.} (émotive) \quad OU \quad \eng{She is very \textbf{sensible} about criticism.} (raisonnable — rare)
\end{enumerate}
\end{corrige}

\begin{savaistu}[Le piège \eng{to assist at} au Bac]
Chaque session, des centaines de candidats écrivent \eng{*“The students assisted at the ceremony”} pour « Les élèves ont assisté à la cérémonie ».
\begin{itemize}
    \item \eng{to assist} = \textbf{aider}, porter assistance (\eng{The nurse assisted the surgeon.})
    \item \eng{to attend} = \textbf{assister à}, être présent à (\eng{to attend a meeting, a conference, a wedding, a ceremony})
\end{itemize}
\textbf{Mnémotechnique :} \eng{Assist} $\rightarrow$ \textbf{A}ide (les deux commencent par A). \eng{Attend} $\rightarrow$ \textbf{A}ttendre / \textbf{A}ssister à (présence).
\textbf{Au Bac :} Si le sujet est un spectateur, un participant passif $\rightarrow$ \textbf{attend}. Si le sujet aide activement $\rightarrow$ \textbf{assist}.
\end{savaistu}

\begin{methode}[Check-list « Zéro interférence française » avant de rendre sa copie]
\begin{enumerate}
    \item \textbf{Négation :} Ai-je mis \eng{not} \textbf{ou} \eng{never/nothing/nobody} (jamais les deux) ?
    \item \textbf{Adjectifs :} Sont-ils tous \textbf{invariables} et \textbf{antéposés} ?
    \item \textbf{Faux amis :} \eng{actually / currently}, \eng{assist / attend}, \eng{sensible / sensitive}, \eng{eventually / possibly}, \eng{pretend / claim}, \eng{demand / ask}, \eng{library / bookshop}, \eng{college / high school}, \eng{pass / take an exam}.
    \item \textbf{Calques :} \eng{am ... years old}, \eng{have/throw a party}, \eng{make a decision}, \eng{be hungry/thirsty/hot/cold}, \eng{explain sth to sb}, \eng{listen to}, \eng{ask sb / ask for sth}.
    \item \textbf{Indénombrables :} \eng{information, advice, furniture, homework, luggage, equipment, news, progress, knowledge, evidence} — pas de \eng{-s}, pas de \eng{a/an}, \eng{much / a lot of}.
    \item \textbf{Parties du corps :} \eng{his/her/their} + nom (jamais \eng{the}).
\end{enumerate}
\end{methode}

\coursec{Application and transformations}

Après avoir identifié les interférences du français et les faux amis dans la section précédente, nous passons maintenant à la mise en pratique. Cette section regroupe six exercices progressifs qui couvrent l'ensemble des points de grammaire révisés : accord sujet-verbe, choix entre \emph{present perfect} et \emph{past simple}, articles et prépositions, transformation de phrases, discours indirect et production écrite type Bac. Chaque exercice est suivi d'une correction commentée pour ancrer les règles.

\courssub{Exercice 1 : Correction du texte de la section 1 (recherche des 8 erreurs cachées)}

\begin{textebox}[Texte original avec 8 erreurs — Section 1 : Observation et diagnostic]
Last weekend, my family and I \textbf{go} to Buea to visit my grandmother. She \textbf{was living} there since 1990. The journey \textbf{took} three hours by bus. When we \textbf{arrived}, she \textbf{prepared} a delicious meal of \textbf{eru} and \textbf{water-fufu}. Everybody \textbf{eat} with appetite. After lunch, my uncle \textbf{tell} us stories about his youth in the 1970s. He \textbf{has worked} in Douala for twenty years before \textbf{retiring} in Buea. We \textbf{stayed} there for two days and we \textbf{enjoyed} every moment.
\end{textebox}

\begin{methode}[Comment corriger un texte : démarche en 4 étapes]
\begin{enumerate}
\item \textbf{Lisez le texte une première fois} pour en comprendre le sens global (qui ? quoi ? quand ? où ?).
\item \textbf{Surlignez chaque verbe} et demandez-vous : quel est le sujet ? Le verbe est-il au bon temps ? L'accord est-il correct ?
\item \textbf{Vérifiez les mots-outils} : articles (\emph{a/an/the/Ø}), prépositions de temps (\emph{since, for, ago, in, on, at}), connecteurs logiques.
\item \textbf{Relisez la phrase corrigée à voix haute} : l'oreille détecte souvent ce que l'œil manque.
\end{enumerate}
\end{methode}

\begin{exoresolu}[Correction guidée du texte]
\textbf{Énoncé :} Repérez et corrigez les 8 erreurs grammaticales dans le texte ci-dessus.

\tcblower
\textbf{Solution détaillée :}

\begin{center}
\begin{tabularx}{\linewidth}{|X|X|X|}
\hline
\rowcolor{popblueL} \textbf{Phrase originale (erreur soulignée)} & \textbf{Correction} & \textbf{Justification (règle)} \\
\hline
\eng{my family and I \textbf{go} to Buea} & \eng{my family and I \textbf{went} to Buea} & \textbf{Past simple} : action terminée dans le passé (\emph{Last weekend} = repère de temps passé). \\
\hline
\eng{She \textbf{was living} there since 1990} & \eng{She \textbf{has lived} there since 1990} & \textbf{Present perfect} : action commencée dans le passé (\emph{since 1990}) et qui continue au présent. \emph{Was living} (past continuous) ne s'utilise pas avec \emph{since} + date. \\
\hline
\eng{The journey \textbf{took} three hours} & \eng{The journey \textbf{took} three hours} & \textbf{Correct} : \emph{took} = past simple de \emph{take}, action passée terminée. \\
\hline
\eng{When we \textbf{arrived}, she \textbf{prepared} a delicious meal} & \eng{When we \textbf{arrived}, she \textbf{had prepared} a delicious meal} & \textbf{Past perfect} : antériorité par rapport à \emph{arrived} (elle avait déjà préparé avant notre arrivée). \\
\hline
\eng{Everybody \textbf{eat} with appetite} & \eng{Everybody \textbf{ate} with appetite} & \textbf{Accord + past simple} : \emph{Everybody} = singulier → verbe au passé (\emph{ate}). \\
\hline
\eng{my uncle \textbf{tell} us stories} & \eng{my uncle \textbf{told} us stories} & \textbf{Past simple} : \emph{told} (irrégulier), action passée terminée. \\
\hline
\eng{He \textbf{has worked} in Douala for twenty years before \textbf{retiring}} & \eng{He \textbf{had worked} in Douala for twenty years before \textbf{retiring}} & \textbf{Past perfect} : \emph{before} + action passée : l'action la plus ancienne se met au plus-que-parfait. \\
\hline
\end{tabularx}
\end{center}

\textbf{Conclusion :} dans ce texte, la plupart des erreurs viennent du mauvais choix du temps (passé simple, present perfect ou plus-que-parfait). Devant chaque verbe, cherchez d'abord le repère de temps (\eng{last weekend}, \eng{since}, \eng{before}) avant de conjuguer.
\end{exoresolu}


\newpage

\coursec{Activité d'intégration}

\coursec{Lesson 2 — Grammar: Revise the problem areas in grammar}

\courssub{Objectifs de la leçon}
À l'issue de cette leçon, l'élève doit être capable de :
\begin{itemize}
    \item \textbf{Identifier} et \textbf{corriger} les erreurs récurrentes d'accord sujet-verbe, de choix temporel (\eng{Present Perfect} vs \eng{Past Simple}), d'articles et de prépositions.
    \item \textbf{Distinguer} les faux-amis lexicaux et les interférences structurelles du français (ordre des mots, prépositions verbales).
    \item \textbf{Produire} des énoncés corrects et cohérents en mobilisant ces points de grammaire dans des tâches d'expression écrite et orale.
\end{itemize}

\courssub{Observation et diagnostic}
\begin{textebox}[Diagnostic initial — Extraits d'élèves de Terminale A]
\eng{1. My brother he works in Douala.} \\
\eng{2. I am living in Yaoundé since three years.} \\
\eng{3. She studied hardly for the exam but she failed.} \\
\eng{4. The majority of students has passed.} \\
\eng{5. We went to the Limbe Zoo the Sunday last.} \\
\eng{6. He is fond to play football.} \\
\eng{7. Actually, I am waiting the bus since 8 am.}
\end{textebox}

\begin{experience}[Activité de classe : « Grammar Clinic »]
En binôme, analysez les 7 phrases ci-dessus. Pour chacune : soulignez l'erreur, nommez la catégorie (Accord, Temps, Article/Prép, Faux-ami) et proposez une correction. Mise en commun au tableau : l'enseignant complète la « Fiche mémo » (section 6) avec les règles évoquées.
\end{experience}

\courssub{Subject-Verb Agreement (Accord sujet-verbe)}

\begin{propriete}[Règle fondamentale]
En anglais, le verbe s'accorde avec le \textbf{sujet grammatical} (pas nécessairement le nom le plus proche).
\begin{itemize}
    \item \textbf{Singulier} : 3\up{e} personne du singulier (il/elle/on) $\to$ Verbe + \eng{s} au \eng{Present Simple} (\eng{he works, she goes, it seems}). Auxiliaires : \eng{is, has, does, was}.
    \item \textbf{Pluriel} / \eng{I, You, We, They} $\to$ Verbe à la forme de base (\eng{they work, we go}). Auxiliaires : \eng{are, have, do, were}.
    \item \textbf{Sujets indéfinis} : \eng{Everyone, Everybody, Someone, Nobody, Each, Either, Neither} = \textbf{Singulier}.
    \item \textbf{Noms collectifs} : \eng{The family, The team, The majority, The government}.
    \begin{itemize}
        \item Unité $\to$ Singulier (\eng{The family \textbf{is} large.}).
        \item Membres individuels $\to$ Pluriel (\eng{The family \textbf{are} arguing among themselves.}).
    \end{itemize}
    \item \textbf{Sujets coordonnés par \eng{and}} $\to$ Pluriel. \eng{My father and my mother \textbf{are} here.}
    \item \textbf{Sujets avec \eng{or / nor / either...or / neither...nor}} $\to$ Accord avec le sujet le plus proche (\eng{proximity rule}).
\end{itemize}
\end{propriete}

\begin{exemplebox}[Exemples traduits]
\eng{Everybody \textbf{was} happy.} « Tout le monde \textbf{était} content. » (Singulier) \\
\eng{The majority of the students \textbf{were} late.} « La majorité des élèves \textbf{were} en retard. » (Pluriel car focus sur \eng{students}) \\
\eng{Neither the teacher nor the students \textbf{know} the answer.} « Ni le prof ni les élèves \textbf{ne savent}... » (Proximité : \eng{students} pluriel) \\
\eng{Either the students or the teacher \textbf{knows} the answer.} « Soit les élèves soit le prof \textbf{sait}... » (Proximité : \eng{teacher} singulier)
\end{exemplebox}

\begin{attention}[Pièges fréquents francophones]
\begin{itemize}
    \item \textbf{Sujet redoublé :} \eng{*My brother he works} $\to$ \eng{My brother \textbf{works}} (pas de pronom sujet redondant).
    \item \textbf{\eng{Everybody/Everyone} + verbe pluriel :} \eng{*Everybody were} $\to$ \eng{Everybody \textbf{was}}.
    \item \textbf{\eng{The majority of} + nom pluriel :} Souvent traité au pluriel en anglais (\eng{The majority of voters \textbf{want} change}), au singulier si focus sur le groupe (\eng{The majority \textbf{is} significant}).
    \item \textbf{\eng{There is/There are}} : \eng{There \textbf{is} a book} / \eng{There \textbf{are} books}. Erreur : \eng{*There are a lot of people} (correct) vs \eng{*There is a lot of people} (incorrect car \eng{people} pluriel).
\end{itemize}
\end{attention}

\begin{exercice}[Entraînement rapide — Accords]
Mettez le verbe à la forme correcte (Present Simple ou Past Simple).
\begin{enumerate}
    \item Everybody \eng{(to know)} the truth about the reunion.
    \item The majority of the cocoa \eng{(to be)} exported.
    \item Neither Ngwa nor his cousins \eng{(to like)} \eng{eru}.
    \item The family \eng{(to gather)} every Christmas in Buea.
    \item Each of the students \eng{(to have)} a textbook.
\end{enumerate}
\end{exercice}

\begin{corrige}[Exercice Accords]
\begin{enumerate}
    \item \eng{knows} (Everybody = singulier)
    \item \eng{is} (Cocoa = indénombrable, singulier)
    \item \eng{like} (Proximité : \eng{cousins} pluriel)
    \item \eng{gathers} (The family = unité, singulier)
    \item \eng{has} (Each of + pluriel = singulier)
\end{enumerate}
\end{corrige}

\courssub{Tenses: Present Perfect vs Past Simple}

\begin{propriete}[Choix du temps : Le critère du lien avec le présent]
\begin{center}
\begin{tabularx}{\linewidth}{|X|X|X|}
\hline
\rowcolor{popblueL} \textbf{Critère} & \textbf{Present Perfect} & \textbf{Past Simple} \\
\hline
\textbf{Lien avec le présent} & \textbf{OUI} : action passée avec conséquence/résultat maintenant. & \textbf{NON} : action terminée, datée, coupée du présent. \\
\hline
\textbf{Marqueurs typiques} & \eng{since, for, already, just, yet, ever, never, how long, recently, lately, up to now, this morning (si pas fini)} & \eng{yesterday, last week/month/year, ... ago, in 2022, on Monday, when I was young, then} \\
\hline
\textbf{Structure} & \eng{have/has + V-en (participe passé)} & \eng{V-ed} (réguliers) / 2\up{e} colonne (irréguliers) \\
\hline
\textbf{Exemple} & \eng{I \textbf{have lived} in Buea \textbf{for} 5 years.} (J'y vis encore) & \eng{I \textbf{lived} in Buea \textbf{in} 2018.} (Fini, je n'y suis plus) \\
\hline
\end{tabularx}
\end{center}
\end{propriete}

\begin{propriete}[Cas particuliers : \eng{It is the first time...} / Superlatifs]
\begin{itemize}
    \item \eng{It is/This is the first/second... time + \textbf{Present Perfect}} (\eng{This is the first time I \textbf{have seen} snow.}).
    \item Au passé (\eng{It was the first time...}) $\to$ \textbf{Past Perfect} (\eng{It was the first time I \textbf{had seen} snow.}).
    \item Superlatif + \eng{ever} $\to$ \textbf{Present Perfect} (\eng{The best eru I \textbf{have ever eaten}.}).
\end{itemize}
\end{propriete}

\begin{exemplebox}[Contraste Present Perfect / Past Simple]
\eng{I \textbf{have lost} my keys.} (Conséquence maintenant : je ne peux pas entrer) \\
\eng{I \textbf{lost} my keys \textbf{yesterday}.} (Récit passé, peut-être retrouvées depuis) \\
\eng{She \textbf{has been} a teacher \textbf{since} 2010.} (Elle l'est encore) \\
\eng{She \textbf{was} a teacher \textbf{from} 2010 \textbf{to} 2015.} (Période révolue) \\
\eng{\textbf{Have you ever} \textbf{been} to Limbe?} (Expérience de vie) \\
\eng{\textbf{Did you go} to Limbe \textbf{last weekend}?} (Action datée)
\end{exemplebox}

\begin{attention}[Erreurs types « Think in French »]
\begin{itemize}
    \item \textbf{\eng{Since} + Present Simple/Continu :} \eng{*I am here since Monday} $\to$ \eng{I \textbf{have been} here \textbf{since} Monday}.
    \item \textbf{\eng{For} + Past Simple (durée terminée) vs \eng{For} + Present Perfect (durée jusqu'à maintenant) :} \eng{I \textbf{studied} French \textbf{for} 5 years} (fini) vs \eng{I \textbf{have studied} English \textbf{for} 5 years} (continue).
    \item \textbf{Marqueur passé + Present Perfect :} \eng{*He has arrived yesterday} $\to$ \eng{He \textbf{arrived} yesterday}.
    \item \textbf{Participe passé irrégulier :} \eng{eat/ate/eaten}, \eng{give/gave/given}, \eng{go/went/gone}, \eng{write/wrote/written}. \eng{*I have ate} $\to$ \eng{I \textbf{have eaten}}.
\end{itemize}
\end{attention}

\begin{exercice}[Entraînement — Choix du temps]
Complétez avec le \eng{Present Perfect} ou le \eng{Past Simple}.
\begin{enumerate}
    \item Ngwa \eng{(write)} his article last night.
    \item He \eng{(already / send)} it to the teacher.
    \item \eng{(You / ever / eat)} \eng{ndolé}?
    \item Yes, I \eng{(eat)} it \eng{(when)} I \eng{(be)} in Douala two years ago.
    \item Mireille \eng{(study)} at UBa \eng{(for)} four years. She \eng{(graduate)} last month.
\end{enumerate}
\end{exercice}

\begin{corrige}[Exercice Temps]
\begin{enumerate}
    \item \eng{wrote} (last night = passé daté)
    \item \eng{has already sent} (already = lien présent)
    \item \eng{Have you ever eaten} (ever = expérience de vie)
    \item \eng{ate / was} (when I was in Douala two years ago = passé daté)
    \item \eng{studied / graduated} (For four years mais graduated last month = période close) \textbf{Note :} Si elle étudie encore, \eng{has studied for}.
\end{enumerate}
\end{corrige}

\courssub{Articles and Prepositions (Articles et prépositions)}

\begin{propriete}[Articles : \eng{A/An, The, \emptyset}]
\begin{itemize}
    \item \textbf{\eng{A/An}} : Singulier comptable, non identifié, première mention. \eng{A reunion, an exam}.
    \item \textbf{\eng{The}} : Défini (identifié, unique, connu, superlatif, ordinal). \eng{The reunion (dont on parle), the sun, the best student, the first time}.
    \item \textbf{\eng{\emptyset} (Zéro article)} :
    \begin{itemize}
        \item Généralités (pluriel / indénombrables) : \eng{\emptyset} \eng{Students work hard.} \eng{\emptyset} \eng{Water fufu is delicious.}
        \item Noms propres : Villes (\eng{Yaoundé}), Pays (\eng{Cameroon}), Montagnes (\eng{Mount Cameroon}), Lacs, Rues, Langues, Repas (\eng{at \emptyset breakfast}), Mois/Jours (\eng{on \emptyset Monday}), Institutions (\eng{at \emptyset school/university/hospital/church} = fonction première).
        \item \eng{By + transport} : \eng{by \emptyset car/bus/train/plane/foot}.
    \end{itemize}
\end{itemize}
\end{propriete}

\begin{propriete}[Prépositions clés : Temps et Lieu]
\begin{center}
\begin{tabularx}{\linewidth}{|l|X|X|}
\hline
\rowcolor{popblueL} \textbf{Prép.} & \textbf{Temps} & \textbf{Lieu} \\
\hline
\eng{At} & Heure précise (\eng{at 8 am}), \eng{at night}, \eng{at the weekend} (BrE), \eng{at Christmas/Easter} (période) & Point précis, adresse, événement (\eng{at the station, at the party, at 10 Main Street}) \\
\hline
\eng{On} & Jour, date, jour+moment (\eng{on Monday, on June 12th, on Monday morning}) & Surface, rue, ligne (\eng{on the table, on Molyko Street, on the coast}) \\
\hline
\eng{In} & Mois, année, saison, moment de journée (\eng{in June, in 2024, in the morning}), durée future (\eng{in 5 minutes}) & Espace clos, ville, pays, quartier (\eng{in Yaoundé, in Molyko, in the garden, in the car}) \\
\hline
\eng{Since} & Point de départ (déclenche Present Perfect) & — \\
\hline
\eng{For} & Durée (tous temps) & — \\
\hline
\eng{During} & Nom (événement/période) : \eng{during the weekend, during the meeting} & — \\
\hline
\end{tabularx}
\end{center}
\end{propriete}

\begin{attention}[Collocations verbe + préposition / Adjectif + préposition (Pièges)]
\begin{center}
\begin{tabular}{|l|l|l|}
\hline
\rowcolor{popblueL} \textbf{Incorrect (Fr $\to$ Ang)} & \textbf{Correct} & \textbf{Règle / Note} \\
\hline
\eng{*Wait someone} & \eng{Wait \textbf{for} someone} & Verbe transitif indirect \\
\eng{*Listen someone} & \eng{Listen \textbf{to} someone} & \\
\eng{*Look someone} & \eng{Look \textbf{at} someone} / \eng{Look \textbf{for} sth (chercher)} & \\
\eng{*Phone to someone} & \eng{Phone someone} (direct) / \eng{Call someone} & Pas de \eng{to} après \eng{phone/call} \\
\eng{*View on} & \eng{View \textbf{of}} & Nom \eng{view} + \eng{of} \\
\eng{*Fond to / Fond of to V} & \eng{Fond \textbf{of} V-ing} / \eng{Fond \textbf{of} n} & Adjectif + \eng{of} \\
\eng{*Interested to} & \eng{Interested \textbf{in} V-ing / n} & \\
\eng{*Good in} & \eng{Good \textbf{at} V-ing / n} & \\
\eng{*Married with} & \eng{Married \textbf{to}} & \\
\eng{*Similar to / than} & \eng{Similar \textbf{to}} & \\
\eng{*Different than} (US) / \eng{*Different to} (BrE rare) & \eng{Different \textbf{from}} (Standard BrE) & \\
\hline
\end{tabular}
\end{center}
\end{attention}

\begin{exercice}[Articles \& Prépositions]
Complétez avec \eng{a, an, the} ou \eng{\emptyset} (barré), et la préposition correcte (\eng{in, on, at, for, since, to, of}).
\begin{enumerate}
    \item We stayed \_\_\_\_ grandfather's house \_\_\_\_ Molyko \_\_\_\_ the weekend.
    \item It is \_\_\_\_ first time \_\_\_\_ three years \_\_\_\_ we meet.
    \item She is fond \_\_\_\_ photography \_\_\_\_ her childhood.
    \item The view \_\_\_\_ Mount Cameroon is stunning.
    \item I have been waiting \_\_\_\_ you \_\_\_\_ 8 o'clock.
\end{enumerate}
\end{exercice}

\begin{corrige}[Exercice Articles/Prép]
\begin{enumerate}
    \item \eng{\emptyset / in / \emptyset (BrE: at the weekend / US: on the weekend)}
    \item \eng{the / for / that} (ou \eng{Ø} si \eng{that} omis)
    \item \eng{of / since}
    \item \eng{of}
    \item \eng{for / since}
\end{enumerate}
\end{corrige}

\courssub{French Interference and False Friends (Interférences et faux-amis)}

\begin{definition}[Faux-amis lexicaux fréquents (Terminale A)]
Mots qui ressemblent à des mots français mais ont un sens différent.
\begin{center}
\begin{tabular}{|l|l|l|l|}
\hline
\rowcolor{popblueL} \textbf{Mot anglais} & \textbf{Sens anglais} & \textbf{Faux-ami français} & \textbf{Traduction correcte du sens fr} \\
\hline
\eng{Hardly} & À peine, à peine si & \eng{Durément} & \eng{Hard} \\
\eng{Actually} & En fait, en réalité & \eng{Actuellement} & \eng{Currently / At present} \\
\eng{Eventually} & Finalement, à la fin & \eng{Éventuellement} & \eng{Possibly / Potentially} \\
\eng{Pass (an exam)} & Réussir & \eng{Passer (un examen)} & \eng{Take / Sit (an exam)} \\
\eng{Assist (to)} & Aider & \eng{Assister (à)} & \eng{Attend} \\
\eng{Attend} & Assister à, être présent & \eng{Attendre} & \eng{Wait for} \\
\eng{Library} & Bibliothèque & \eng{Librairie} & \eng{Bookshop / Bookstore} \\
\eng{Parents} & Parents (père/mère) & \eng{Parents (famille élargie)} & \eng{Relatives} \\
\eng{Reunion} & Retrouvailles & \eng{Réunion (meeting)} & \eng{Meeting} \\
\eng{Graduation} & Remise de diplômes & \eng{Graduation (échelle)} & \eng{Graduation / Scale} \\
\eng{Funny} & Drôle (qui fait rire) & \eng{Amusant (agréable)} & \eng{Fun / Enjoyable} \\
\eng{Sensible} & Raisonnable, sensé & \eng{Sensible (émotif)} & \eng{Sensitive} \\
\hline
\end{tabular}
\end{center}
\end{definition}

\begin{attention}[Interférences structurelles (Syntaxique)]
\begin{itemize}
    \item \textbf{But / For + but :} \eng{*Organised for celebrate} $\to$ \eng{Organised \textbf{to celebrate}} (But = \eng{to} + infinitif). \eng{For} + V-ing = fonction (\eng{A tool for cutting}).
    \item \textbf{Adverbe de manière :} \eng{*Studied hardly} $\to$ \eng{Studied \textbf{hard}} (\eng{Hard} = adverbe irrégulier = durément ; \eng{Hardly} = à peine).
    \item \textbf{Verbe + objet direct :} \eng{*Wait the bus} $\to$ \eng{Wait \textbf{for} the bus}. \eng{*Phone to him} $\to$ \eng{Phone him}.
    \item \textbf{Since / For :} \eng{*Since three years} (durée) $\to$ \eng{\textbf{For} three years}. \eng{Since} = point de départ (\eng{since 2020, since Monday}).
    \item \textbf{Jour de la semaine :} \eng{*The Sunday} $\to$ \eng{\textbf{On} Sunday} (unique) / \eng{On Sundays} (habitude).
    \item \textbf{Être de retour :} \eng{*I am back since Monday} $\to$ \eng{I \textbf{have been} back \textbf{since} Monday}.
\end{itemize}
\end{attention}

\begin{exercice}[Faux-amis \& Structure]
Corrigez les erreurs (vocabulaire ou structure) dans ces phrases.
\begin{enumerate}
    \item Actually, I am living in Buea since 2020.
    \item She studied hardly to pass her exam.
    \item We attended the meeting and it was very funny.
    \item My parents (father and mother) are farmers, but my relatives (uncles) are teachers.
    \item He eventually arrived at 10 pm.
\end{enumerate}
\end{exercice}

\begin{corrige}[Exercice Faux-amis]
\begin{enumerate}
    \item \eng{Currently / At present, I \textbf{have been living} in Buea \textbf{since} 2020.} (\eng{Actually} $\to$ \eng{Currently} ; \eng{am living since} $\to$ \eng{have been living since})
    \item \eng{She studied \textbf{hard} to \textbf{take/sit} her exam.} (\eng{hardly} $\to$ \eng{hard} ; \eng{pass} $\to$ \eng{take/sit} pour « passer »)
    \item \eng{We \textbf{attended} the meeting and it was very \textbf{fun/enjoyable}.} (\eng{funny} $\to$ \eng{fun/enjoyable})
    \item \eng{My \textbf{parents} are farmers, but my \textbf{relatives} are teachers.} (\eng{Parents} = père/mère ; \eng{Relatives} = famille élargie).
    \item \eng{He \textbf{finally} arrived at 10 pm.} (\eng{Eventually} = finalement (après attente) ; \eng{Finalement} (résultat) = \eng{Finally/In the end}). Nuance : \eng{Eventually} implique un processus long.
\end{enumerate}
\end{corrige}

\courssub{Application and Transformations : Grammar Detectives (Intégration)}

\courssub{Situation}
\emph{Contexte :} Nous sommes à Yaoundé, au lycée bilingue d'Etoug-Ebe. Ngwa, élève en Terminale A, est le rédacteur en chef du \eng{“Family \& School Life Newsletter”}, un journal mural mensuel affiché dans le hall du lycée et partagé sur le groupe WhatsApp des parents d'élèves. Pour le numéro de juin, il a rédigé un article sur la grande réunion de famille (\eng{family reunion}) organisée à Buea le mois dernier à l'occasion de la remise de diplôme de sa sœur aînée, Mireille, à l'Université de Buea (UBa).

Ngwa a écrit son article directement en anglais, mais en pensant en français. Le texte contient de nombreuses erreurs typiques des francophones : accords oubliés, confusion \eng{Present Perfect}/\eng{Past Simple}, mauvais choix d'articles et de prépositions, et faux-amis. Le professeur d'anglais, M. Fokou, a demandé à la classe de jouer les \eng{“Grammar Detectives”} pour corriger le texte avant publication.

\begin{textebox}[Ngwa's Draft Article — “A Memorable Reunion in Buea”]
Last month, my family and I travelled to Buea for a big family reunion. It was the first time since three years that all the family gathered together. The reunion was organised for celebrate my sister Mireille's graduation at the University of Buea. She studied hardly for four years and finally obtained her Master's degree in Agricultural Engineering.

Everybody were very happy to see each other. My uncles and aunts came from Douala, Bamenda and even from Nigeria. We stayed at my grandfather's house at Molyko. The house is very large and have a beautiful view on Mount Cameroon. During the weekend, we eaten many traditional dishes like eru and water fufu, ndolé and plantains. My mother has cooked the best eru I have ever ate.

On Saturday, there was a big ceremony at the university. The Vice-Chancellor gaved a speech and gave the diplomas. Mireille was very proud. She told me: “I am waiting this moment since I was a little girl.” After the ceremony, we took many photos in front of the faculty. My father has bought a new camera for the occasion. He is very fond of photography since his childhood.

The Sunday, we went to the Limbe Botanical Garden and the Zoo. It was very funny. We saw gorillas, chimpanzees and many birds. In the evening, we shared gifts. My grandfather gived me a beautiful traditional hat. He said: “This hat is for you, it will protect you from the sun during your studies.”

Now, I am back in Yaoundé since Monday. I have already started my revisions for the Baccalauréat. I miss my family already. I hope we will organise another reunion soon.
\end{textebox}

\courssub{Consignes}
\emph{Travail en binôme (45 min) puis mise en commun (15 min).}

\begin{exercice}[Intégration : Grammar Detectives]
\begin{enumerate}
    \item \textbf{Analyse et diagnostic (Compréhension \& Grammaire).} Lisez le texte de Ngwa. Surlignez ou listez \textbf{15 erreurs grammaticales ou lexicales} distinctes. Pour chaque erreur, indiquez la \textbf{catégorie} correspondante parmi :
    \begin{itemize}
        \item[A] Accord sujet-verbe (Subject-Verb Agreement)
        \item[B] Temps verbaux : \eng{Present Perfect} vs \eng{Past Simple} (Marqueurs de temps, durée vs action terminée)
        \item[C] Articles (\eng{a/an, the, \emptyset}) et Prépositions (\eng{in, on, at, for, since, to, from, of})
        \item[D] Faux-amis / Interférence française (Vocabulaire, structure de phrase, ordre des mots)
    \end{itemize}
    
    \item \textbf{Correction ciblée (Grammaire).} Corrigez les \textbf{10 phrases suivantes} extraites du texte (ou inspirées de celui-ci). Réécrivez la phrase correcte complète et justifiez votre correction en une phrase (règle + mot-clé).
    \begin{enumerate}
        \item It was the first time since three years that all the family gathered together.
        \item The reunion was organised for celebrate my sister's graduation.
        \item She studied hardly for four years.
        \item Everybody were very happy to see each other.
        \item The house is very large and have a beautiful view on Mount Cameroon.
        \item We eaten many traditional dishes...
        \item My mother has cooked the best eru I have ever ate.
        \item She told me: “I am waiting this moment since I was a little girl.”
        \item My father has bought a new camera for the occasion.
        \item The Sunday, we went to the Limbe Botanical Garden.
    \end{enumerate}

    \item \textbf{Réécriture cohérente (Expression écrite).} Réécrivez le \textbf{deuxième paragraphe} (\eng{“Everybody were very happy... My father has bought a new camera...”}) en corrigeant toutes les erreurs et en améliorant la fluidité (connecteurs logiques, variété syntaxique). Nombre de mots visé : 60-80 mots.

    \item \textbf{Production autonome (Expression écrite).} Imaginez que vous êtes un camarade de classe de Ngwa (ou un cousin vivant à Douala). Écrivez un \textbf{court commentaire (email ou message WhatsApp formel)} de 50-70 mots pour répondre à sa newsletter. Vous devez \textbf{impérativement utiliser} :
    \begin{itemize}
        \item Au moins un verbe au \eng{Present Perfect} avec \eng{since/for}.
        \item Au moins un verbe au \eng{Past Simple} avec un marqueur de temps passé.
        \item Un sujet collectif ou indéfini avec accord correct (\eng{everyone, nobody, the family, the majority}).
        \item Une préposition de lieu/temps correcte (\eng{in, on, at}).
        \item Un mot du champ lexical de la famille/études \textbf{qui n'est pas un faux-ami} (ex. éviter \eng{actually} pour “actuellement”, \eng{eventually} pour “éventuellement”, \eng{pass} pour “réussir un examen”).
    \end{itemize}
\end{enumerate}
\end{exercice}

\courssub{Ressources à mobiliser — Fiche Mémo}
\begin{center}
\begin{tabularx}{\linewidth}{|X|X|X|}
\hline
\rowcolor{popblueL} \textbf{Notion} & \textbf{Règle clé / Structure} & \textbf{Piège francophone / Exemple correct} \\
\hline
\textbf{Accord Sujet-Verbe} & Sujet singulier $\to$ Verbe + \eng{s} (PV). Sujet pluriel / \eng{I/You} $\to$ Verbe base. \eng{Everyone/Everybody/Nobody/Each} = Singulier. \eng{The family} (unité) = Singulier ; \eng{The family} (membres) = Pluriel possible. & \eng{*Everybody were} $\to$ \eng{Everybody \textbf{was}}. \eng{*The family have} $\to$ \eng{The family \textbf{has}} (unité). \\
\hline
\textbf{Present Perfect vs Past Simple} & \eng{Present Perfect} : Lien avec le présent (\eng{since, for, already, just, ever, never, how long}). \eng{Past Simple} : Action terminée, datée (\eng{yesterday, last month, in 2022, ago}). & \eng{*I am here since Monday} $\to$ \eng{I \textbf{have been} here \textbf{since} Monday}. \eng{*He has bought it yesterday} $\to$ \eng{He \textbf{bought} it \textbf{yesterday}}. \\
\hline
\textbf{First time / Superlatif} & \eng{It is the first time + Present Perfect}. \eng{It was the first time + Past Perfect}. Superlatif + \eng{ever} + \eng{Present Perfect}. & \eng{*It was the first time I see} $\to$ \eng{It was the first time I \textbf{had seen}}. \eng{*The best I ever ate} $\to$ \eng{The best I \textbf{have ever eaten}}. \\
\hline
\textbf{Articles} & \eng{The} : défini (connu, unique, superlatif). \eng{A/An} : indéfini (singulier comptable). \emptyset : généralités (pluriel/indénombrable), noms propres (villes, pays, montagnes), repas, \eng{by + transport}, \eng{at \emptyset school/university} (fonction). & \eng{*The Sunday} $\to$ \eng{\emptyset \textbf{Sunday}} / \eng{\textbf{On} Sunday}. \eng{At \emptyset university} (étudiant) vs \eng{At \textbf{the} university} (bâtiment précis). \\
\hline
\textbf{Prépositions} & \eng{At} (point précis, événement, adresse). \eng{In} (ville, pays, rue, mois, année, matin/aprèm). \eng{On} (jour, date, jour+matin, rue). \eng{View \textbf{of}} (pas \eng{on}). \eng{Wait \textbf{for}} (pas \eng{wait} direct). \eng{Fond \textbf{of}} (pas \eng{fond to}). \eng{During} + nom (\eng{during the weekend}). & \eng{*View on} $\to$ \eng{View \textbf{of}}. \eng{*Wait someone} $\to$ \eng{Wait \textbf{for}} someone. \eng{*Fond to} $\to$ \eng{Fond \textbf{of}}. \eng{*The Sunday} $\to$ \eng{On Sunday}. \\
\hline
\textbf{Faux-amis / Interférences} & \eng{Hardly} = à peine (pas \eng{durément} $\to$ \eng{hard}). \eng{Actually} = en fait (pas \eng{actuellement} $\to$ \eng{currently}). \eng{Eventually} = finalement (pas \eng{éventuellement} $\to$ \eng{possibly}). \eng{Pass an exam} = réussir (pas \eng{passer} $\to$ \eng{take/sit}). \eng{Graduation} = remise de diplômes. \eng{Reunion} = retrouvailles. \eng{Funny} = drôle ; \eng{Fun} = amusant. & \eng{*She studied hardly} $\to$ \eng{She studied \textbf{hard}}. \eng{*He passed his exam} (ambigu) $\to$ \eng{He \textbf{passed/succeeded in}} his exam. \eng{*It was very funny} (pour « c'était amusant ») $\to$ \eng{It was very \textbf{fun/enjoyable}}. \\
\hline
\end{tabularx}
\end{center}

\begin{exoresolu}[Corrigé détaillé de l'activité « Grammar Detectives »]
\textbf{1. Analyse et diagnostic (15 erreurs identifiées)}

\begin{center}
\begin{tabularx}{\linewidth}{|c|l|l|X|}
\hline
\rowcolor{popblueL} \textbf{N°} & \textbf{Extrait erroné} & \textbf{Cat.} & \textbf{Explication / Correction} \\
\hline
1 & \eng{since three years} & B & \eng{Since} + point de départ ; \eng{For} + durée. $\to$ \eng{\textbf{for} three years}. \\
\hline
2 & \eng{for celebrate} & D & \eng{For} + V-ing (fonction) ; But = \eng{To} + infinitif. $\to$ \eng{\textbf{to celebrate}}. \\
\hline
3 & \eng{studied hardly} & D & Faux-ami. \eng{Hardly} = à peine. $\to$ \eng{studied \textbf{hard}}. \\
\hline
4 & \eng{Everybody were} & A & \eng{Everybody} = singulier. $\to$ \eng{Everybody \textbf{was}}. \\
\hline
5 & \eng{have a beautiful view on} & C & Préposition : \eng{view \textbf{of}}. \\
\hline
6 & \eng{We eaten} & B & Passé simple requis (contexte passé daté \eng{During the weekend}). $\to$ \eng{We \textbf{ate}}. \\
\hline
7 & \eng{have ever ate} & B & Participe passé irrégulier : \eng{eaten} (pas \eng{ate}). $\to$ \eng{have ever \textbf{eaten}}. \\
\hline
8 & \eng{gaved a speech} & B & Verbe irrégulier : \eng{gave} (pas \eng{gaved}). \\
\hline
9 & \eng{waiting this moment since} & C/D & \eng{Wait \textbf{for}} + objet. \eng{Since} $\to$ \eng{Present Perfect Continuous}. $\to$ \eng{\textbf{have been waiting for} this moment \textbf{since}}. \\
\hline
10 & \eng{has bought ... for the occasion} & B & Action passée datée (cérémonie de samedi). $\to$ \eng{\textbf{bought}}. \\
\hline
11 & \eng{is very fond ... since} & B/D & \eng{Fond of} + \eng{Present Perfect} pour durée depuis le passé. $\to$ \eng{\textbf{has been}} very fond of... \textbf{since}. \\
\hline
12 & \eng{The Sunday} & C & Pas d'article devant jour de la semaine (sauf habitude \eng{on Sundays}). $\to$ \eng{\textbf{On} Sunday} ou \eng{\emptyset \textbf{Sunday}}. \\
\hline
13 & \eng{very funny} & D & \eng{Funny} = drôle (qui fait rire). \eng{Fun} = amusant, agréable. $\to$ \eng{very \textbf{fun}} / \eng{enjoyable}. \\
\hline
14 & \eng{gived me} & B & Verbe irrégulier : \eng{gave}. \\
\hline
15 & \eng{I am back ... since Monday} & B & \eng{Be back} + \eng{since} $\to$ \eng{Present Perfect}. $\to$ \eng{I \textbf{have been} back \textbf{since} Monday}. \\
\hline
\end{tabularx}
\end{center}

\textbf{2. Correction ciblée des 10 phrases (avec justifications)}

\begin{enumerate}
    \item \eng{It was the first time \textbf{for} three years that all the family \textbf{had gathered} together.} (Cat. B : \eng{For} pour la durée ; \eng{Past Perfect} après \eng{It was the first time} au passé).
    \item \eng{The reunion was organised \textbf{to celebrate} my sister's graduation.} (Cat. D : \eng{To} + infinitif pour exprimer le but).
    \item \eng{She studied \textbf{hard} for four years.} (Cat. D : \eng{Hard} = adverbe de manière ; \eng{hardly} = à peine).
    \item \eng{Everybody \textbf{was} very happy to see \textbf{one another}.} (Cat. A : \eng{Everybody} singulier ; \eng{one another} pour groupe $>2$).
    \item \eng{The house is very large and \textbf{has} a beautiful view \textbf{of} Mount Cameroon.} (Cat. A : \eng{The house} = 3\up{e} pers. sing. $\to$ \eng{has} ; Cat. C : \eng{view of}).
    \item \eng{We \textbf{ate} many traditional dishes...} (Cat. B : \eng{Past Simple} action terminée \eng{During the weekend}).
    \item \eng{My mother \textbf{cooked} the best eru I \textbf{have ever eaten}.} (Cat. B : \eng{Cooked} = passé simple (action terminée week-end dernier) ; \eng{Have eaten} = \eng{Present Perfect} après superlatif \eng{best ... ever} lien présent).
    \item \eng{She told me: “I \textbf{have been waiting for} this moment \textbf{since} I was a little girl.”} (Cat. B/C : \eng{Wait for} + \eng{Present Perfect Continuous} pour durée depuis le passé avec \eng{since}).
    \item \eng{My father \textbf{bought} a new camera for the occasion.} (Cat. B : \eng{Past Simple} avec marqueur temporel implicite \eng{for the occasion} = samedi dernier).
    \item \eng{\textbf{On} Sunday, we went to the Limbe Botanical Garden.} (Cat. C : \eng{On} + jour de la semaine ; pas \eng{The}).
\end{enumerate}

\textbf{3. Réécriture du deuxième paragraphe (Modèle — 72 mots)}

\begin{textebox}[Rewritten Paragraph — Corrected \& Improved]
\textbf{Everybody \textbf{was}} delighted to see \textbf{one another}. My uncles and aunts \textbf{had come} from Douala, Bamenda, and even Nigeria. We \textbf{stayed} at my grandfather's house in Molyko, a spacious building \textbf{with} a stunning view \textbf{of} Mount Cameroon. \textbf{Over the weekend,} we \textbf{ate} traditional dishes such as \eng{eru} and water fufu, \eng{ndolé} and plantains. My mother \textbf{had cooked} the most delicious \eng{eru} I \textbf{have ever eaten}. On Saturday, the graduation ceremony \textbf{took place} at the university. The Vice-Chancellor \textbf{gave} a speech and \textbf{awarded} the diplomas. Mireille \textbf{was} beaming with pride; she \textbf{whispered}, “I \textbf{have been waiting for} this moment \textbf{since} childhood.” Afterwards, we \textbf{captured} the moment in front of the faculty. My father \textbf{had bought} a new camera especially for the event, as he \textbf{has been keen on} photography \textbf{since} his childhood.
\end{textebox}

\textbf{Commentaire des choix de langue :}
\begin{itemize}
    \item \eng{\textbf{Was} / \textbf{Had come} / \textbf{Stayed} / \textbf{Ate} / \textbf{Gave} / \textbf{Awarded} / \textbf{Took place} / \textbf{Captured} / \textbf{Bought}} : \textbf{Past Simple} (ou \eng{Past Perfect} pour l'antériorité \eng{had come, had cooked, had bought}) car toutes ces actions sont situées dans le passé révolu (\eng{Last month, On Saturday, Over the weekend}).
    \item \eng{\textbf{Have ever eaten} / \textbf{Have been waiting for} / \textbf{Has been keen on}} : \textbf{Present Perfect (Simple ou Continuous)} car lien avec le présent : expérience de vie (\eng{ever}), durée depuis le passé avec \eng{since} (\eng{waiting for, keen on}).
    \item \eng{\textbf{View of}} : Préposition correcte après \eng{view}.
    \item \eng{\textbf{One another}} : Réciproque pour groupe $>2$ (préférable à \eng{each other} ici).
    \item \eng{\textbf{Hard} / \textbf{Fond of} / \textbf{Keen on}} : Correction des faux-amis (\eng{hardly}) et collocations (\eng{fond of, keen on}).
    \item \eng{\textbf{On} Sunday / \textbf{In} Molyko} : Prépositions de temps/jour et lieu/ville correctes.
\end{itemize}

\textbf{4. Production autonome — Modèle de réponse (Email/WhatsApp)}

\begin{textebox}[Model Reply — WhatsApp Message to Ngwa]
\textbf{From:} Cousin Pierre (Douala) \\
\textbf{To:} Ngwa (Yaoundé) \\
\textbf{Date:} June 12th

Hi Ngwa,

\textbf{Everyone in the family \textbf{was}} thrilled to read your newsletter! \textbf{The majority of us \textbf{have not seen}} each other \textbf{since} the Buea reunion. I \textbf{arrived} back in Douala \textbf{last Tuesday} and \textbf{I have been revising} my Physics notes \textbf{since} Wednesday. \textbf{Currently}, I am focusing on the \eng{Baccalauréat} oral exams. \textbf{Nobody \textbf{has forgotten}} Mireille's speech; it \textbf{was} truly inspiring.

Hope to see you \textbf{in} Yaoundé soon \textbf{for} the holidays.

Best regards,
Pierre
\end{textebox}

\textbf{Vérification des contraintes imposées :}
\begin{itemize}
    \item \eng{Present Perfect + since/for} : \eng{\textbf{have not seen} ... \textbf{since}} ; \eng{\textbf{have been revising} ... \textbf{since}} ; \eng{\textbf{has forgotten}}.
    \item \eng{Past Simple + marqueur passé} : \eng{\textbf{arrived} ... \textbf{last Tuesday}} ; \eng{\textbf{was} truly inspiring}.
    \item \textbf{Sujet collectif/indéfini + accord} : \eng{\textbf{Everyone} \textbf{was}} (singulier) ; \eng{\textbf{The majority of us} \textbf{have}} (pluriel car \eng{us}) ; \eng{\textbf{Nobody} \textbf{has}} (singulier).
    \item \textbf{Préposition correcte} : \eng{\textbf{in} Yaoundé} (ville) ; \eng{\textbf{for} the holidays} (but) ; \eng{\textbf{since} Wednesday}.
    \item \textbf{Lexique sans faux-ami} : \eng{Currently} (actuellement — pas \eng{actually}) ; \eng{Focus on} (se concentrer sur) ; \eng{Inspiring} ; \eng{Oral exams} (épreuves orales — pas \eng{oral} seul comme nom). Pas d'utilisation de \eng{actually, eventually, pass} (pour réussir), \eng{hardly}.
\end{itemize}
\end{exoresolu}

\courssub{Bilan de l'activité}
\begin{aretenir}[Auto-évaluation — Suis-je prêt pour le Baccalauréat ?]
Cochez les cases correspondant à votre niveau de maîtrise après cette activité.
\begin{itemize}
    \item[$\square$] \textbf{Diagnostic :} Je sais identifier une erreur d'accord, de temps, d'article/préposition ou de faux-ami dans un texte de pair.
    \item[$\square$] \textbf{Correction :} Je peux corriger une phrase erronée et citer la règle précise (ex. \eng{“Everybody takes a singular verb”}, \eng{“Since triggers Present Perfect”}, \eng{“View of, not on”}).
    \item[$\square$] \textbf{Réécriture :} Je parviens à réécrire un paragraphe en gérant la cohérence des temps (\eng{Past Simple} pour le récit, \eng{Present Perfect} pour le lien présent) et les accords.
    \item[$\square$] \textbf{Production :} J'écris un court message personnel en intégrant naturellement les 5 contraintes grammaticales/lexicales imposées.
\end{itemize}
\textbf{Plan d'action si case non cochée :} Revoir la fiche \eng{“Ressources à mobiliser”} ci-dessus ; refaire les exercices d'application de la leçon (p. X-Y du manuel) ; demander un entretien de remédiation au professeur sur le point précis (ex. \eng{Present Perfect Continuous} avec \eng{since/for}).
\end{aretenir}

\begin{methode}[Pour réussir l'épreuve de Grammaire au Baccalauréat (Terminale A)]
\begin{enumerate}
    \item \textbf{Lisez le texte support 2 fois :} 1\up{re} fois pour le sens global (Qui ? Quoi ? Quand ?), 2\up{e} fois \textbf{au crayon} pour traquer les \textbf{indices temporels} (\eng{yesterday, since, for, last week, now, already}) et les \textbf{sujets} (singulier/pluriel, indéfinis).
    \item \textbf{Classez mentalement les erreurs :} Accord (A) ? Temps (B) ? Article/Prép (C) ? Lexique/Faux-ami (D) ? Cela active la bonne règle.
    \item \textbf{Ne traduisez pas mot à mot :} Pensez en anglais. \eng{“J'attends depuis...”} $\neq$ \eng{I wait since...} $\to$ \eng{I have been waiting for... since...}
    \item \textbf{Vérifiez les “Petits Mots” tueurs :} \eng{a/the}, \eng{in/on/at}, \eng{s} (3\up{e} pers. sing.), \eng{ed/en} (participe passé), \eng{hard/hardly}, \eng{fun/funny}.
    \item \textbf{Relisez votre production à l'envers} (dernière phrase $\to$ première) pour repérer les fautes d'inattention (orthographe, majuscules, ponctuation).
\end{enumerate}
\end{methode}

\newpage

\coursec{Méthodes et exercices résolus}

\coursec{Lesson 2 — Grammar: Revise the problem areas in grammar}

\courssub{Subject-Verb Agreement: Diagnostic and Rules}

\begin{methode}[Identifier et appliquer la concordance sujet-verbe (Concord)]
\textbf{Objectif :} Ne plus commettre d'erreurs d'accord entre le sujet et le verbe (nombre, personne, sujets complexes).
\textbf{Démarche en 4 étapes :}
\begin{enumerate}
    \item \textbf{Isoler le noyau du sujet} : Ignorer les groupes prépositionnels intercalés (\eng{of}, \eng{with}, \eng{as well as}, \eng{together with}). \eng{The box \textbf{of chocolates} is open.} (Sujet = \eng{The box}, singulier).
    \item \textbf{Analyser la nature du sujet} :
    \begin{itemize}
        \item Sujets \eng{each}, \eng{every}, \eng{none}, \eng{neither}, \eng{either}, \eng{everyone}, \eng{nobody} $\rightarrow$ \textbf{singulier}.
        \item \eng{A number of} / \eng{A variety of} + pluriel $\rightarrow$ \textbf{pluriel} ; \eng{The number of} $\rightarrow$ \textbf{singulier}.
        \item Noms collectifs (\eng{team, family, committee, government, audience}) : \textbf{singulier} si unité (décision unique), \textbf{pluriel} si membres agissent individuellement (surtout UK). \eng{The team \textbf{is} winning.} / \eng{The team \textbf{are} arguing among themselves.}
        \item Sujets joints par \eng{or / nor / either...or / neither...nor} : accord avec le sujet \textbf{le plus proche} (règle de proximité). \eng{Neither the students nor the teacher \textbf{knows}.}
        \item \eng{There is / There are} : accord avec le nom qui suit.
    \end{itemize}
    \item \textbf{Vérifier les pièges de l'inversion} : Après \eng{Here comes...}, \eng{There goes...}, \eng{Only then did he...}, le sujet suit le verbe.
    \item \textbf{Relire en remplaçant le sujet par le pronom personnel} (he/she/it $\rightarrow$ -s ; they $\rightarrow$ base verbale) pour valider l'oreille.
\end{enumerate}
\cle{Règle d'or : Le verbe s'accorde avec le SUJET, jamais avec le complément.}
\end{methode}

\begin{exoresolu}[Concordance : Choix, correction et transformation (Type Bac)]
\textbf{Énoncé :}
\begin{enumerate}
    \item Choose the correct verb form.
    \begin{enumerate}
        \item The quality of the mangoes (was / were) excellent.
        \item Neither the coach nor the players (has / have) arrived.
        \item A number of candidates (is / are) waiting outside.
        \item The committee (has / have) decided to postpone the meeting. (Context: acting as a unit).
    \end{enumerate}
    \item Correct the errors in the following sentences.
    \begin{enumerate}
        \item Each of the students have their own textbook.
        \item The news about the fuel scarcity are alarming.
        \item My friend and colleague are coming to dinner.
        \item Ten thousand francs CFA are a large sum for a student.
    \end{enumerate}
    \item Transform the sentences starting with the words given.
    \begin{enumerate}
        \item The police are investigating the case. \eng{The police...}
        \item Both solutions are acceptable. \eng{Neither solution...}
    \end{enumerate}
\end{enumerate}

\tcblower
\textbf{Solution détaillée :}
\textbf{1. Choix du verbe (Justification grammaticale) :}
\begin{enumerate}
    \item \textbf{was}. Sujet = \eng{The quality} (singulier). \eng{of the mangoes} est un complément du nom, ignoré pour l'accord.
    \item \textbf{have}. Règle de proximité avec \eng{nor} : le sujet le plus proche est \eng{the players} (pluriel).
    \item \textbf{are}. \eng{A number of} = “beaucoup de” $\rightarrow$ pluriel. (Contraste : \eng{The number of candidates is}...).
    \item \textbf{has}. Nom collectif \eng{committee} agissant comme une unité unique (décision officielle) $\rightarrow$ singulier (usage standard international).
\end{enumerate}
\textbf{2. Correction (Erreur $\rightarrow$ Correction + Explication) :}
\begin{enumerate}
    \item \eng{Each of the students \textbf{has} his/her own textbook.} (\eng{Each} = singulier ; \eng{their} évité en grammaire formelle au singulier, préférer \eng{his/her} ou reformuler).
    \item \eng{The news about the fuel scarcity \textbf{is} alarming.} (\eng{News} est un \textbf{singulier invariable} en \eng{-s} : \eng{news, mathematics, physics, economics, measles}).
    \item \eng{My friend and colleague \textbf{is} coming to dinner.} (Un seul article \eng{My} $\rightarrow$ une seule personne $\rightarrow$ singulier. Si deux personnes : \eng{My friend and \textbf{my} colleague are...}).
    \item \eng{Ten thousand francs CFA \textbf{is} a large sum...} (Sommes d'argent, distances, durées, poids considérées comme une \textbf{unité globale} $\rightarrow$ singulier).
\end{enumerate}
\textbf{3. Transformation :}
\begin{enumerate}
    \item \eng{The police \textbf{are} investigating the case.} (\eng{Police} est un \textbf{pluriel invariable} $\rightarrow$ toujours \eng{are}. Attention : \eng{A police officer is...}).
    \item \eng{Neither solution \textbf{is} acceptable.} (\eng{Neither} + singulier $\rightarrow$ verbe singulier. Sens négatif conservé).
\end{enumerate}
\textbf{Conclusion :} La maîtrise du noyau du sujet et des quantificateurs (\eng{each, neither, a number of, the number of}) est \textbf{décisive} pour le Bac.
\end{exoresolu}

\courssub{Present Perfect vs Past Simple: The Time Boundary}

\begin{methode}[Choisir entre Present Perfect et Past Simple]
\textbf{Objectif :} Distinguer action achevée dans un passé révolu (Past Simple) et passé lié au présent (Present Perfect).
\textbf{Démarche en 3 étapes (Test de décision) :}
\begin{enumerate}
    \item \textbf{Repérer les marqueurs temporels (Time Adverbials)} :
    \begin{center}
    \begin{tabularx}{\linewidth}{|X|X|}\hline
    \rowcolor{popblueL} \textbf{Past Simple (Passé révolu)} & \textbf{Present Perfect (Pont présent-passé)} \\ \hline
    \eng{yesterday, last week/year, ... ago, in 2010,} & \eng{just, already, yet, ever, never, since, for,} \\
    \eng{when I was young, once upon a time} & \eng{recently, lately, so far, up to now, this week} \\ \hline
    \end{tabularx}
    \end{center}
    \item \textbf{Poser la question : « Le moment est-il terminé ? »}
    \begin{itemize}
        \item \textbf{OUI} (hier, l'an dernier, en 2010) $\rightarrow$ \textbf{Past Simple}. \eng{I \textbf{saw} him \textbf{yesterday}.}
        \item \textbf{NON} (aujourd'hui, cette semaine, depuis 2010, toute ma vie) $\rightarrow$ \textbf{Present Perfect}. \eng{I \textbf{have seen} him \textbf{today}.}
    \end{itemize}
    \item \textbf{Test du « When ? »} : Si on peut demander \eng{When did it happen?} $\rightarrow$ Past Simple. Si la question est \eng{How long?} ou \eng{Have you ever?} $\rightarrow$ Present Perfect.
\end{enumerate}
\cle{Exception :} \eng{Present Perfect} pour annonce de nouvelle (\eng{Breaking news: The President \textbf{has resigned}.}) puis bascule en \eng{Past Simple} pour les détails (\eng{He \textbf{announced} it \textbf{last night}.}).
\end{methode}

\begin{exoresolu}[Mise au temps, justification et traduction (Type Bac)]
\textbf{Énoncé :}
\begin{enumerate}
    \item Put the verbs in brackets into the \textbf{Past Simple} or the \textbf{Present Perfect}. Justify each choice.
    \begin{enumerate}
        \item Look! Somebody (break) that window.
        \item My father (work) in Douala for ten years. Now he works in Yaoundé.
        \item Shakespeare (write) many plays.
        \item I (lose) my keys. I cannot open the door.
        \item When (you / arrive) in Buea?
    \end{enumerate}
    \item Translate into English.
    \begin{enumerate}
        \item Je connais Paul depuis 2015.
        \item Il est parti il y a deux heures.
        \item Nous n'avons pas encore fini nos devoirs.
    \end{enumerate}
\end{enumerate}

\tcblower
\textbf{Solution détaillée :}
\textbf{1. Conjugaison et Justification :}
\begin{enumerate}
    \item \eng{\textbf{has broken}}. \textbf{Justification :} Résultat visible \textbf{maintenant} (\eng{Look!}). Pas de date passée. Lien fort avec le présent.
    \item \eng{\textbf{worked}}. \textbf{Justification :} Période \textbf{terminée} (\eng{Now he works in Yaoundé} marque la coupure). \eng{For ten years} ici = durée dans le passé révolu.
    \item \eng{\textbf{wrote}}. \textbf{Justification :} Shakespeare est \textbf{mort} $\rightarrow$ période de vie révolue. Past Simple obligatoire pour personnes décédées / époques terminées.
    \item \eng{\textbf{have lost}}. \textbf{Justification :} Conséquence \textbf{actuelle} (\eng{I cannot open the door now}). \eng{Just / already / yet} implicite.
    \item \eng{\textbf{did you arrive}}. \textbf{Justification :} Question \eng{When} + verbe d'action ponctuelle $\rightarrow$ demande une date précise passée $\rightarrow$ Past Simple.
\end{enumerate}
\textbf{2. Traduction (Pièges francophones) :}
\begin{enumerate}
    \item \eng{I \textbf{have known} Paul \textbf{since} 2015.} (\textbf{Piège :} \eng{Since} + date $\rightarrow$ Present Perfect. \eng{Connaître} = stative verb $\rightarrow$ pas de \eng{been knowing}. \eng{For} 9 years aussi possible).
    \item \eng{He \textbf{left} \textbf{two hours ago}.} (\textbf{Piège :} \eng{Il y a} = \eng{ago} $\rightarrow$ \textbf{Past Simple} obligatoire. Jamais \eng{has left ago}).
    \item \eng{We \textbf{have not finished} our homework \textbf{yet}.} (\textbf{Piège :} \eng{Pas encore} = \eng{not... yet} en fin de phrase $\rightarrow$ Present Perfect. \eng{Already} en affirmatif, \eng{yet} en négatif/interrogatif).
\end{enumerate}
\textbf{Conclusion :} L'opposition \eng{ago} (Passé Simple) vs \eng{since/for} (Present Perfect) est la plus testée à l'écrit.
\end{exoresolu}

\courssub{Articles and Prepositions: Specific Difficulties}

\begin{methode}[Maîtriser l'usage des articles (a/an, the, $\emptyset$) et prépositions pièges]
\textbf{Objectif :} Éviter les omissions ou ajouts d'articles (absence d'équivalent exact en français) et choisir la bonne préposition après verbes/adjectifs.
\textbf{Démarche en 3 étapes :}
\begin{enumerate}
    \item \textbf{Articles : Appliquer la règle du « Connu vs Inconnu / Général vs Spécifique »} :
    \begin{itemize}
        \item \eng{A/An} : Première mention, classification (\eng{He is a teacher}), prix/vitesse (\eng{50 km an hour}).
        \item \eng{The} : Déjà mentionné, unique (\eng{the sun}), défini par relatif (\eng{the book \textbf{I gave you}}), superlatif (\eng{the best}), noms de fleuves/montagnes plurielles/mers (\eng{the Nile, the Alps, the Atlantic}).
        \item \textbf{Zéro article ($\emptyset$)} : Noms non dénombrables/généraux (\eng{Life is hard}, \eng{I like coffee}), noms propres (villes, pays \eng{Cameroon}, rues \eng{Messassi Street}), repas (\eng{at breakfast}), institutions (\eng{at school, in hospital, at church, in prison, at university} — \textbf{sans the} quand on y va pour leur fonction première).
    \end{itemize}
    \item \textbf{Prépositions : Mémoriser les collocations (Verbe/Adj + Prep) par cœur} :
    \begin{center}
    \begin{tabularx}{\linewidth}{|X|X|X|}\hline
    \rowcolor{popblueL} \textbf{Verbe / Adjectif} & \textbf{Préposition} & \textbf{Exemple} \\ \hline
    \eng{interested, keen} & \eng{in} & \eng{interested in politics} \\
    \eng{good, bad, excellent} & \eng{at} & \eng{good at maths} \\
    \eng{afraid, proud, ashamed} & \eng{of} & \eng{afraid of snakes} \\
    \eng{responsible, famous} & \eng{for} & \eng{responsible for the error} \\
    \eng{married, engaged} & \eng{to} & \eng{married to a doctor} \\
    \eng{depend} & \eng{on} & \eng{it depends on the weather} \\
    \eng{listen, wait} & \eng{for} & \eng{wait for the bus} \\
    \eng{look} & \eng{at / for / after} & \eng{look at} (regarder), \eng{look for} (chercher), \eng{look after} (soigner) \\ \hline
    \end{tabularx}
    \end{center}
    \item \textbf{Vérifier les prépositions de temps/lieu} : \eng{in} (mois, années, matin/après-midi, pays, grandes villes), \eng{on} (jours, dates, surfaces), \eng{at} (heures précises, adresses précises, petits lieux, événements).
\end{enumerate}
\end{methode}

\begin{exoresolu}[Remplissage articles/prépositions + Analyse d'erreurs (Type Bac)]
\textbf{Énoncé :}
\begin{enumerate}
    \item Fill in the blanks with \eng{a, an, the} or $\emptyset$ (zero article).
    \begin{enumerate}
        \item \_\_\_\_ honesty is \_\_\_\_ best policy.
        \item My brother is in \_\_\_\_ prison for \_\_\_\_ theft. (He is a prisoner).
        \item \_\_\_\_ Mount Cameroon is \_\_\_\_ highest peak in \_\_\_\_ West Africa.
        \item We had \_\_\_\_ lunch at \_\_\_\_ noon.
    \end{enumerate}
    \item Choose the correct preposition.
    \begin{enumerate}
        \item The students are interested \_\_\_\_ (in / on / at) learning English.
        \item Cameroon is rich \_\_\_\_ (in / with / of) natural resources.
        \item He insisted \_\_\_\_ (on / to / for) paying the bill.
        \item She is good \_\_\_\_ (at / in / for) singing.
    \end{enumerate}
    \item Correct the following sentences.
    \begin{enumerate}
        \item I go to the school by foot every morning.
        \item He plays the football very well.
        \item We discussed about the problem.
        \item I prefer tea than coffee.
    \end{enumerate}
\end{enumerate}

\tcblower
\textbf{Solution détaillée :}
\textbf{1. Articles :}
\begin{enumerate}
    \item $\emptyset$ honesty is \textbf{the} best policy. (\eng{Honesty} = concept abstrait/général $\rightarrow$ $\emptyset$. \eng{The best} = superlatif $\rightarrow$ \eng{the}).
    \item My brother is in $\emptyset$ prison for $\emptyset$ theft. (Institution $\rightarrow$ $\emptyset$ prison / $\emptyset$ hospital / $\emptyset$ school. \eng{Theft} non dénombrable ici $\rightarrow$ $\emptyset$).
    \item $\emptyset$ Mount Cameroon is \textbf{the} highest peak in $\emptyset$ West Africa. (Monts isolés $\rightarrow$ $\emptyset$ Mount. Superlatif $\rightarrow$ \eng{the}. Régions cardinales + nom propre $\rightarrow$ $\emptyset$ West Africa).
    \item We had $\emptyset$ lunch at $\emptyset$ noon. (Repas et moments de la journée $\rightarrow$ $\emptyset$).
\end{enumerate}
\textbf{2. Prépositions :}
\begin{enumerate}
    \item \textbf{in} (\eng{interested in}).
    \item \textbf{in} (\eng{rich in}).
    \item \textbf{on} (\eng{insist on} + V-ing).
    \item \textbf{at} (\eng{good at}).
\end{enumerate}
\textbf{3. Correction (Erreurs types francophones) :}
\begin{enumerate}
    \item \eng{I go to $\emptyset$ school on foot every morning.} (\textbf{The} school = le bâtiment ; $\emptyset$ school = l'institution. \eng{By foot} $\rightarrow$ \eng{on foot}).
    \item \eng{He plays $\emptyset$ football very well.} (Sports/Jeux $\rightarrow$ $\emptyset$ article. \eng{Play the piano} pour instruments).
    \item \eng{We discussed the problem.} (\eng{Discuss} est transitif direct $\rightarrow$ pas de \eng{about}. \eng{Talk about} / \eng{Speak about}).
    \item \eng{I prefer tea \textbf{to} coffee.} (\eng{Prefer A to B} pour noms. \eng{Would rather} + V + \eng{than} pour verbes).
\end{enumerate}
\textbf{Conclusion :} Les collocations \eng{prefer to}, \eng{discuss} (sans prép), \eng{on foot}, $\emptyset$ \eng{school/prison/hospital} sont des « points chauds » du Bac.
\end{exoresolu}

\courssub{French Interference: False Friends and Calques}

\begin{methode}[Éviter les faux-amis et les calques syntaxiques]
\textbf{Objectif :} Identifier les mots trompeurs (faux-amis) et ne pas traduire mot-à-mot les structures françaises.
\textbf{Démarche en 3 étapes :}
\begin{enumerate}
    \item \textbf{Lister et mémoriser le « Top 20 » des faux-amis fréquents au Bac} :
    \begin{center}
    \begin{tabularx}{\linewidth}{|X|X|X|}\hline
    \rowcolor{popblueL} \textbf{Mot anglais} & \textbf{Faux sens (FR)} & \textbf{Vrai sens (EN $\rightarrow$ FR)} \\ \hline
    \eng{Actually} & Actuellement & \textbf{En fait, réellement} (\eng{Currently} = actuellement) \\
    \eng{Eventually} & Éventuellement & \textbf{Finalement, au bout du compte} (\eng{Possibly} = éventuellement) \\
    \eng{Sensible} & Sensible (émotions) & \textbf{Raisonnable, sensé} (\eng{Sensitive} = sensible) \\
    \eng{Library} & Librairie & \textbf{Bibliothèque} (\eng{Bookshop} = librairie) \\
    \eng{Receipt} & Recette (cuisine) & \textbf{Reçu (paiement)} (\eng{Recipe} = recette) \\
    \eng{Preservatives} & Préservatifs & \textbf{Conservateurs (alimentaires)} (\eng{Condoms} = préservatifs) \\
    \eng{Assist} & Assister (regarder) & \textbf{Aider} (\eng{Attend / Watch} = assister à) \\
    \eng{Realise} & Réaliser (construire) & \textbf{Prendre conscience, se rendre compte} (\eng{Achieve / Build} = réaliser) \\
    \eng{Parents} & Parents (famille élargie) & \textbf{Père et mère uniquement} (\eng{Relatives} = parents/parenté) \\
    \eng{Argument} & Dispute, querelle & \textbf{Raison, argumentation (dans un débat)} (\eng{Quarrel / Row} = dispute) \\ \hline
    \end{tabularx}
    \end{center}
    \item \textbf{Éviter les 5 calques syntaxiques majeurs} :
    \begin{enumerate}
        \item \eng{« J'ai 18 ans »} $\neq$ \eng{I have 18 years} $\rightarrow$ \textbf{I \eng{am} 18 (years old).}
        \item \eng{« Il fait chaud »} $\neq$ \eng{It makes hot} $\rightarrow$ \textbf{It \eng{is} hot.}
        \item \eng{« Je lui ai téléphoné »} $\neq$ \eng{I phoned to him} $\rightarrow$ \textbf{I phoned him.} (Verbe transitif direct).
        \item \eng{« Elle est mariée avec lui »} $\neq$ \eng{She is married with him} $\rightarrow$ \textbf{She is married \eng{to} him.}
        \item \eng{« Cela dépend de... »} $\neq$ \eng{It depends of...} $\rightarrow$ \textbf{It depends \eng{on}...}
    \end{enumerate}
    \item \textbf{Stratégie de révision} : Pour tout mot ressemblant au français, \textbf{vérifier au dictionnaire} avant d'écrire. Reformuler la phrase en anglais simple (SVO) plutôt que de calquer le français.
\end{enumerate}
\end{methode}

\begin{exoresolu}[Correction de faux-amis et reformulation (Type Bac)]
\textbf{Énoncé :}
\begin{enumerate}
    \item Identify and correct the false friend in each sentence.
    \begin{enumerate}
        \item The meeting is \textbf{actually} in progress. (Context: Right now).
        \item He is a very \textbf{sensible} boy; he cries easily.
        \item I went to the \textbf{library} to buy a novel.
        \item We \textbf{assisted} at the conference yesterday.
        \item The \textbf{opportunity} of rain is high today.
    \end{enumerate}
    \item Rewrite the sentences correctly (avoid French calques).
    \begin{enumerate}
        \item She has 20 years.
        \item It makes cold today.
        \item I asked to him the question.
        \item They entered in the room.
        \item He explained me the lesson.
    \end{enumerate}
    \item Translate into natural English: « \textit{Finalement, il a réalisé son rêve de devenir médecin.} »
\end{enumerate}

\tcblower
\textbf{Solution détaillée :}
\textbf{1. Faux-amis (Correction + Vrai mot pour le sens voulu) :}
\begin{enumerate}
    \item \eng{actually} $\rightarrow$ \textbf{currently} / \eng{at the moment}. (\eng{Actually} = « en fait »).
    \item \eng{sensible} $\rightarrow$ \textbf{sensitive}. (\eng{Sensible} = « raisonnable »).
    \item \eng{library} $\rightarrow$ \textbf{bookshop} / \eng{bookstore}. (\eng{Library} = « bibliothèque »).
    \item \eng{assisted at} $\rightarrow$ \textbf{attended}. (\eng{Assist} = « aider »).
    \item \eng{opportunity} $\rightarrow$ \textbf{likelihood} / \eng{probability} / \eng{chance}. (\eng{Opportunity} = « occasion favorable »).
\end{enumerate}
\textbf{2. Calques syntaxiques (Correction + Règle) :}
\begin{enumerate}
    \item \eng{She \textbf{is} 20 (years old).} (\eng{Age} = \eng{be} + number).
    \item \eng{It \textbf{is} cold today.} (\eng{Weather} = \eng{It is} + adj).
    \item \eng{I asked him the question.} / \eng{I asked the question \textbf{to} him.} (\eng{Ask} qqn qqch / ask qqch \textbf{to} qqn. Pas \eng{ask to} qqn).
    \item \eng{They entered the room.} (\eng{Enter} = transitif direct, pas de \eng{in}).
    \item \eng{He explained the lesson \textbf{to} me.} (\eng{Expliquer qqch à qqn} = \eng{explain sth \textbf{to} sb}. Jamais \eng{explain me}).
\end{enumerate}
\textbf{3. Traduction :}
\eng{\textbf{Eventually}, he \textbf{achieved} / \textbf{fulfilled} his dream of becoming a doctor.}
\begin{itemize}
    \item \eng{Finalement} = \eng{Eventually} (pas \eng{Finally} qui implique une liste).
    \item \eng{Réaliser (un rêve)} = \eng{achieve} / \eng{fulfill} / \eng{accomplish} (pas \eng{realise} = prendre conscience).
\end{itemize}
\textbf{Conclusion :} La confusion \eng{Realise/Achieve} et \eng{Eventually/Finally} coûte très cher en expression écrite.
\end{exoresolu}

\courssub{Reported Speech and Question Tags: Structural Revisions}

\begin{methode}[Transformer le discours direct en indirect et former les Question Tags]
\textbf{Objectif :} Maîtriser le \eng{Backshift} (décalage des temps), les changements de pronoms/adverbes, et la formation des \eng{Question Tags}.
\textbf{Démarche en 4 étapes (Discours Rapporté) :}
\begin{enumerate}
    \item \textbf{Identifier le verbe introducteur} : S'il est au \textbf{Past} (\eng{said, told, asked}), appliquer le \textbf{Backshift}. S'il est au \textbf{Present} (\eng{says, thinks}), \textbf{pas de backshift}.
    \item \textbf{Appliquer le Backshift (Tableau de correspondance)} :
    \begin{center}
    \begin{tabularx}{\linewidth}{|X|X|}\hline
    \rowcolor{popblueL} \textbf{Direct Speech} & \textbf{Reported Speech (Backshift)} \\ \hline
    Present Simple (\eng{I like}) & Past Simple (\eng{he liked}) \\
    Present Continuous (\eng{I am reading}) & Past Continuous (\eng{he was reading}) \\
    Present Perfect (\eng{I have done}) & Past Perfect (\eng{he had done}) \\
    Past Simple (\eng{I went}) & Past Perfect (\eng{he had gone}) \\
    \eng{will} (\eng{I will go}) & \eng{would} (\eng{he would go}) \\
    \eng{can} (\eng{I can swim}) & \eng{could} (\eng{he could swim}) \\
    \eng{must} (\eng{I must go}) & \eng{had to} (\eng{he had to go}) \\
    \eng{may} & \eng{might} \\ \hline
    \end{tabularx}
    \end{center}
    \item \textbf{Adapter pronoms, possessifs, adverbes temps/lieu} :
    \eng{I $\rightarrow$ he/she}, \eng{my $\rightarrow$ his/her}, \eng{now $\rightarrow$ then}, \eng{today $\rightarrow$ that day}, \eng{yesterday $\rightarrow$ the day before / the previous day}, \eng{tomorrow $\rightarrow$ the next day / the following day}, \eng{here $\rightarrow$ there}, \eng{this $\rightarrow$ that}, \eng{these $\rightarrow$ those}, \eng{ago $\rightarrow$ before}.
    \item \textbf{Structure des questions indirectes} : \eng{He asked \textbf{where I lived}.} (Ordre affirmatif \textbf{Sujet + Verbe}, pas d'auxiliaire \eng{do/does/did}, pas de point d'interrogation). \eng{If / Whether} pour questions fermées.
\end{enumerate}
\textbf{Question Tags (Règle de l'alternance) :}
\begin{enumerate}
    \item Phrase affirmative $\rightarrow$ Tag \textbf{négatif}. \eng{You are tired, \textbf{aren't you}?}
    \item Phrase négative $\rightarrow$ Tag \textbf{affirmatif}. \eng{You aren't tired, \textbf{are you}?}
    \item Verbe du tag = Auxiliaire/Modal de la phrase (ou \eng{do/does/did}). Sujet du tag = \textbf{Pronom personnel}.
    \item \textbf{Cas particuliers} : \eng{I am $\rightarrow$ aren't I?} / \eng{Let's $\rightarrow$ shall we?} / \eng{Don't $\rightarrow$ will you?} / \eng{Imperative $\rightarrow$ will/would/can/could you?} / \eng{Nobody/Everyone $\rightarrow$ they?} / \eng{Nothing $\rightarrow$ it?}.
\end{enumerate}
\end{methode}

\begin{exoresolu}[Transformation Reported Speech + Question Tags (Type Bac)]
\textbf{Énoncé :}
\begin{enumerate}
    \item Turn the following sentences into \textbf{Reported Speech}.
    \begin{enumerate}
        \item "I will visit the museum tomorrow," she said.
        \item "Did you finish your homework?" the teacher asked me.
        \item "Don't make noise," the librarian told the students.
    \end{enumerate}
    \item Add the correct \textbf{Question Tags}.
    \begin{enumerate}
        \item Let's go for a walk, \_\_\_\_\_\_\_\_ ?
        \item Nobody called me yesterday, \_\_\_\_\_\_\_\_ ?
        \item I am right, \_\_\_\_\_\_\_\_ ?
        \item Don't forget to lock the door, \_\_\_\_\_\_\_\_ ?
        \item This is your book, \_\_\_\_\_\_\_\_ ?
    \end{enumerate}
    \item Combine: Report the question + add a tag to the reporting clause.
    \eng{John said: "Are you coming?"} $\rightarrow$ John asked \_\_\_\_\_\_\_\_, \_\_\_\_\_\_\_\_ ?
\end{enumerate}

\tcblower
\textbf{Solution détaillée :}
\textbf{1. Reported Speech (Backshift + Changements) :}
\begin{enumerate}
    \item \eng{She said (that) she \textbf{would visit} the museum \textbf{the next day} / \textbf{the following day}.} (\eng{will $\rightarrow$ would}, \eng{tomorrow $\rightarrow$ the next day}).
    \item \eng{The teacher asked me \textbf{if / whether I had finished} my homework \textbf{the day before} / \textbf{the previous day}.} (Question fermée $\rightarrow$ \eng{if/whether}. Ordre affirmatif \eng{I had finished}. \eng{yesterday $\rightarrow$ the day before}. Past Simple $\rightarrow$ Past Perfect).
    \item \eng{The librarian told the students \textbf{not to make} noise.} (Impératif négatif $\rightarrow$ \eng{not to} + BV. \eng{Told} + objet + \eng{not to}).
\end{enumerate}
\textbf{2. Question Tags :}
\begin{enumerate}
    \item \eng{Let's go for a walk, \textbf{shall we}?} (Cas particulier \eng{Let's}).
    \item \eng{Nobody called me yesterday, \textbf{did they}?} (Sujet indéfini négatif $\rightarrow$ tag positif, pronom \eng{they}).
    \item \eng{I am right, \textbf{aren't I}?} (Cas particulier \eng{I am} $\rightarrow$ \eng{aren't I} ? formel : \eng{am I not?}).
    \item \eng{Don't forget to lock the door, \textbf{will you}?} (Impératif négatif $\rightarrow$ tag positif \eng{will you}).
    \item \eng{This is your book, \textbf{isn't it}?} (Sujet \eng{This} $\rightarrow$ pronom \eng{it}).
\end{enumerate}
\textbf{3. Combinaison :}
\eng{John asked \textbf{if/whether I was coming}, \textbf{didn't he}?}
\begin{itemize}
    \item Discours rapporté : \eng{Are you coming?} $\rightarrow$ \eng{if/whether I was coming} (Present Cont $\rightarrow$ Past Cont, \eng{you $\rightarrow$ I}).
    \item Question tag sur la principale : \eng{John asked...} (Affirmatif, Past Simple) $\rightarrow$ \eng{didn't he}? (Auxiliaire \eng{did} au passé).
\end{itemize}
\textbf{Conclusion :} Le \eng{Backshift} et l'ordre \eng{Sujet+Verbe} dans les questions indirectes sont systématiquement testés. Les tags \eng{aren't I}, \eng{shall we}, \eng{will you} (impératif) sont des « bonus » faciles si mémorisés.
\end{exoresolu}

\courssub{Conditional Sentences and Wish/If only: Hypothetical Meaning}

\begin{methode}[Maîtriser les Conditionnels (Type 0, 1, 2, 3, Mixte) et Wish/If only]
\textbf{Objectif :} Construire correctement les phrases hypothétiques selon le degré de réalité/possibilité.
\textbf{Démarche en 3 étapes :}
\begin{enumerate}
    \item \textbf{Identifier le TYPE de conditionnel selon le sens (Réalité vs Irréalité)} :
    \begin{center}
    \begin{tabularx}{\linewidth}{|c|X|X|}\hline
    \rowcolor{popblueL} \textbf{Type} & \textbf{Sens / Temps} & \textbf{Structure (If-clause / Main clause)} \\ \hline
    \textbf{0} & Vérité générale, scientifique & \eng{If + Present Simple, ... Present Simple} \\
    \textbf{1} & Possible / Probable (Futur) & \eng{If + Present Simple, ... will/can/may + BV} \\
    \textbf{2} & Irréel Présent / Peu probable & \eng{If + Past Simple, ... would/could/might + BV} \\
    \textbf{3} & Irréel Passé (Regret) & \eng{If + Past Perfect, ... would/could/might + have + V3} \\
    \textbf{Mixte} & Passé $\rightarrow$ Présent & \eng{If + Past Perfect, ... would + BV} \\
    & Présent $\rightarrow$ Passé & \eng{If + Past Simple, ... would + have + V3} \\ \hline
    \end{tabularx}
    \end{center}
    \item \textbf{Règles d'or du \eng{If}-clause} : \textbf{JAMAIS} de \eng{will / would / would have} dans la proposition \eng{If}. \eng{If I \textbf{would know}...} $\rightarrow$ \textbf{FAUX}. \eng{If I \textbf{knew}...} (Type 2) ou \eng{If I \textbf{had known}...} (Type 3).
    \item \textbf{Wish / If only (Regrets)} : Suivent la logique des Types 2 et 3.
    \begin{itemize}
        \item Regret Présent/Futur (Type 2) : \eng{I wish / If only + Past Simple / could + BV}. \eng{I wish I \textbf{knew} the answer.}
        \item Regret Passé (Type 3) : \eng{I wish / If only + Past Perfect}. \eng{I wish I \textbf{had come} earlier.}
        \item \eng{Wish + would} : Seulement pour exprimer un souhait de \textbf{changement de comportement d'autrui} (agacement). \eng{I wish he \textbf{would stop} smoking.}
    \end{itemize}
\end{enumerate}
\end{methode}

\begin{exoresolu}[Conditionnels mixtes, Wish et Transformation (Type Bac)]
\textbf{Énoncé :}
\begin{enumerate}
    \item Put the verbs in brackets into the correct tense.
    \begin{enumerate}
        \item If you (heat) ice, it (melt).
        \item If I (be) you, I (accept) the offer.
        \item If she (study) harder last year, she (pass) her exams.
        \item If they (not / miss) the bus, they (be) here by now.
        \item I wish I (know) his phone number. (Present regret)
        \item I wish I (not / spend) all my money yesterday. (Past regret)
    \end{enumerate}
    \item Rewrite using \eng{Wish} or \eng{If only}.
    \begin{enumerate}
        \item I don't have a laptop. (Present)
        \item He didn't wake up early. (Past)
        \item She won't listen to me. (Annoyance/Behaviour)
    \end{enumerate}
    \item Correct the error: \eng{If I would have known, I would have told you.}
\end{enumerate}

\tcblower
\textbf{Solution détaillée :}
\textbf{1. Conjugaison (Identification du type) :}
\begin{enumerate}
    \item \eng{If you \textbf{heat} ice, it \textbf{melts}.} (Type 0 : Vérité scientifique. Present / Present).
    \item \eng{If I \textbf{were} you, I \textbf{would accept} the offer.} (Type 2 : Conseil/irréel présent. \eng{Were} pour toutes les personnes. \eng{Would} + BV).
    \item \eng{If she \textbf{had studied} harder last year, she \textbf{would have passed} her exams.} (Type 3 : Regret passé. Past Perfect / \eng{Would have} + V3).
    \item \eng{If they \textbf{hadn't missed} the bus, they \textbf{would be} here by now.} (Mixte Type 3 $\rightarrow$ 2 : Cause passée, conséquence présente. Past Perfect / \eng{Would} + BV).
    \item \eng{I wish I \textbf{knew} his phone number.} (Regret présent $\rightarrow$ Past Simple / \eng{could}).
    \item \eng{I wish I \textbf{hadn't spent} all my money yesterday.} (Regret passé $\rightarrow$ Past Perfect).
\end{enumerate}
\textbf{2. Reformulation \eng{Wish / If only} :}
\begin{enumerate}
    \item \eng{I wish / If only I \textbf{had} a laptop.} (Ou \eng{could have}).
    \item \eng{He wishes / If only he \textbf{had woken} up early.}
    \item \eng{I wish / If only she \textbf{would listen} to me.} (\eng{Would} pour comportement d'autrui).
\end{enumerate}
\textbf{3. Correction :}
\eng{If I \textbf{had known}, I would have told you.}
\begin{itemize}
    \item Erreur : \eng{would have} dans le \eng{If}-clause (interdit).
    \item Correction : Type 3 $\rightarrow$ \eng{If + Past Perfect} (\eng{had known}).
\end{itemize}
\textbf{Conclusion :} La distinction Type 2 / Type 3 et le \eng{Mixed Conditional} (cause passée $\rightarrow$ conséquence présente) sont des classiques du Bac. \eng{Wish + would} = agacement.
\end{exoresolu}

\begin{attention}[Les 5 erreurs qui coûtent des points au Bac]
\begin{enumerate}
    \item \textbf{Concordance Sujet-Verbe (Concord) :} Oublier le noyau du sujet (\eng{The *list of items are* long $\rightarrow$ \textbf{is}}). Confondre \eng{Each/Every/None} (singulier) et \eng{A number of} (pluriel). Traiter \eng{Police/News/Mathematics} comme pluriels.
    \item \textbf{Present Perfect vs Past Simple :} Utiliser \eng{Present Perfect} avec \eng{ago, yesterday, last year} (\eng{*I have seen him yesterday*}). Oublier \eng{Present Perfect} avec \eng{since, for, just, already, yet} quand le lien au présent existe.
    \item \textbf{Faux-amis et Calques (Vocabulaire/Syntaxe) :} \eng{Actually} pour « actuellement », \eng{Sensible} pour « sensible », \eng{Library} pour « librairie », \eng{Assist} pour « assister à ». Calques : \eng{*I have 20 years*}, \eng{*It makes cold*}, \eng{*Discuss about*}, \eng{*Explain me*}, \eng{*Prefer ... than*}.
    \item \textbf{Reported Speech (Backshift) :} Ne pas décaler les temps (\eng{*He said he will come*}). Garder l'ordre interrogatif dans la question indirecte (\eng{*He asked where did I go*}). Oublier les changements \eng{now/then, here/there, this/that}.
    \item \textbf{Conditionnels et Wish :} Mettre \eng{will/would} dans le \eng{If}-clause (\eng{*If I would know...*}). Confondre Type 2 (\eng{If + Past, would + BV}) et Type 3 (\eng{If + Past Perfect, would have + V3}). Utiliser \eng{Wish + Present} pour un regret (\eng{*I wish I have...*}) au lieu de \eng{Past Simple / Past Perfect}.
\end{enumerate}
\cle{Astuce finale :} Relisez votre copie \textbf{uniquement pour ces 5 points} (verbes, marqueurs temporels, faux-amis, discours rapporté, \eng{if}-clauses). 80\% des fautes de grammaire y sont concentrées.
\end{attention}

\newpage

\coursec{Exercices}

\courssub{Je vérifie mes connaissances}

\begin{exercice}[QCM : zones de difficulté grammaticale]
\emph{Choisissez la bonne réponse parmi les trois propositions. Une seule réponse est correcte.}
\begin{enumerate}
    \item My uncle \eng{has been} in Douala \eng{since} / \eng{for} / \eng{from} 2015.
    \item Look! The children \eng{play} / \eng{are playing} / \eng{have played} football in the courtyard right now.
    \item \eng{—} / \eng{A} / \eng{The} honesty is a virtue highly valued in our community.
    \item Neither the president nor his assistants \eng{was} / \eng{were} / \eng{are} present at the meeting yesterday.
    \item She asked me where \eng{I live} / \eng{I lived} / \eng{did I live}.
    \item If he \eng{will come} / \eng{comes} / \eng{came} early, we will go to the market together.
    \item My father \eng{has gone} / \eng{has been} / \eng{went} to Buea; he will be back tomorrow.
    \item There isn't \eng{some} / \eng{any} / \eng{no} water left in the jar.
    \item The news \eng{are} / \eng{is} / \eng{were} very encouraging for the students.
    \item I am used to \eng{eat} / \eng{eating} / \eng{ate} fufu and eru on Sundays.
\end{enumerate}
\end{exercice}

\begin{exercice}[Vrai ou Faux grammatical : justifiez]
\emph{Lisez les phrases ci-dessous. Écrivez « Correct » ou « Incorrect ». Si la phrase est incorrecte, réécrivez-la correctement et citez la règle violée (accord, temps, article, préposition, ordre des mots).}
\begin{enumerate}
    \item My brother and his friend is coming to Yaoundé next week.
    \item She has been living in Bamenda since three years.
    \item The police is investigating the matter.
    \item He suggested me to go to the hospital.
    \item Despite the rain was heavy, we attended the ceremony.
    \item My mother she cooks delicious ndolé.
    \item I prefer tea than coffee.
    \item It is high time we leave this place.
\end{enumerate}
\end{exercice}

\begin{exercice}[Vocabulaire : Family and Social Life]
\emph{Complétez les phrases avec le mot approprié choisi dans la liste. Attention : deux mots sont en trop.}
\begin{center}
\textbf{extended family — bride price — kinship — nuclear family — lineage — widow — inheritance — foster — sibling — ceremony}
\end{center}
\begin{enumerate}
    \item In many Cameroonian cultures, the payment of the \_\_\_\_\_\_\_\_\_\_\_\_ seals the marriage alliance between two families.
    \item My \_\_\_\_\_\_\_\_\_\_\_\_ includes my parents, my two brothers and myself; my grandparents live in the village.
    \item The \_\_\_\_\_\_\_\_\_\_\_\_ system traces descent through the father's side in the Beti tradition.
    \item After the death of her husband, the \_\_\_\_\_\_\_\_\_\_\_\_ wore mourning clothes for a year.
    \item The \_\_\_\_\_\_\_\_\_\_\_\_ to the throne follows strict traditional rules in the Grassfields.
    \item During the \_\_\_\_\_\_\_\_\_\_\_\_, the elders pour libations to honour the ancestors.
\end{enumerate}
\end{exercice}

\begin{exercice}[Conjugaison et accord : mise en forme]
\emph{Mettez les verbes entre parenthèses au temps qui convient (Present Simple, Present Continuous, Present Perfect, Past Simple, Past Continuous). Respectez l'accord sujet-verbe.}
\begin{enumerate}
    \item Every morning, the bendskin driver \_\_\_\_\_\_\_\_\_\_\_\_ (stop) in front of the school gate.
    \item Listen! Somebody \_\_\_\_\_\_\_\_\_\_\_\_ (play) the balafon near the palace.
    \item My elder sister \_\_\_\_\_\_\_\_\_\_\_\_ (just / finish) her degree at the University of Buea.
    \item While we \_\_\_\_\_\_\_\_\_\_\_\_ (have) dinner last night, the power \_\_\_\_\_\_\_\_\_\_\_\_ (go) out.
    \item How long \_\_\_\_\_\_\_\_\_\_\_\_ (you / know) Mr. Nkeng? — Since 2010.
    \item The number of students in this class \_\_\_\_\_\_\_\_\_\_\_\_ (increase) every year.
    \item By the time the chief arrived, the dancers \_\_\_\_\_\_\_\_\_\_\_\_ (already / perform) for an hour.
\end{enumerate}
\end{exercice}

\courssub{Je m'entraîne}

\begin{exercice}[Traduction ciblée : interférences français-anglais]
\emph{Traduisez les phrases suivantes en anglais. Attention aux faux-amis, aux prépositions, à l'ordre des mots et aux temps.}
\begin{enumerate}
    \item Mon père est professeur depuis vingt ans.
    \item Il est parti à Douala hier matin.
    \item Elle m'a dit qu'elle viendrait demain.
    \item Je n'ai personne à qui parler.
    \item C'est la première fois que je mange du ndolé.
    \item Il a suggéré d'attendre le bus.
    \item Malgré qu'il était fatigué, il a travaillé.
    \item Elle s'est mariée avec un médecin.
    \item Je lui ai demandé ce qu'il faisait.
    \item Il y a longtemps que je ne t'ai vu.
\end{enumerate}
\end{exercice}

\begin{exercice}[Transformations structurales]
\emph{Effectuez les transformations demandées.}
\begin{enumerate}
    \item \textbf{Discours indirect :} Tournez les phrases suivantes au discours rapporté (introducteur au passé).
    \begin{enumerate}
        \item “I will visit my grandparents in the village next week,” she said.
        \item “Have you finished your homework?” the teacher asked me.
        \item “Don't make noise during the ceremony,” the elder told the children.
        \item “Where does your brother work?” he wanted to know.
        \item “Let's organize a fundraising for the school,” they proposed.
    \end{enumerate}
    \item \textbf{Forme négative et interrogative :} Mettez les phrases suivantes à la forme négative, puis à la forme interrogative (question totale).
    \begin{enumerate}
        \item They have already paid the bride price.
        \item My uncle used to tell us folktales.
        \item She is looking forward to the wedding.
        \item There were many guests at the reception.
        \item He must respect the traditions.
    \end{enumerate}
\end{enumerate}
\end{exercice}

\begin{exercice}[Texte à trous : articles et prépositions]
\emph{Remplissez les blancs par \eng{a}, \eng{an}, \eng{the}, \eng{—} (zéro article) ou la préposition appropriée (\eng{in}, \eng{on}, \eng{at}, \eng{to}, \eng{for}, \eng{with}, \eng{of}, \eng{by}, \eng{since}, \eng{during}).}
\begin{textebox}[A Traditional Ceremony in the Grassfields]
\_\_\_\_\_\_\_\_\_\_\_\_ \eng{Fon}'s palace is \_\_\_\_\_\_\_\_\_\_\_\_ centre of social life. \_\_\_\_\_\_\_\_\_\_\_\_ last Saturday, \_\_\_\_\_\_\_\_\_\_\_\_ grand ceremony took place \_\_\_\_\_\_\_\_\_\_\_\_ honour \_\_\_\_\_\_\_\_\_\_\_\_ the new title holders. People came \_\_\_\_\_\_\_\_\_\_\_\_ far and near \_\_\_\_\_\_\_\_\_\_\_\_ witness the event. The ceremony started \_\_\_\_\_\_\_\_\_\_\_\_ 10 a.m. \_\_\_\_\_\_\_\_\_\_\_\_ the sound of drums. \_\_\_\_\_\_\_\_\_\_\_\_ notables, dressed \_\_\_\_\_\_\_\_\_\_\_\_ their traditional \eng{toghu}, sat \_\_\_\_\_\_\_\_\_\_\_\_ the main courtyard. \_\_\_\_\_\_\_\_\_\_\_\_ \eng{Fon} addressed \_\_\_\_\_\_\_\_\_\_\_\_ crowd \_\_\_\_\_\_\_\_\_\_\_\_ a loud voice. He spoke \_\_\_\_\_\_\_\_\_\_\_\_ unity and development. \_\_\_\_\_\_\_\_\_\_\_\_ ceremony lasted \_\_\_\_\_\_\_\_\_\_\_\_ three hours. Everyone enjoyed \_\_\_\_\_\_\_\_\_\_\_\_ meal served \_\_\_\_\_\_\_\_\_\_\_\_ banana leaves.
\end{textebox}
\end{exercice}

\begin{exercice}[Appariement : structures et fonctions]
\emph{Associez chaque énoncé (1--8) à la notion grammaticale qu'il illustre (A--H).}
\begin{center}
\begin{tabular}{|c|l|c|l|}
\hline
\textbf{Énoncé} & & \textbf{Notion} & \\
\hline
1. & \eng{If I were you, I would accept the offer.} & A. & \eng{Present Perfect (expérience)} \\
2. & \eng{She has been waiting for two hours.} & B. & \eng{Second Conditional} \\
3. & \eng{The meeting was postponed by the chairman.} & C. & \eng{Passive Voice} \\
4. & \eng{He denied stealing the money.} & D. & \eng{Present Perfect Continuous} \\
5. & \eng{Not only did he come late, but he also forgot the gift.} & E. & \eng{Gerund after verb} \\
6. & \eng{I wish I knew the answer.} & F. & \eng{Inversion after negative adverbial} \\
7. & \eng{Having finished his work, he went home.} & G. & \eng{Wish + Past Simple (regret présent)} \\
8. & \eng{This is the man whose daughter won the prize.} & H. & \eng{Relative clause (possessive)} \\
\hline
\end{tabular}
\end{center}
\end{exercice}

\courssub{Je me prépare au Bac}

\begin{exercice}[Type Bac : Correction d'erreurs et questions de langue]
\emph{Le paragraphe ci-dessous contient \textbf{10 erreurs grammaticales} (accord, temps, articles, prépositions, ordre des mots, orthographe). Recopiez le texte en corrigeant les erreurs (soulignez vos corrections). Ensuite, répondez aux questions.}
\begin{textebox}[Social Life in Modern Cameroon]
In the modern Cameroon, the family structure is changing rapid. The traditional extended family, who used to be the backbone of the society, is gradually replaced by the nuclear family. Many young peoples leave the villages to go in towns like Yaoundé or Douala for search of jobs. This situation has created news challenges. For example, the elderly peoples are often left alone in the village without nobody to care for them. However, the strong kinship ties still exist. During the holidays, the city-dwellers returns to their villages to participate to family meetings and ceremonies. The bride price, although criticised by some, remains an important custom who strengthens the bonds between two families. It is high time the government take measures to support the family unit.
\end{textebox}
\begin{enumerate}
    \item Identifiez dans le texte corrigé : deux exemples de \eng{Present Perfect}, un exemple de \eng{Passive Voice}, un exemple de \eng{Relative Clause}.
    \item Donnez la forme du \eng{Base Form} (infinitif sans \eng{to}) des verbes soulignés dans le texte original : \eng{is changing}, \eng{used to be}, \eng{is replaced}, \eng{has created}, \eng{are left}, \eng{returns}, \eng{strengthens}.
    \item Transformez la phrase \eng{“Many young peoples leave the villages...”} en phrase négative et en question totale.
\end{enumerate}
\end{exercice}

\begin{exercice}[Type Bac : Traduction Thème]
\emph{Traduisez le texte suivant en anglais. Votre traduction doit être naturelle et respecter les règles de grammaire étudiées (temps, articles, prépositions, accord).}
\begin{quote}
« Au Cameroun, le mariage n'est pas seulement l'union de deux individus, c'est une alliance entre deux familles. Avant la cérémonie officielle, la famille du fiancé doit verser la dot, qui consiste souvent en de l'argent, des tissus et des animaux. Les anciens des deux familles se réunissent pour négocier. Une fois la dot payée, la mariée rejoint la maison de son mari. Aujourd'hui, cependant, certains jeunes couples préfèrent le mariage civil à la mairie, suivi d'une petite fête. Mais même dans ce cas, les traditions ancestrales ne sont pas oubliées. »
\end{quote}
\end{exercice}

\begin{exercice}[Type Bac : Expression écrite guidée]
\emph{Traitez \textbf{un seul} des deux sujets suivants. Votre composition comportera \textbf{environ 150 mots}. Elle sera évaluée sur la correction de la langue (grammaire, vocabulaire, orthographe), la cohérence et l'organisation des idées.}
\begin{enumerate}
    \item \textbf{Subject A:} \eng{“Write an article for your school magazine describing an important family event you attended recently (a wedding, a birth celebration, a funeral, a traditional ceremony). Use the \textbf{Present Perfect} at least twice and the \textbf{Past Simple} to narrate the events. Mention the role of the extended family.”}
    \item \textbf{Subject B:} \eng{“Some people think that traditional family values are disappearing in modern Cameroonian society. Others believe they are adapting to new realities. Write an essay giving your opinion. Use linking words (\eng{however, furthermore, consequently, in conclusion}) and at least one \eng{If-clause} (Conditionnel).”}
\end{enumerate}
\end{exercice}

\newpage

\coursec{Corrigés des exercices}

\coursec{Corrigés des exercices — Lesson 2 : Grammar Revision}
\courssub{Chapter 1 : Family and Social Life}

\begin{corrige}[Exercice 1 — QCM : zones de difficulté grammaticale]
\begin{enumerate}
    \item \eng{My uncle has been in Douala \textbf{since} 2015.} — \cle{since} + point de départ (une date) ; \eng{for} + durée (\eng{for ten years}).
    \item \eng{Look! The children \textbf{are playing} football in the courtyard right now.} — \eng{Look!} et \eng{right now} signalent une action en cours : \cle{Present Continuous}.
    \item \eng{\textbf{—} Honesty is a virtue highly valued in our community.} — article \cle{zéro} devant un nom abstrait non comptable pris dans un sens général.
    \item \eng{Neither the president nor his assistants \textbf{were} present at the meeting yesterday.} — avec \eng{neither... nor}, le verbe s'accorde avec le sujet le plus proche (\eng{his assistants}, pluriel).
    \item \eng{She asked me where \textbf{I lived}.} — question indirecte : ordre sujet + verbe (pas d'inversion) et concordance des temps (\eng{asked} $\rightarrow$ \eng{lived}).
    \item \eng{If he \textbf{comes} early, we will go to the market together.} — après \eng{if} (condition réelle), jamais de \eng{will} : \cle{Present Simple} dans la subordonnée, \eng{will} dans la principale.
    \item \eng{My father \textbf{has gone} to Buea; he will be back tomorrow.} — \eng{has gone to} = il y est allé et y est encore ; \eng{has been to} = il y est allé et en est revenu.
    \item \eng{There isn't \textbf{any} water left in the jar.} — \eng{any} dans les phrases négatives ; \eng{some} à l'affirmative.
    \item \eng{The news \textbf{is} very encouraging for the students.} — \eng{news} est un faux pluriel : nom \cle{singulier} en anglais.
    \item \eng{I am used to \textbf{eating} fufu and eru on Sundays.} — \eng{to be used to} + \cle{V-ing} (ici \eng{to} est une préposition, pas l'infinitif).
\end{enumerate}
\textbf{Barème indicatif :} 1 point par réponse correcte, soit 10 points.
\end{corrige}

\begin{corrige}[Exercice 2 — Vrai ou Faux grammatical : justifiez]
\begin{enumerate}
    \item \textbf{Incorrect.} \eng{My brother and his friend \textbf{are} coming to Yaoundé next week.} — Règle violée : \cle{accord sujet-verbe}. Deux sujets reliés par \eng{and} = sujet pluriel.
    \item \textbf{Incorrect.} \eng{She has been living in Bamenda \textbf{for} three years.} — Règle violée : \cle{préposition}. \eng{for} + durée ; \eng{since} + point de départ (\eng{since 2021}).
    \item \textbf{Incorrect.} \eng{The police \textbf{are} investigating the matter.} — Règle violée : \cle{accord}. \eng{police} est un nom collectif toujours pluriel en anglais.
    \item \textbf{Incorrect.} \eng{He suggested \textbf{that I (should) go} to the hospital.} (ou : \eng{He suggested going to the hospital.}) — Règle violée : construction de \eng{suggest}. On ne dit jamais \eng{suggest somebody to do} ; \eng{suggest} + \eng{that}-proposition ou + V-ing.
    \item \textbf{Incorrect.} \eng{\textbf{Although} the rain was heavy, we attended the ceremony.} (ou : \eng{\textbf{Despite the heavy rain}, we attended the ceremony.}) — Règle violée : \cle{connecteur}. \eng{Despite} + nom ; \eng{although} + proposition (sujet + verbe). Le français « malgré que » n'existe pas en anglais.
    \item \textbf{Incorrect.} \eng{My mother cooks delicious ndolé.} — Règle violée : \cle{double sujet}. En anglais, un seul sujet : \eng{My mother} \emph{ou} \eng{She}, pas les deux.
    \item \textbf{Incorrect.} \eng{I prefer tea \textbf{to} coffee.} — Règle violée : \cle{préposition}. \eng{prefer A to B} ; \eng{than} est réservé aux comparatifs (\eng{better than}).
    \item \textbf{Incorrect.} \eng{It is high time we \textbf{left} this place.} — Règle violée : \cle{temps}. Après \eng{It is (high) time}, on emploie le \eng{Past Simple} à valeur de présent (subjonctif irréel).
\end{enumerate}
\textbf{Barème indicatif :} 1 point par identification correcte + 1 point par correction exacte, soit 16 points.
\end{corrige}

\begin{corrige}[Exercice 3 — Vocabulaire : Family and Social Life]
\begin{enumerate}
    \item \eng{In many Cameroonian cultures, the payment of the \textbf{bride price} seals the marriage alliance between two families.} — \eng{bride price} (nom composé) : la dot versée par la famille du fiancé.
    \item \eng{My \textbf{nuclear family} includes my parents, my two brothers and myself; my grandparents live in the village.} — \eng{nuclear family} : famille nucléaire (parents + enfants), par opposition à \eng{extended family}.
    \item \eng{The \textbf{lineage} system traces descent through the father's side in the Beti tradition.} — \eng{lineage} (nom) : lignage, filiation.
    \item \eng{After the death of her husband, the \textbf{widow} wore mourning clothes for a year.} — \eng{widow} (nom) : veuve (le masculin est \eng{widower}).
    \item \eng{The \textbf{inheritance} to the throne follows strict traditional rules in the Grassfields.} — \eng{inheritance} (nom) : héritage, succession.
    \item \eng{During the \textbf{ceremony}, the elders pour libations to honour the ancestors.} — \eng{ceremony} (nom) : cérémonie.
\end{enumerate}
\textbf{Mots en trop :} \eng{extended family}, \eng{kinship}, \eng{foster}, \eng{sibling}. (Quatre mots excédentaires étaient en réalité présents dans la liste.)
\textbf{Barème indicatif :} 1 point par mot correctement placé, soit 6 points.
\end{corrige}

\begin{corrige}[Exercice 4 — Conjugaison et accord : mise en forme]
\begin{enumerate}
    \item \eng{Every morning, the bendskin driver \textbf{stops} in front of the school gate.} — \cle{Present Simple} (habitude, marqueur \eng{every morning}) ; 3\textsuperscript{e} personne du singulier : \eng{stop} + \eng{-s}.
    \item \eng{Listen! Somebody \textbf{is playing} the balafon near the palace.} — \cle{Present Continuous} (action en cours, marqueur \eng{Listen!}) ; \eng{somebody} = singulier.
    \item \eng{My elder sister \textbf{has just finished} her degree at the University of Buea.} — \cle{Present Perfect} avec \eng{just} (action récente au résultat présent).
    \item \eng{While we \textbf{were having} dinner last night, the power \textbf{went} out.} — \cle{Past Continuous} (action longue en arrière-plan) interrompu par le \cle{Past Simple} (action brève) ; \eng{go} $\rightarrow$ \eng{went} (irrégulier).
    \item \eng{How long \textbf{have you known} Mr. Nkeng? — Since 2010.} — \cle{Present Perfect} avec \eng{How long...?} et \eng{since} ; \eng{know} $\rightarrow$ \eng{known}.
    \item \eng{The number of students in this class \textbf{increases} every year.} — \cle{Present Simple} ; attention : le sujet est \eng{the number} (singulier), pas \eng{students}.
    \item \eng{By the time the chief arrived, the dancers \textbf{had already been performing} for an hour.} — \cle{Past Perfect Continuous} : action antérieure à un repère passé (\eng{arrived}) avec idée de durée. (\eng{had already performed} est aussi acceptable si l'on insiste sur l'antériorité simple.)
\end{enumerate}
\textbf{Barème indicatif :} 1 point par verbe correctement conjugué (temps + accord), soit 8 points.
\end{corrige}

\begin{corrige}[Exercice 5 — Traduction ciblée : interférences français-anglais]
\begin{enumerate}
    \item \eng{My father has been a teacher for twenty years.} — Durée qui continue : \cle{Present Perfect} + \eng{for} ; jamais le présent simple comme en français.
    \item \eng{He left for Douala yesterday morning.} — Repère passé daté : \cle{Past Simple} ; \eng{leave for} + destination.
    \item \eng{She told me (that) she would come the next day.} — Concordance des temps : \eng{will} $\rightarrow$ \eng{would} ; \eng{tomorrow} $\rightarrow$ \eng{the next day}.
    \item \eng{I have nobody to talk to.} (ou : \eng{I don't have anybody to talk to.}) — Une seule négation par phrase ; la préposition finale \eng{to} est obligatoire.
    \item \eng{This is the first time I have eaten ndolé.} — Après \eng{This is the first time}, \cle{Present Perfect} obligatoire.
    \item \eng{He suggested waiting for the bus.} (ou : \eng{He suggested that we (should) wait for the bus.}) — \eng{suggest} + V-ing ; \eng{wait for} + complément.
    \item \eng{Although he was tired, he worked.} (ou : \eng{Despite being tired, he worked.} / \eng{Despite his tiredness, he worked.}) — « Malgré que » ne se traduit jamais par \eng{despite} + proposition.
    \item \eng{She got married to a doctor.} (ou : \eng{She married a doctor.}) — \eng{marry} est transitif direct : pas de \eng{with} ; avec \eng{get married}, la préposition est \eng{to}.
    \item \eng{I asked him what he was doing.} — Question indirecte : ordre sujet + verbe, sans inversion ni auxiliaire \eng{did}.
    \item \eng{I haven't seen you for a long time.} (ou : \eng{It has been a long time since I last saw you.}) — \cle{Present Perfect} pour une situation qui dure jusqu'au présent.
\end{enumerate}
\textbf{Points de vigilance :} le piège majeur est le calque du présent français avec \eng{since/for} ; vérifier aussi la transitivité des verbes (\eng{marry}, \eng{wait for}, \eng{talk to}).
\textbf{Barème indicatif :} 2 points par phrase (1 point pour le temps/la structure, 1 point pour la préposition ou l'ordre des mots), soit 20 points.
\end{corrige}

\begin{corrige}[Exercice 6 — Transformations structurales]
\textbf{1. Discours indirect :}
\begin{enumerate}
    \item \eng{She said (that) she would visit her grandparents in the village the following week.} — \eng{will} $\rightarrow$ \eng{would} ; \eng{next week} $\rightarrow$ \eng{the following week}.
    \item \eng{The teacher asked me if/whether I had finished my homework.} — Question totale : \eng{if/whether} + ordre sujet-verbe ; \eng{Present Perfect} $\rightarrow$ \eng{Past Perfect}.
    \item \eng{The elder told the children not to make noise during the ceremony.} — Ordre rapporté : \eng{tell + objet + not + to}-infinitif.
    \item \eng{He wanted to know where my brother worked.} — Question partielle rapportée : mot interrogatif + sujet + verbe ; \eng{does... work} $\rightarrow$ \eng{worked}.
    \item \eng{They proposed organizing a fundraising for the school.} (ou : \eng{They proposed that they (should) organize a fundraising for the school.}) — \eng{Let's} $\rightarrow$ \eng{suggest/propose} + V-ing.
\end{enumerate}

\textbf{2. Forme négative et interrogative :}
\begin{enumerate}
    \item Négative : \eng{They have not paid the bride price yet.} Interrogative : \eng{Have they paid the bride price yet?} — \eng{already} $\rightarrow$ \eng{yet} en fin de phrase.
    \item Négative : \eng{My uncle did not use to tell us folktales.} (ou : \eng{My uncle never used to tell us folktales.}) Interrogative : \eng{Did your uncle use to tell you folktales?} — Avec l'auxiliaire \eng{did}, on écrit \eng{use to} (sans \eng{-d}).
    \item Négative : \eng{She is not looking forward to the wedding.} Interrogative : \eng{Is she looking forward to the wedding?}
    \item Négative : \eng{There were not many guests at the reception.} Interrogative : \eng{Were there many guests at the reception?}
    \item Négative (interdiction) : \eng{He must not respect the traditions.} Négative (absence d'obligation) : \eng{He doesn't have to respect the traditions.} Interrogative : \eng{Must he respect the traditions?}
\end{enumerate}
\textbf{Barème indicatif :} 1 point par transformation correcte, soit 15 points.
\end{corrige}

\begin{corrige}[Exercice 7 — Texte à trous : articles et prépositions]
\begin{textebox}[A Traditional Ceremony in the Grassfields — version corrigée]
\textbf{The} \eng{Fon}'s palace is \textbf{the} centre of social life. \textbf{—} last Saturday, \textbf{a} grand ceremony took place \textbf{in} honour \textbf{of} the new title holders. People came \textbf{from} far and near \textbf{to} witness the event. The ceremony started \textbf{at} 10 a.m. \textbf{to/with} the sound of drums. \textbf{The} notables, dressed \textbf{in} their traditional \eng{toghu}, sat \textbf{in} the main courtyard. \textbf{The} \eng{Fon} addressed \textbf{the} crowd \textbf{in} a loud voice. He spoke \textbf{of/about} unity and development. \textbf{The} ceremony lasted \textbf{for} three hours. Everyone enjoyed \textbf{the} meal served \textbf{on} banana leaves.
\end{textebox}
\textbf{Justifications essentielles :}
\begin{itemize}
    \item \eng{The Fon's palace} : article défini, le \eng{Fon} est unique et identifié ; \eng{the centre of} : complément du nom déterminatif.
    \item \eng{— last Saturday} : article zéro devant les expressions de temps avec \eng{last/next}.
    \item \eng{a grand ceremony} : première mention d'un élément non identifié $\rightarrow$ article indéfini.
    \item \eng{in honour of} : locution figée ; \eng{from far and near} : provenance ; \eng{to witness} : infinitif de but.
    \item \eng{at 10 a.m.} : \eng{at} + heure précise ; \eng{dressed in} : vêtu de ; \eng{in the courtyard} : lieu délimité.
    \item \eng{in a loud voice} : manière ; \eng{speak of/about} : parler de ; \eng{lasted for three hours} : \eng{for} + durée ; \eng{on banana leaves} : support (servi sur).
\end{itemize}
\textbf{Barème indicatif :} 0,5 point par blanc correct, soit 10 points.
\end{corrige}

\begin{corrige}[Exercice 8 — Appariement : structures et fonctions]
\begin{center}
\begin{tabular}{|c|l|}
\hline
\rowcolor{popblueL} \textbf{Énoncé} & \textbf{Réponse} \\
\hline
1. \eng{If I were you, I would accept the offer.} & \textbf{B} — \eng{Second Conditional} (\eng{If + Past Simple, would + V}) \\
\hline
2. \eng{She has been waiting for two hours.} & \textbf{D} — \eng{Present Perfect Continuous} (durée jusqu'au présent) \\
\hline
3. \eng{The meeting was postponed by the chairman.} & \textbf{C} — \eng{Passive Voice} (\eng{be} + participe passé + \eng{by}) \\
\hline
4. \eng{He denied stealing the money.} & \textbf{E} — \eng{Gerund after verb} (\eng{deny} + V-ing) \\
\hline
5. \eng{Not only did he come late, but he also forgot the gift.} & \textbf{F} — \eng{Inversion after negative adverbial} (\eng{Not only did he...}) \\
\hline
6. \eng{I wish I knew the answer.} & \textbf{G} — \eng{Wish + Past Simple} (regret portant sur le présent) \\
\hline
7. \eng{Having finished his work, he went home.} & — participe parfait (proposition participiale d'antériorité) ; \emph{notion non listée : voir remarque ci-dessous.} \\
\hline
8. \eng{This is the man whose daughter won the prize.} & \textbf{H} — \eng{Relative clause (possessive)} avec \eng{whose} \\
\hline
\end{tabular}
\end{center}
\textbf{Remarque méthodologique :} l'énoncé 7 illustre la \cle{proposition participiale} (\eng{perfect participle}) ; la notion \textbf{A} (\eng{Present Perfect}) correspondrait à une phrase du type \eng{I have visited Buea twice.} En situation d'examen, si une notion reste sans énoncé, signalez-le clairement dans votre copie : cela fait partie de la compétence d'analyse grammaticale.
\textbf{Barème indicatif :} 1 point par appariement correct, soit 8 points.
\end{corrige}

\begin{corrige}[Exercice 9 — Type Bac : Correction d'erreurs et questions de langue]
\textbf{Texte corrigé (les corrections sont en gras) :}
\begin{textebox}[Social Life in Modern Cameroon — corrected version]
In \textbf{—} modern Cameroon, the family structure is changing \textbf{rapidly}. The traditional extended family, \textbf{which} used to be the backbone of \textbf{—} society, is gradually \textbf{being} replaced by the nuclear family. Many young \textbf{people} leave the villages to go \textbf{to} towns like Yaoundé or Douala \textbf{in} search of jobs. This situation has created \textbf{new} challenges. For example, \textbf{—} elderly people are often left alone in the village without \textbf{anybody} to care for them. However, \textbf{—} strong kinship ties still exist. During the holidays, the city-dwellers \textbf{return} to their villages to participate \textbf{in} family meetings and ceremonies. The bride price, although criticised by some, remains an important custom \textbf{which/that} strengthens the bonds between two families. It is high time the government \textbf{took} measures to support the family unit.
\end{textebox}
\textbf{Détail des erreurs principales (14 fautes repérables, 10 attendues) :}
\begin{enumerate}
    \item \eng{In the modern Cameroon} $\rightarrow$ \eng{In modern Cameroon} : article zéro devant un nom propre de pays (même précédé d'un adjectif).
    \item \eng{changing rapid} $\rightarrow$ \eng{changing rapidly} : il faut un adverbe (en \eng{-ly}) pour modifier le verbe.
    \item \eng{family, who used to be} $\rightarrow$ \eng{which used to be} : \eng{who} pour les personnes, \eng{which} pour les choses/notions.
    \item \eng{the backbone of the society} $\rightarrow$ \eng{of society} : article zéro devant \eng{society} pris dans un sens général.
    \item \eng{is gradually replaced} $\rightarrow$ \eng{is gradually being replaced} : passif au présent continu (processus en cours).
    \item \eng{young peoples} $\rightarrow$ \eng{young people} : \eng{people} est déjà pluriel (le pluriel \eng{peoples} = « peuples »).
    \item \eng{go in towns} $\rightarrow$ \eng{go to towns} : mouvement vers $\rightarrow$ \eng{to}.
    \item \eng{for search of} $\rightarrow$ \eng{in search of} : locution figée.
    \item \eng{news challenges} $\rightarrow$ \eng{new challenges} : confusion \eng{news} (nouvelles, informations) / \eng{new} (nouveau).
    \item \eng{elderly peoples} $\rightarrow$ \eng{elderly people} ; \eng{without nobody} $\rightarrow$ \eng{without anybody} : double négation interdite.
    \item \eng{the strong kinship ties} $\rightarrow$ \eng{strong kinship ties} : article zéro devant un pluriel général ; \eng{city-dwellers returns} $\rightarrow$ \eng{return} : sujet pluriel.
    \item \eng{participate to} $\rightarrow$ \eng{participate in} : préposition correcte.
    \item \eng{a custom who strengthens} $\rightarrow$ \eng{which/that strengthens} : relatif pour une chose.
    \item \eng{It is high time the government take} $\rightarrow$ \eng{took} : \eng{Past Simple} après \eng{It is high time}.
\end{enumerate}
(\emph{Toute dizaine de corrections exactes et justifiées vaut le maximum.})
\textbf{Questions de langue :}
\begin{enumerate}
    \item Exemples attendus dans le texte corrigé : \\
    \eng{Present Perfect} : \eng{This situation \textbf{has created} new challenges}. \\
    \eng{Passive Voice} : \eng{is gradually being replaced by the nuclear family} ; \eng{the elderly people are often left alone} ; \eng{although criticised by some}. \\
    \eng{Relative Clause} : \eng{which used to be the backbone of society} ; \eng{which/that strengthens the bonds between two families}.
    \item Formes de base (infinitif sans \eng{to}) : \eng{is changing} $\rightarrow$ \eng{change} ; \eng{used to be} $\rightarrow$ \eng{be} ; \eng{is replaced} $\rightarrow$ \eng{replace} ; \eng{has created} $\rightarrow$ \eng{create} ; \eng{are left} $\rightarrow$ \eng{leave} ; \eng{returns} $\rightarrow$ \eng{return} ; \eng{strengthens} $\rightarrow$ \eng{strengthen}.
    \item Phrase corrigée : \eng{Many young people leave the villages.} \\
    Négative : \eng{Many young people do not leave the villages.} \\
    Interrogative : \eng{Do many young people leave the villages?}
\end{enumerate}
\textbf{Points de vigilance :} ne pas « corriger » ce qui est juste ; souligner uniquement les modifications ; pour les questions de langue, recopier la phrase entière et mettre la structure en évidence.
\textbf{Barème indicatif :} 10 points pour la correction (1 point par erreur), 4 points pour la question 1, 3,5 points pour la question 2 (0,5 par base verbale), 2,5 points pour la question 3, soit 20 points.
\end{corrige}

\begin{corrige}[Exercice 10 — Type Bac : Traduction Thème]
\textbf{Proposition de traduction :}
\begin{textebox}[Model translation]
In Cameroon, marriage is not only the union of two individuals; it is an alliance between two families. Before the official ceremony, the fiancé's family must pay the bride price, which often consists of money, cloth and animals. The elders of both families meet to negotiate. Once the bride price has been paid, the bride moves into her husband's house. Nowadays, however, some young couples prefer a civil marriage at the town hall, followed by a small party. But even in that case, the ancestral traditions are not forgotten.
\end{textebox}
\textbf{Points de vigilance :}
\begin{itemize}
    \item « la dot » se traduit \eng{the bride price} (ou \eng{bride wealth}) ; éviter le calque \eng{the dowry} si l'on veut être précis, mais \eng{dowry} reste accepté.
    \item « consister en » = \eng{consist of} (jamais \eng{consist in} pour énumérer des éléments).
    \item « les anciens » = \eng{the elders} ; « les deux familles » = \eng{both families} ou \eng{the two families}.
    \item « une fois la dot payée » = \eng{once the bride price has been paid} (\eng{Present Perfect} passif) ou \eng{once the bride price is paid}.
    \item « à la mairie » = \eng{at the town hall} ; « suivi de » = \eng{followed by}.
    \item Accord et articles : \eng{marriage is}, \eng{the traditions are not forgotten} (passif au présent).
\end{itemize}
\textbf{Barème indicatif :} 20 points — exactitude grammaticale (8), fidélité du sens (6), lexique approprié (4), fluidité et orthographe (2). Toute traduction correcte et naturelle différente du modèle est acceptée.
\end{corrige}

\begin{corrige}[Exercice 11 — Type Bac : Expression écrite guidée]
\textbf{Subject A — proposition de rédaction modèle (environ 150 mots) :}
\begin{textebox}[A Wedding to Remember — model article]
Last month, I attended my cousin's traditional wedding in Bafoussam, and it was one of the most memorable family events I have ever experienced. The ceremony took place in my uncle's large courtyard, which the extended family had decorated with colourful fabrics and flowers. Early in the morning, the elders gathered to negotiate the final details of the bride price, while the women prepared enormous pots of ndolé and achu. When the bride finally appeared, dressed in a magnificent toghu outfit, everybody ululated with joy. The Fon poured libations to honour our ancestors, and the drummers played until midnight. I have never seen my grandmother so happy: she danced with all her grandchildren and told us stories about her own wedding fifty years ago. This celebration reminded me that the extended family remains the heart of our social life. I will never forget that day.
\end{textebox}

\textbf{Subject B — proposition de rédaction modèle (environ 150 mots) :}
\begin{textebox}[Traditional Values: Disappearing or Adapting? — model essay]
Many people claim that traditional family values are disappearing in modern Cameroon. In my opinion, they are not disappearing; they are simply adapting to new realities. It is true that the extended family is less present in big cities, where young couples often live alone in small flats. However, kinship ties remain extremely strong. During holidays, thousands of city-dwellers travel back to their villages to attend family meetings and ceremonies. Furthermore, modern technology even strengthens these bonds: family members who live abroad send money home and take part in discussions through telephone calls. Consequently, values such as solidarity and respect for elders survive, even if their forms change. If we abandoned these values completely, our society would lose its identity. In conclusion, tradition and modernity are not enemies: the Cameroonian family is transforming itself, but its deep roots remain alive.
\end{textebox}

\textbf{Points de vigilance (critères d'évaluation) :}
\begin{itemize}
    \item \textbf{Respect de la tâche :} sujet A = récit d'un événement réel ou plausible, avec au moins deux \eng{Present Perfect} (\eng{I have ever experienced}, \eng{I have never seen}) et une narration au \eng{Past Simple} ; sujet B = prise de position argumentée avec connecteurs et au moins une \eng{If-clause} (\eng{If we abandoned..., our society would lose...}).
    \item \textbf{Langue :} accords sujet-verbe, temps cohérents, articles et prépositions corrects, orthographe soignée (britannique).
    \item \textbf{Organisation :} introduction, développement en paragraphes, conclusion ; connecteurs variés (\eng{however, furthermore, consequently, in conclusion}).
    \item \textbf{Longueur :} environ 150 mots (une marge de 10 \% est tolérée).
\end{itemize}
\textbf{Barème indicatif :} 20 points — respect de la tâche et contenu (6), correction grammaticale (6), vocabulaire (4), cohérence et organisation (4).
\end{corrige}

\newpage

\coursec{Fiche bilan}

\begin{popfigure}
\begin{tikzpicture}[mindmap, grow cyclic, every node/.style={concept, align=center, font=\small},
root concept/.append style={concept color=popblue, font=\bfseries},
level 1/.append style={level distance=3.4cm, sibling angle=51.4},
level 2/.append style={level distance=2.2cm, font=\scriptsize}]
\node[root concept, align=center]{Grammar Problem Areas\\ Tle A}
child[concept color=poporange]{ node{Subject-Verb Agreement}
  child[concept color=poporange!60]{ node{3\textsuperscript{e} pers. -s} }
  child[concept color=poporange!60]{ node{each, everyone + sing.} }
}
child[concept color=popgreen]{ node{Present Perfect vs Past Simple}
  child[concept color=popgreen!60]{ node{for/since/just/yet} }
  child[concept color=popgreen!60]{ node{ago/yesterday $\to$ past} }
}
child[concept color=popblue!70]{ node{Articles}
  child[concept color=popblue!40]{ node{a / an / the / $\varnothing$} }
}
child[concept color=popgold]{ node{Prepositions}
  child[concept color=popgold!60]{ node{in / on / at} }
  child[concept color=popgold!60]{ node{depend on, listen to} }
}
child[concept color=poppink]{ node{False Friends}
  child[concept color=poppink!60]{ node{actually $\neq$ actuellement} }
}
child[concept color=poppurple]{ node{Word Order}
  child[concept color=poppurple!60]{ node{S + V + C} }
  child[concept color=poppurple!60]{ node{adverbe avant verbe} }
}
child[concept color=popgreen!80!popblue]{ node{Transformations}
  child[concept color=popgreen!50]{ node{affirm. $\to$ nég. / question} }
};
\end{tikzpicture}
\legende{Carte mentale : les zones à problèmes de la grammaire anglaise en Tle A}
\end{popfigure}

\begin{bilan}
\textbf{1. Lexique clé de la vie familiale et sociale}

\begin{center}
\begin{tabularx}{\linewidth}{|X|X|}
\hline
\rowcolor{popblueL} English & Français \\
\hline
relative / extended family & parent / famille élargie \\
\hline
upbringing & éducation (reçue dans l'enfance) \\
\hline
household chores & tâches ménagères \\
\hline
generation gap & fossé des générations \\
\hline
to get on with somebody & bien s'entendre avec quelqu'un \\
\hline
to take after somebody & ressembler à (un parent) \\
\hline
\end{tabularx}
\end{center}

\textbf{2. Règles à connaître par cœur}

\begin{center}
\begin{tabularx}{\linewidth}{|X|X|}
\hline
\rowcolor{popblueL} Règle & Exemple \\
\hline
Accord sujet-verbe : 3\textsuperscript{e} pers. du singulier en \textbf{-s} au present simple ; \eng{each, every, everyone, nobody} + verbe au singulier. & \eng{Everybody in my family speaks English.} \\
\hline
\cle{Present perfect} : bilan présent, pas de repère passé précis (\eng{just, already, yet, ever, never, for, since}). & \eng{I have lived in Yaoundé since 2015.} \\
\hline
\cle{Past simple} : action datée, terminée (\eng{ago, yesterday, last year, when...}). & \eng{My uncle visited us two weeks ago.} \\
\hline
Articles : \eng{a/an} (indéfini singulier), \eng{the} (défini), $\varnothing$ (généralités plurielles et noms abstraits). & \eng{Family life is changing.} \\
\hline
Prépositions figées : \eng{depend on, listen to, be interested in, be good at, get married to}. & \eng{Children depend on their parents.} \\
\hline
Ordre des mots : Sujet + Verbe + Complément ; l'adverbe de fréquence se place avant le verbe principal. & \eng{She often visits her grandmother.} \\
\hline
\end{tabularx}
\end{center}

\textbf{3. Méthodes express}
\begin{itemize}
\item Repérer d'abord le \cle{sujet} (et sa nature : singulier ou pluriel) avant de conjuguer.
\item Chercher un \cle{marqueur temporel} : \eng{ago, last, yesterday} $\to$ past simple ; \eng{since, for, just, yet} $\to$ present perfect.
\item Pour les transformations : identifier le temps, garder le même temps dans la forme négative ou interrogative, ajouter l'auxiliaire \eng{do/does/did} si nécessaire.
\item Traduire mentalement puis vérifier les \cle{faux amis} avant d'écrire.
\end{itemize}
\end{bilan}

\begin{aretenir}[Checklist avant le Bac]
Avant de rendre votre copie, vérifiez chaque point :
\begin{itemize}
\item[$\square$] Chaque verbe au present simple à la 3\textsuperscript{e} personne du singulier porte bien un \textbf{-s} (\eng{she works}, pas \eng{she work}).
\item[$\square$] Je n'ai pas utilisé le present perfect avec un repère de temps passé (\eng{ago, yesterday, last week}).
\item[$\square$] Je n'ai pas mis \eng{the} devant une généralité (\eng{Life is hard}, pas \eng{the life is hard}).
\item[$\square$] Mes prépositions sont correctes : \eng{depend on}, \eng{married to}, \eng{interested in}, \eng{good at}.
\item[$\square$] Mes adverbes de fréquence sont placés avant le verbe (\eng{I always help}, pas \eng{I help always}).
\item[$\square$] J'ai évité les faux amis : \eng{actually} = en fait, \eng{sensible} = raisonnable, \eng{a library} = une bibliothèque.
\item[$\square$] Mes questions et négations utilisent bien l'auxiliaire \eng{do/does/did} ou \eng{have/has}.
\item[$\square$] L'ordre Sujet + Verbe + Complément est respecté dans chaque phrase.
\item[$\square$] \eng{Everyone, each, nobody, somebody} sont suivis d'un verbe au singulier.
\item[$\square$] J'ai relu ma production à voix basse pour détecter les accords fautifs.
\end{itemize}
\end{aretenir}

\findecours

\end{document}
